% Lutte contre le VIH/Sida — SVT 1ère TI — Cours signé OnBuch+
% Programme officiel MINESEC. Compiler avec : tectonic NCT04-lutte-vih-sida-c-ti.tex
\newcommand{\DOCMATIERE}{SVT}
\newcommand{\DOCNIVEAU}{1ère TI}
\newcommand{\DOCMODULE}{Module 2 : L'Éducation à la Santé}
\newcommand{\DOCLECON}{NCT04}
\newcommand{\DOCTITRE}{Lutte contre le VIH/Sida}
% OnBuch+ — préambule « cours signés » ENRICHI (compilé avec tectonic / XeLaTeX)
% Système visuel « Pop » d'OnBuch+ : fond crème (jamais blanc), bordures encre
% épaisses, ombres « dures » décalées, titres Archivo Black en capitales.
% Version MATHÉMATIQUES : graphes (pgfplots/tikz), boîtes pédagogiques
% (définition, propriété, méthode, attention, le savais-tu, exercice résolu,
% exercice, TP, bilan), page de garde et signature OnBuch+ ancrée en bas.
%
% Macros attendues dans main.tex AVANT \input{preamble} :
%   \DOCMATIERE  \DOCNIVEAU  \DOCMODULE  \DOCLECON (numéro)  \DOCTITRE
\documentclass[11pt]{article}

\usepackage{fontspec}
\usepackage[french]{babel}
\usepackage[a4paper,margin=2cm,top=3.4cm,bottom=2.8cm,headheight=1.4cm,footskip=1.3cm]{geometry}
\usepackage{amsmath,amssymb,amsfonts}
\usepackage{mathtools} % \xmapsto, \coloneqq... souvent utilisés par les modèles
\usepackage{cancel} % simplifications barrées (bilans de forces)
% \abs{z} (module, valeur absolue) et \norm{u} (norme), utilisés par les modèles
\providecommand{\abs}{}\let\abs\relax\DeclarePairedDelimiter{\abs}{\lvert}{\rvert}
\providecommand{\norm}{}\let\norm\relax\DeclarePairedDelimiter{\norm}{\lVert}{\rVert}
\usepackage[table]{xcolor} % [table] : fournit aussi \rowcolors (alternance de lignes)
\usepackage{pagecolor}
\usepackage{tikz}
\usetikzlibrary{arrows.meta,positioning,calc,shapes.geometric,shapes.misc,shapes.arrows,decorations.pathmorphing,decorations.pathreplacing,decorations.markings,patterns,shadows,mindmap,trees,fit,backgrounds,matrix,intersections,angles,quotes}
\usepackage{pgfplots}
\usepackage[version=4]{mhchem} % équations nucléaires \ce{^{235}_{92}U}
\usepackage[european,siunitx]{circuitikz} % schémas électriques
\pgfplotsset{compat=1.18}
\usepgfplotslibrary{fillbetween,groupplots} % aires entre courbes (calcul intégral)
\usepackage{enumitem}
\usepackage{fancyhdr}
% [most] charge toutes les bibliothèques tcolorbox (listings, minted, raster,
% documentation...) et gonfle inutilement la mémoire TeX sur un document long
% (~300 boîtes) : on ne charge que ce qui est réellement utilisé.
\usepackage[skins,breakable]{tcolorbox}
\usepackage{graphicx}
\usepackage{colortbl}
\usepackage{booktabs}
\usepackage{tabularx}
\usepackage{diagbox} % case d'en-tête diagonale des tableaux à double entrée
% Colonnes extensibles centrée / gauche / droite (C, L, R), souvent utilisées par les modèles
\newcolumntype{C}{>{\centering\arraybackslash}X}
\newcolumntype{L}{>{\raggedright\arraybackslash}X}
\newcolumntype{R}{>{\raggedleft\arraybackslash}X}
\usepackage{array}
\usepackage{multicol}
\usepackage{siunitx}
% parse-numbers=false : accepte aussi \SI{2,5 \times 10^{-3}}{...}
\sisetup{locale=FR,output-decimal-marker={,},group-minimum-digits=4,parse-numbers=false}
\DeclareSIUnit\ohm{\ensuremath{\symup{\Omega}}} % U+2126 absent de la police
\DeclareSIUnit\division{div} % réglages d'oscilloscope (ms/div, V/div)
\DeclareSIUnit\tr{tr} % tours (vitesse de rotation, ex. \SI{1500}{\tr\per\minute})
\DeclareSIUnit\ch{ch} % cheval-vapeur (puissance moteur, unité usuelle hors SI)
\DeclareSIUnit\franccfa{FCFA} % franc CFA (coûts, contexte camerounais)
\DeclareSIUnit\dioptre{\ensuremath{\delta}} % vergence d'une lentille (dioptrie)
\usepackage{qrcode}
% \degree : commande fréquemment utilisée par les modèles pour un angle ou une
% température en dehors de \SI{}{} (non fournie par les paquets chargés).
\providecommand{\degree}{^{\circ}}
% \qed (fin de démonstration), utilisé par les modèles sans amsthm
\providecommand{\qed}{\hfill$\square$}
% \barre : double barre des valeurs interdites dans un tableau de variations (tkz-tab)
\providecommand{\barre}{\ensuremath{\|}}
% environnement proof (amsthm non chargé) utilisé par les modèles
\ifdefined\proof\else\newenvironment{proof}[1][Démonstration]{\par\noindent\textit{#1.}\ }{\qed\par}\fi
\usepackage{lastpage}
\usepackage{hyperref}
\hypersetup{hidelinks}

% ============================================================
% Palette « Pop » OnBuch+ — mêmes hex que la classe P (plus_ui.dart)
% ============================================================
\definecolor{poporange}{HTML}{F59321}
\definecolor{popdark}{HTML}{E07A0C}
\definecolor{popcream}{HTML}{FFF6E8}
\definecolor{popink}{HTML}{1C1714}
\definecolor{poppurple}{HTML}{7A5AE0}
\definecolor{popgreen}{HTML}{2E9E6A}
\definecolor{poppink}{HTML}{E0457B}
\definecolor{popblue}{HTML}{2F7DE1}
\definecolor{popgold}{HTML}{C98A12}
\definecolor{popmuted}{HTML}{8A7F76}
\definecolor{popline}{HTML}{EADFD2}
\definecolor{popgray}{HTML}{8A7F76}
\colorlet{popgrey}{popgray}
% Alias tolérants (le modèle nomme parfois les couleurs différemment)
\colorlet{popwhite}{white}
\colorlet{popblack}{popink}
\colorlet{popred}{poppink}
\colorlet{popyellow}{popgold}
% Teintes claires pour fonds de tableaux / zones
\colorlet{popblueL}{popblue!10!white}
\colorlet{poporangeL}{poporange!14!white}
\colorlet{popgreenL}{popgreen!12!white}
\colorlet{poppurpleL}{poppurple!12!white}
\colorlet{poppinkL}{poppink!12!white}
\colorlet{popgoldL}{popgold!14!white}

\pagecolor{popcream}

% ============================================================
% Typographie
% ============================================================
\defaultfontfeatures{Ligatures=TeX}
\setmainfont{PlusJakartaSans-Medium.ttf}[Path=../../../fonts/,BoldFont=PlusJakartaSans-Bold.ttf]
% Maths en sans-serif (Fira Math) avec chiffres et lettres droites de Plus
% Jakarta, pour que les formules se fondent dans le texte.
\usepackage[math-style=ISO,bold-style=ISO]{unicode-math}
\setmathfont{FiraMath-Regular.otf}[Path=../../../fonts/]
\setmathfont{PlusJakartaSans-Medium.ttf}[Path=../../../fonts/,range={up/num,up/{Latin,latin},\mathup}]
\usepackage{icomma} % virgule décimale sans espace en mode math
% Caractères mathématiques Unicode absents des polices de texte (Plus Jakarta,
% Archivo Black) : rendus par leur équivalent mathématique, en texte comme en
% formule — sinon ils s'affichent en carré vide (ex. « Arithmétique dans ℤ »).
\usepackage{newunicodechar}
\newunicodechar{ℝ}{\ensuremath{\mathbb{R}}}
\newunicodechar{ℤ}{\ensuremath{\mathbb{Z}}}
\newunicodechar{ℂ}{\ensuremath{\mathbb{C}}}
\newunicodechar{ℕ}{\ensuremath{\mathbb{N}}}
\newunicodechar{ℚ}{\ensuremath{\mathbb{Q}}}
\newunicodechar{𝔻}{\ensuremath{\mathbb{D}}}
\newunicodechar{α}{\ensuremath{\alpha}}
\newunicodechar{β}{\ensuremath{\beta}}
\newunicodechar{γ}{\ensuremath{\gamma}}
\newunicodechar{δ}{\ensuremath{\delta}}
\newunicodechar{ε}{\ensuremath{\varepsilon}}
\newunicodechar{θ}{\ensuremath{\theta}}
\newunicodechar{λ}{\ensuremath{\lambda}}
\newunicodechar{μ}{\ensuremath{\mu}}
\newunicodechar{σ}{\ensuremath{\sigma}}
\newunicodechar{τ}{\ensuremath{\tau}}
\newunicodechar{φ}{\ensuremath{\varphi}}
\newunicodechar{ψ}{\ensuremath{\psi}}
\newunicodechar{ω}{\ensuremath{\omega}}
\newunicodechar{Σ}{\ensuremath{\Sigma}}
% majuscules produites par \MakeUppercase dans les titres (α-aminés -> Α)
\newunicodechar{Α}{\ensuremath{\alpha}}
\newunicodechar{Β}{\ensuremath{\beta}}
\newunicodechar{Γ}{\ensuremath{\Gamma}}
\newunicodechar{Δ}{\ensuremath{\Delta}}
\newunicodechar{Ε}{\ensuremath{\varepsilon}}
\newunicodechar{Θ}{\ensuremath{\Theta}}
\newunicodechar{Λ}{\ensuremath{\Lambda}}
\newunicodechar{Μ}{\ensuremath{\mu}}
\newunicodechar{Τ}{\ensuremath{\tau}}
\newunicodechar{Φ}{\ensuremath{\Phi}}
\newunicodechar{Ψ}{\ensuremath{\Psi}}
\newunicodechar{Ω}{\ensuremath{\Omega}}
\newunicodechar{↦}{\ensuremath{\mapsto}}
\newunicodechar{∈}{\ensuremath{\in}}
\newunicodechar{∉}{\ensuremath{\notin}}
\newunicodechar{∧}{\ensuremath{\wedge}}
\newunicodechar{⇔}{\ensuremath{\Leftrightarrow}}
\newunicodechar{⟺}{\ensuremath{\Longleftrightarrow}}
\newunicodechar{⇒}{\ensuremath{\Rightarrow}}
\newunicodechar{⟹}{\ensuremath{\Longrightarrow}}
\newunicodechar{∘}{\ensuremath{\circ}}
\newunicodechar{∪}{\ensuremath{\cup}}
\newunicodechar{∩}{\ensuremath{\cap}}
\newunicodechar{≡}{\ensuremath{\equiv}}
\newunicodechar{⊂}{\ensuremath{\subset}}
\newunicodechar{⊥}{\ensuremath{\perp}}
\newunicodechar{∃}{\ensuremath{\exists}}
\newunicodechar{∀}{\ensuremath{\forall}}
\newunicodechar{∛}{\ensuremath{\sqrt[3]{\,}}}
\newunicodechar{∜}{\ensuremath{\sqrt[4]{\,}}}
\newunicodechar{⃗}{\ensuremath{{}^{\to}}}
\newunicodechar{ⁿ}{\ifmmode^{n}\else\textsuperscript{\ensuremath{n}}\fi}
\newunicodechar{ˣ}{\ifmmode^{x}\else\textsuperscript{\ensuremath{x}}\fi}
\newunicodechar{ᵇ}{\ifmmode^{b}\else\textsuperscript{\ensuremath{b}}\fi}
\newunicodechar{ʳ}{\ifmmode^{r}\else\textsuperscript{\ensuremath{r}}\fi}
\newunicodechar{ᵘ}{\ifmmode^{u}\else\textsuperscript{\ensuremath{u}}\fi}
\newunicodechar{ʸ}{\ifmmode^{y}\else\textsuperscript{\ensuremath{y}}\fi}
\newunicodechar{ᵃ}{\ifmmode^{a}\else\textsuperscript{\ensuremath{a}}\fi}
\newunicodechar{ᵉ}{\ifmmode^{e}\else\textsuperscript{\ensuremath{e}}\fi}
\newunicodechar{ᵏ}{\ifmmode^{k}\else\textsuperscript{\ensuremath{k}}\fi}
\newunicodechar{ᵖ}{\ifmmode^{p}\else\textsuperscript{\ensuremath{p}}\fi}
\newunicodechar{ᴺ}{\ifmmode^{N}\else\textsuperscript{\ensuremath{N}}\fi}
\newunicodechar{ᶜ}{\ifmmode^{c}\else\textsuperscript{\ensuremath{c}}\fi}
\newunicodechar{ʰ}{\ifmmode^{h}\else\textsuperscript{\ensuremath{h}}\fi}
\newunicodechar{ᵠ}{\ifmmode^{q}\else\textsuperscript{\ensuremath{q}}\fi}
\newunicodechar{ₙ}{\ifmmode_{n}\else\textsubscript{\ensuremath{n}}\fi}
\newunicodechar{ₐ}{\ifmmode_{a}\else\textsubscript{\ensuremath{a}}\fi}
\newunicodechar{ᵢ}{\ifmmode_{i}\else\textsubscript{\ensuremath{i}}\fi}
\newunicodechar{ᵣ}{\ifmmode_{r}\else\textsubscript{\ensuremath{r}}\fi}
\newunicodechar{ₖ}{\ifmmode_{k}\else\textsubscript{\ensuremath{k}}\fi}
\newunicodechar{ᵦ}{\ifmmode_{\beta}\else\textsubscript{\ensuremath{\beta}}\fi}
\newunicodechar{ₓ}{\ifmmode_{x}\else\textsubscript{\ensuremath{x}}\fi}
\newfontfamily\popdisplay{ArchivoBlack-Regular.ttf}[Path=../../../fonts/]
\newfontfamily\poplogo{SpaceGrotesk-Bold.ttf}[Path=../../../fonts/]
\newfontfamily\popsemi{PlusJakartaSans-SemiBold.ttf}[Path=../../../fonts/]
\newfontfamily\popxbold{PlusJakartaSans-ExtraBold.ttf}[Path=../../../fonts/]
\newfontfamily\popxboldls{PlusJakartaSans-ExtraBold.ttf}[Path=../../../fonts/,LetterSpace=3.0]

\setlist[enumerate]{leftmargin=1.7em,itemsep=5pt,topsep=4pt}
\setlist[itemize]{leftmargin=1.7em,itemsep=4pt,topsep=4pt,label={\color{poporange}\textbullet}}
\setlength{\parindent}{0pt}
\setlength{\parskip}{5pt}
\renewcommand{\arraystretch}{1.25}

% pgfplots : style Pop par défaut
\pgfplotsset{
  every axis/.append style={axis line style={popink,line width=1pt},tick style={popink},
    grid style={popline,line width=0.5pt},label style={font=\small},tick label style={font=\footnotesize}},
  popcurve/.style={poporange,line width=1.8pt,no markers,smooth},
}
% Même style disponible dans un \draw TikZ simple (hors pgfplots)
\tikzset{popcurve/.style={poporange,line width=1.8pt}}
\tikzset{popgrid/.style={popline,very thin}}
\colorlet{popcurve}{poporange} % le nom du style est parfois utilisé comme couleur

% ============================================================
% Pill / squiggle
% ============================================================
\newcommand{\poppill}[3]{%
  \tikz[baseline=-0.55ex]\node[draw=popink,line width=1.2pt,rounded corners=2.6pt,%
    fill=#2,inner xsep=6.5pt,inner ysep=3pt]{%
    \popxboldls\fontsize{7}{7}\selectfont\color{#3}\MakeUppercase{#1}};%
}
\newcommand{\popsquiggle}[1]{%
  \tikz[baseline=-0.4ex]{
    \draw[#1,line width=2.6pt,line cap=round]
      plot [smooth] coordinates {(0,0.09) (0.34,0.01) (0.68,0.21) (1.02,0.01) (1.36,0.10)};
  }%
}
% Mot-clé surligné (marqueur orange sous le texte)
\newcommand{\cle}[1]{{\popxbold\color{popdark}#1}}

% ============================================================
% En-tête / pied
% ============================================================
\pagestyle{fancy}
\fancyhf{}
\renewcommand{\headrulewidth}{0pt}
\renewcommand{\footrulewidth}{0pt}
\lhead{\raisebox{-0.15ex}{\poplogo\fontsize{13}{13}\selectfont\color{popink}OnBuch\color{poporange}+}}
\rhead{\poppill{\DOCMATIERE\ \textbullet\ \DOCNIVEAU\ \textbullet\ Leçon \DOCLECON}{white}{popink}}
\lfoot{\popxbold\fontsize{7.6}{7.6}\selectfont\color{popmuted}\MakeUppercase{Cours signé OnBuch+}\ \textbullet\ {\popsemi\DOCTITRE}}
\rfoot{\tikz[baseline=-0.6ex]\node[rounded corners=8.5pt,fill=popink,minimum height=17pt,minimum width=17pt,inner xsep=5pt,inner sep=0]{\popxbold\fontsize{8}{8}\selectfont\color{popcream}\thepage\,/\,\pageref*{LastPage}};}

% ============================================================
% Titres
% ============================================================
\newcommand{\chaptitre}[2]{%
  \begin{tcolorbox}[enhanced,colback=poporange,colframe=popink,boxrule=2.6pt,arc=11pt,
      drop shadow=popink,left=15pt,right=15pt,top=12pt,bottom=11pt,before skip=3pt,after skip=18pt]
    \poppill{#2}{popink}{popcream}\par\vspace{9pt}
    {\raggedright\hyphenpenalty=10000\exhyphenpenalty=10000\popdisplay\fontsize{22}{24}\selectfont\color{popcream}\MakeUppercase{#1}\par}
  \end{tcolorbox}
}
% Section numérotée : pastille orange avec le numéro + titre Archivo
\newcounter{popsec}
\newcommand{\coursec}[1]{%
  \stepcounter{popsec}\setcounter{popsub}{0}%
  \par\vspace{16pt}\noindent
  \begin{minipage}{\linewidth}
  \tikz[baseline=-0.7ex]\node[draw=popink,line width=1.6pt,fill=poporange,rounded corners=4pt,minimum size=22pt,inner sep=0]{\popdisplay\fontsize{12}{12}\selectfont\color{popcream}\thepopsec};%
  \hspace{8pt}{\popdisplay\fontsize{15}{17}\selectfont\color{popink}\MakeUppercase{#1}}\par
  \vspace{4pt}\noindent{\color{poporange}\rule{36pt}{3.6pt}}
  \end{minipage}\par\nopagebreak\vspace{8pt}
}
\newcounter{popsub}
\newcommand{\courssub}[1]{%
  \stepcounter{popsub}%
  \par\vspace{9pt}\noindent{\popxbold\fontsize{11.5}{13}\selectfont\color{popink}{\color{poporange}\thepopsec.\thepopsub}\hspace{6pt}#1}\par\nopagebreak\vspace{4pt}
}

% ============================================================
% Boîtes pédagogiques — même recette : fond blanc, bordure épaisse colorée,
% ombre dure encre, badge plein. Titre optionnel [#1].
% ============================================================
\tcbset{popbase/.style={breakable,enhanced,colback=white,boxrule=2.3pt,arc=9pt,
  shadow={3pt}{-3pt}{0pt}{popink},left=14pt,right=14pt,top=11pt,bottom=11pt,before skip=11pt,after skip=15pt,
  fontupper=\popsemi\fontsize{10.6}{14.5}\selectfont\color{popink},
  fontlower=\popsemi\fontsize{10.6}{14.5}\selectfont\color{popink}}}
% Coche dessinée (le glyphe ✓ n'existe pas dans les polices du document)
\newcommand{\popcheck}[1]{\tikz[baseline=-0.5ex]\draw[#1,line width=1.6pt,line cap=round,line join=round](-0.13,0.0)--(-0.04,-0.09)--(0.13,0.1);}
\providecommand{\coche}{\popcheck{popgreen}}
\providecommand{\resultat}[1]{\par\smallskip\noindent\fcolorbox{popgreen}{popgreenL}{\parbox{\dimexpr\linewidth-2\fboxsep-2\fboxrule\relax}{\textbf{Résultat.} #1}}\par\smallskip}
\newcommand{\popboxhead}[3]{\poppill{#1}{#2}{white}\ifx\relax#3\relax\else\hspace{6pt}{\popxbold\fontsize{10.6}{12}\selectfont\color{popink}#3}\fi\par\vspace{7pt}}

\newtcolorbox{aretenir}[1][]{popbase,colframe=popgreen,before upper={\popboxhead{À retenir}{popgreen}{#1}}}
\newtcolorbox{exemplebox}[1][]{popbase,colframe=poporange,before upper={\popboxhead{Exemple}{poporange}{#1}}}
\newtcolorbox{definition}[1][]{popbase,colframe=popblue,before upper={\popboxhead{Définition}{popblue}{#1}}}
\newtcolorbox{propriete}[1][]{popbase,colframe=poppurple,before upper={\popboxhead{Propriété}{poppurple}{#1}}}
\newtcolorbox{methode}[1][]{popbase,colframe=popgold,before upper={\popboxhead{Méthode}{popgold}{#1}}}
\newtcolorbox{attention}[1][]{popbase,colframe=poppink,before upper={\popboxhead{Attention}{poppink}{#1}}}
\newtcolorbox{experience}[1][]{popbase,colframe=popblue,colback=popblueL,before upper={\popboxhead{Expérience}{popblue}{#1}}}
% « Le savais-tu ? » : carte violette pleine (contexte, culture, Cameroun)
\newtcolorbox{savaistu}[1][]{popbase,colback=poppurple,colframe=popink,
  fontupper=\popsemi\fontsize{10.4}{14.2}\selectfont\color{white},
  before upper={\poppill{Le savais-tu\,?}{white}{poppurple}\ifx\relax#1\relax\else\hspace{6pt}{\popxbold\color{white}#1}\fi\par\vspace{7pt}}}
% Exercice résolu : énoncé en haut, \tcblower puis la solution sur fond crème
\newcounter{popexo}\newcounter{popexr}
\newtcolorbox{exoresolu}[1][]{popbase,colframe=poporange,colbacklower=popcream,
  segmentation style={popink,line width=1.2pt,dashed},
  before upper={\stepcounter{popexr}\popboxhead{Exercice résolu \thepopexr}{poporange}{#1}},
  before lower={\poppill{Solution}{popink}{popcream}\par\vspace{6pt}}}
% Exercice (non résolu en place) — la correction est dans la partie Corrigés
\newtcolorbox{exercice}[1][]{popbase,colframe=popink,
  before upper={\stepcounter{popexo}\popboxhead{Exercice \thepopexo}{popink}{#1}}}
% Corrigé
\newtcolorbox{corrige}[1][]{popbase,colframe=popgreen,colback=popgreenL,
  before upper={\popboxhead{Corrigé}{popgreen}{#1}}}
% Objectifs / prérequis / bilan
\newtcolorbox{objectifs}{popbase,colframe=popink,colback=poporangeL,
  before upper={\popboxhead{Objectifs de la leçon}{popink}{}}}
\newtcolorbox{prerequis}{popbase,colframe=popink,colback=popblueL,
  before upper={\popboxhead{Prérequis}{popblue}{}}}
\newtcolorbox{bilan}{popbase,colframe=popink,colback=white,boxrule=2.8pt,
  before upper={\popboxhead{Fiche bilan}{poporange}{L'essentiel à connaître pour l'examen}}}
% Figure encadrée avec légende
\newtcolorbox{popfigure}[1][]{enhanced,breakable=false,colback=white,colframe=popline,boxrule=1.4pt,arc=8pt,
  left=10pt,right=10pt,top=10pt,bottom=8pt,before skip=10pt,after skip=12pt,halign=center,
  lower separated=false,halign lower=center,
  fontlower=\popsemi\fontsize{9}{11.5}\selectfont\color{popmuted},
  #1}
\newcommand{\legende}[1]{\par\vspace{6pt}{\popsemi\fontsize{9}{11.5}\selectfont\color{popmuted}#1}}

% ============================================================
% Signature OnBuch+
% ============================================================
\newcommand{\poplogoinline}[1]{{\poplogo\fontsize{#1}{#1}\selectfont\color{popink}OnBuch\color{poporange}+}}
\newcommand{\signatureonbuch}{%
  \begin{tcolorbox}[enhanced,colback=popink,colframe=popink,boxrule=2.2pt,arc=10pt,
      drop shadow=poporange,left=14pt,right=14pt,top=10pt,bottom=10pt,before skip=0pt,after skip=0pt]
    \tikz[baseline=-0.7ex]\node[circle,fill=popgreen,draw=popcream,line width=1.3pt,minimum size=22pt,inner sep=0]{\popcheck{white}};%
    \hspace{9pt}\begin{minipage}[c]{0.62\linewidth}
      {\popxbold\fontsize{12}{13}\selectfont\color{popcream}Signé\ \textbullet\ {\poplogo OnBuch\color{poporange}+}}\par\vspace{2pt}
      {\raggedright\popsemi\fontsize{8}{10.5}\selectfont\color{popline}\MakeUppercase{Équipe pédagogique OnBuch+}\\\MakeUppercase{Conforme au programme officiel MINESEC}\par}
    \end{minipage}\hfill
    {\poplogo\fontsize{22}{22}\selectfont\color{popcream}OnBuch\color{poporange}+}
  \end{tcolorbox}%
}

% ============================================================
% Page de garde
% #1 titre · #2 matière · #3 niveau · #4 module · #5 n° leçon · #6 durée ·
% #7 description / objectifs (texte ou itemize)
% ============================================================
\newcommand{\pagegarde}[7]{%
  \thispagestyle{empty}
  \newgeometry{a4paper,margin=1.7cm,top=1.6cm,bottom=1.4cm}%
  \begin{tikzpicture}[remember picture,overlay]
    \fill[poporange!18] ([xshift=-3.2cm,yshift=-3.6cm]current page.north east) circle (4.6cm);
    \fill[poppurple!14] ([xshift=2.4cm,yshift=7.6cm]current page.south west) circle (3.8cm);
  \end{tikzpicture}%
  {\poplogo\fontsize{30}{30}\selectfont\color{popink}OnBuch\color{poporange}+}\hfill
  \raisebox{0.6ex}{\poppill{Cours signé \textbullet\ Premium}{popink}{popcream}}\par
  \vspace{1pt}\popsquiggle{poppurple}\par
  \vspace{8pt}
  \begin{tcolorbox}[enhanced,colback=poporange,colframe=popink,boxrule=2.8pt,arc=13pt,
      drop shadow=popink,left=18pt,right=18pt,top=12pt,bottom=13pt,after skip=10pt]
    \poppill{#2 \textbullet\ #3}{popink}{popcream}\hspace{5pt}\poppill{#4}{white}{popink}\par\vspace{8pt}
    {\popdisplay\fontsize{14}{14}\selectfont\color{popink}LEÇON #5}\par\vspace{4pt}
    {\raggedright\hyphenpenalty=10000\exhyphenpenalty=10000\popdisplay\fontsize{25}{27}\selectfont\color{popcream}\MakeUppercase{#1}\par}
    \vspace{8pt}
    {\popxbold\fontsize{9.5}{9.5}\selectfont\color{popink}\MakeUppercase{Durée conseillée : #6}}
  \end{tcolorbox}
  \begin{tcolorbox}[enhanced,colback=white,colframe=popink,boxrule=2.2pt,arc=10pt,
      drop shadow=popink,left=16pt,right=16pt,top=10pt,bottom=10pt,after skip=8pt]
    \poppill{Dans cette leçon}{poppurple}{white}\par\vspace{5pt}
    \popsemi\fontsize{9.3}{12.6}\selectfont\color{popink}#7
  \end{tcolorbox}
  \noindent\begin{minipage}[t]{0.487\linewidth}
    \begin{tcolorbox}[enhanced,colback=popgreen,colframe=popink,boxrule=2.2pt,arc=10pt,
        drop shadow=popink,left=11pt,right=11pt,top=10pt,bottom=10pt,height=3.1cm,valign=center]
      \tcbox[colback=white,colframe=popink,boxrule=1.3pt,arc=5pt,boxsep=0pt,nobeforeafter,
        left=3pt,right=3pt,top=3pt,bottom=3pt,valign=center]{\qrcode[height=1.9cm]{https://play.google.com/store/apps/details?id=com.luvvix.onbuch&hl=fr}}%
      \hspace{8pt}\begin{minipage}[c]{0.5\linewidth}
        {\popdisplay\fontsize{11}{12.5}\selectfont\color{white}\MakeUppercase{Télécharge OnBuch}}\par\vspace{4pt}
        {\raggedright\popsemi\fontsize{8.4}{11}\selectfont\color{white}Gratuit \textbullet\ Google Play\\Tuteur IA, annales, concours, résultats\par}
      \end{minipage}
    \end{tcolorbox}
  \end{minipage}\hfill
  \begin{minipage}[t]{0.487\linewidth}
    \begin{tcolorbox}[enhanced,colback=poppurple,colframe=popink,boxrule=2.2pt,arc=10pt,
        drop shadow=popink,left=11pt,right=11pt,top=10pt,bottom=10pt,height=3.1cm,valign=center]
      \tikz[baseline=-0.7ex]\node[circle,fill=white,draw=popink,line width=1.3pt,minimum size=1.6cm,inner sep=0]{\popdisplay\fontsize{24}{24}\selectfont\color{poppurple}+};%
      \hspace{8pt}\begin{minipage}[c]{0.6\linewidth}
        {\popdisplay\fontsize{11}{12.5}\selectfont\color{white}\MakeUppercase{Passe à OnBuch+}}\par\vspace{4pt}
        {\raggedright\popsemi\fontsize{8.4}{11}\selectfont\color{white}Tous les cours signés, Tuteur IA illimité, fiches PDF — onglet OnBuch+ de l'app\par}
      \end{minipage}
    \end{tcolorbox}
  \end{minipage}
  \vspace{8pt}
  \popcontenu
  \vfill
  \signatureonbuch
  \restoregeometry
  \newpage
}

% Rangée « Ce cours contient » (page de garde)
\newcommand{\poptile}[3]{% #1 couleur #2 gros texte #3 légende
  \begin{tcolorbox}[enhanced,colback=#1,colframe=popink,boxrule=2pt,arc=9pt,shadow={2.5pt}{-2.5pt}{0pt}{popink},
      left=6pt,right=6pt,top=9pt,bottom=9pt,halign=center,valign=center,height=1.75cm,width=\linewidth,nobeforeafter]
    {\popdisplay\fontsize{12.5}{13}\selectfont\color{white}\MakeUppercase{#2}}\par\vspace{3pt}
    {\centering\popsemi\fontsize{7.8}{9.5}\selectfont\color{white}#3\par}
  \end{tcolorbox}}
\newcommand{\popcontenu}{%
  \par\noindent{\popxboldls\fontsize{7.5}{7.5}\selectfont\color{popmuted}CE COURS SIGNÉ CONTIENT}\par\vspace{7pt}
  \noindent\begin{minipage}[t]{0.235\linewidth}\poptile{popblue}{Cours}{complet, illustré, pas à pas}\end{minipage}\hfill
  \begin{minipage}[t]{0.235\linewidth}\poptile{poporange}{Méthodes}{et exercices résolus}\end{minipage}\hfill
  \begin{minipage}[t]{0.235\linewidth}\poptile{poppink}{Exercices}{type examen + corrigés}\end{minipage}\hfill
  \begin{minipage}[t]{0.235\linewidth}\poptile{popgreen}{Fiche bilan}{l'essentiel à retenir}\end{minipage}\par}

% Sommaire
\newtcolorbox{sommaire}{enhanced,colback=white,colframe=popink,boxrule=2.2pt,arc=10pt,
  drop shadow=popink,left=18pt,right=18pt,top=13pt,bottom=13pt,before skip=6pt,after skip=20pt,
  before upper={\poppill{Sommaire}{poppurple}{white}\par\vspace{9pt}\popsemi\fontsize{10.6}{14.5}\selectfont\color{popink}}}

% Signature de fin de cours — poussée tout en bas de la dernière page
\newcommand{\findecours}{%
  \par\vfill\nopagebreak
  \begin{center}\popsquiggle{poporange}\end{center}\vspace{4pt}
  \signatureonbuch
}
\AtBeginDocument{\def\llbracket{\mathopen{[\mkern-3.5mu[}}\def\rrbracket{\mathclose{]\mkern-3.5mu]}}} % intervalles d'entiers (glyphe ⟦ absent de la police)
% Glyphes absents de Fira Math : redéfinis à partir de glyphes disponibles
\AtBeginDocument{%
  \def\setminus{\mathbin{\backslash}}\let\smallsetminus\setminus
  \def\vdots{\mathord{\vbox{\baselineskip3.2pt\lineskiplimit0pt\kern4pt\hbox{.}\hbox{.}\hbox{.}}}}%
  \def\ddots{\mathinner{\mkern1mu\raise6.5pt\hbox{.}\mkern2mu\raise3.6pt\hbox{.}\mkern2mu\raise0.7pt\hbox{.}\mkern1mu}}%
  \def\triangle{\mathord{\scalebox{0.9}{$\symup{\Delta}$}}}%
  \def\nabla{\mathord{\raisebox{\depth}{\scalebox{1}[-1]{$\symup{\Delta}$}}}}%
  \def\checkmark{\popcheck{popdark}}%
  \def\star{\mathbin{\ast}}\def\bigstar{\mathord{\ast}}%
  \def\diamond{\mathbin{\scalebox{0.6}{$\Diamond$}}}%
  \def\bigcup{\mathop{\mathchoice{\raisebox{-0.3em}{\scalebox{1.7}{$\cup$}}}{\scalebox{1.25}{$\cup$}}{\cup}{\cup}}\displaylimits}%
  \def\bigcap{\mathop{\mathchoice{\raisebox{-0.3em}{\scalebox{1.7}{$\cap$}}}{\scalebox{1.25}{$\cap$}}{\cap}{\cap}}\displaylimits}%
  \def\lesseqqgtr{\mathrel{\lessgtr}}\def\bigoplus{\mathop{\oplus}\displaylimits}%
}
\providecommand{\ssi}{si et seulement si} % abréviation fréquente
\AtBeginDocument{\providecommand{\wideparen}{\overparen}} % arc de cercle

%% FIGFIX-BEGIN v5 — correctifs globaux des figures (docs/onbuch-generation/tools/figfix)
\usepackage{adjustbox}\usepackage{etoolbox}\tcbuselibrary{raster}
% toute figure TikZ/pgfplots plus large que la ligne est réduite à la largeur disponible
\BeforeBeginEnvironment{tikzpicture}{\begin{adjustbox}{max width=\linewidth}}
\AfterEndEnvironment{tikzpicture}{\end{adjustbox}}
% légendes pgfplots lisibles par-dessus les courbes
\pgfplotsset{every axis/.append style={legend style={fill=white,fill opacity=0.92,text opacity=1,draw=black!15,inner sep=3pt}}}
% étiquettes d'arêtes (syntaxe edge["..."]) avec fond blanc
\tikzset{every edge quotes/.append style={fill=white,inner sep=1.2pt}}
%% FIGFIX-END
\begin{document}

\pagegarde{Lutte contre le VIH/Sida}{SVT}{1ère TI}{Module 2 : L'Éducation à la Santé}{NCT04}{6 h}{Cette leçon aborde la biologie du VIH, sa transmission, son évolution naturelle et les stratégies de prévention et de prise en charge. Elle vise à former des citoyens responsables, capables d'argumenter scientifiquement pour le dépistage volontaire et l'observance thérapeutique.
\vspace{4pt}
\begin{multicols}{2}\small
\begin{itemize}
\item Décrire la structure du VIH et expliquer les grandes étapes de son cycle de réplication dans la cellule CD4.
\item Lister les liquides biologiques contaminants et expliquer les trois modes de transmission (sexuel, sanguin, materno-fœtal) en les reliant aux pratiques à risque.
\item Différencier les moyens de prévention primaire (préservatifs, PrEP, Tasp, éducation) et secondaire (PEP, dépistage).
\item Reconstituer l'évolution naturelle de l'infection : phase aiguë (séroconversion), phase chronique asymptomatique, phase Sida (maladies opportunistes) à l'aide des marqueurs CD4 et charge virale.
\item Interpréter des courbes d'évolution de la charge virale et des lymphocytes T CD4 sous traitement antirétroviral (ARV) ou sans traitement.
\item Expliquer le principe des tests de dépistage (recherche anticorps/antigène p24) et la fenêtre sérologique.
\end{itemize}\end{multicols}}

\begin{sommaire}
\begin{multicols}{2}
\begin{enumerate}
\item 1. Le VIH : structure et cycle infectieux
\item 2. Modes de transmission et prévention
\item 3. Évolution naturelle de l'infection et marqueurs biologiques
\item 4. Dépistage : principes, stratégies et algorithmes
\item 5. Prise en charge : Traitements Antirétroviraux (ARV) et suivi
\item 6. Aspects psychosociaux, droits et riposte nationale
\item Activité d'intégration
\item Méthodes et exercices résolus
\item Exercices
\item Corrigés des exercices
\item Fiche bilan
\end{enumerate}
\end{multicols}
\end{sommaire}

\begin{objectifs}
\begin{itemize}
\item Décrire la structure du VIH et expliquer les grandes étapes de son cycle de réplication dans la cellule CD4.
\item Lister les liquides biologiques contaminants et expliquer les trois modes de transmission (sexuel, sanguin, materno-fœtal) en les reliant aux pratiques à risque.
\item Différencier les moyens de prévention primaire (préservatifs, PrEP, Tasp, éducation) et secondaire (PEP, dépistage).
\item Reconstituer l'évolution naturelle de l'infection : phase aiguë (séroconversion), phase chronique asymptomatique, phase Sida (maladies opportunistes) à l'aide des marqueurs CD4 et charge virale.
\item Interpréter des courbes d'évolution de la charge virale et des lymphocytes T CD4 sous traitement antirétroviral (ARV) ou sans traitement.
\item Expliquer le principe des tests de dépistage (recherche anticorps/antigène p24) et la fenêtre sérologique.
\item Argumenter en faveur du dépistage volontaire, de l'accès précoce aux ARV et de l'observance pour atteindre l'indétectabilité (= intransmissibilité).
\item Analyser des données épidémiologiques simples (prévalence, incidence, cascade de prise en charge) pour le Cameroun.
\end{itemize}
\end{objectifs}

\begin{prerequis}
\begin{itemize}
\item Structure d'un virus (capside, enveloppe, acide nucléique) et différence avec une bactérie (Seconde).
\item Rôle des lymphocytes T CD4 (coordinateurs de l'immunité adaptative) et notion d'immunité cellulaire/humorale (1ère, début année).
\item Notion de transcription inverse (ADN -> ARN) et dogme central de la biologie moléculaire (1ère).
\item Principes de base de la transmission des IST et moyens de prévention (Seconde/1ère).
\item Lecture et interprétation de graphiques simples (courbes, histogrammes) et calculs de pourcentages.
\end{itemize}
\end{prerequis}

\newpage

\coursec{Le VIH : structure et cycle infectieux}

Au Cameroun, malgré les progrès considérables accomplis ces deux dernières décennies, le VIH demeure un enjeu de santé publique majeur. La prévalence estimée à 2,7\,\% chez les adultes de 15 à 49 ans (EDS-Cameroun 2018) masque des disparités régionales et par genre : les jeunes filles de 15 à 24 ans sont trois à quatre fois plus touchées que les garçons du même âge. Chaque jour, dans nos centres de santé de district comme dans les hôpitaux de référence de Yaoundé ou Douala, de jeunes Camerounais découvrent leur séropositivité souvent tardivement, au stade Sida, par manque d'information, peur du dépistage ou stigmatisation. Comment ce virus microscopique parvient-il à détruire notre système de défense, le plus sophistiqué de la biologie ? Pourquoi l'utilisation correcte et systématique du préservatif, associée aujourd'hui à la PrEP (prophylaxie pré-exposition), reste-t-elle notre meilleure arme préventive ? Comment un biologiste interprète-t-il une prise de sang montrant une charge virale à 150\,000 copies/mL et un taux de CD4 à 180/mm³ pour décider de l'urgence thérapeutique ? C'est pour répondre à ces questions vitales que nous devons d'abord comprendre l'ennemi : sa structure, sa stratégie de réplication et ses failles.

\courssub{Classification et structure du virion VIH-1}

\begin{definition}[Rétrovirus et Transcriptase Inverse]
Un \cle{rétrovirus} est un virus à ARN monobrin de polarité positive (+) dont le cycle de réplication implique une étape obligatoire de \cle{transcription inverse} : son génome ARN est rétrotranscrit en ADN double brin par une enzyme virale spécifique, la \cle{transcriptase inverse} (TI), avant d'être intégré dans le génome de la cellule hôte. La famille des \emph{Retroviridae} comprend deux sous-familles (\emph{Orthoretrovirinae} et \emph{Spumaretrovirinae}) ; le VIH appartient au genre \cle{Lentivirus} (lenteurs de l'évolution clinique), caractérisé par une longue période de latence clinique, un tropisme pour les cellules du système immunitaire et nerveux, et une pathogénicité forte.
\end{definition}

Le VIH existe sous deux types : le \cle{VIH-1} (pandémique, responsable de l'immense majorité des infections dans le monde et au Cameroun) et le \cle{VIH-2} (endémique en Afrique de l'Ouest, moins transmissible, moins pathogène, naturellement résistant aux INNRTI). Nous nous concentrons ici sur la structure du VIH-1, le type circulant au Cameroun (principalement du groupe M, sous-types A, G, CRF02\_AG).

\begin{popfigure}
\begin{tikzpicture}[scale=1.2, every node/.style={font=\small}]
    % Enveloppe
    \draw[popblue, very thick] (0,0) circle (3cm);
    % Glycoprotéines de surface (gp120/gp41)
    \foreach \angle/\label in {30/gp120, 150/gp120, 210/gp120, 330/gp120, 90/gp41, 270/gp41} {
        \draw[poporange, thick] (\angle:3cm) -- (\angle:3.6cm);
        \node[poporange, anchor=center] at (\angle:3.85cm) {\label};
    }
    % Légende enveloppe
    \node[popblue, align=center] at (0, 3.8) {Enveloppe lipidique\\ (dérivée membrane cellulaire)};
    
    % Matrice (p17) - ligne fine sous enveloppe
    \draw[popmuted, dashed] (0,0) circle (2.7cm);
    \node[popmuted, rotate=45] at (2.2, 1.5) {Matrice p17};
    
    % Capside (p24) - cône tronqué simplifié en cercle
    \draw[popdark, very thick] (0,0) circle (1.8cm);
    \node[popdark] at (0, -1.2) {Capside (p24)};
    
    % ARN génomique (deux brins)
    \draw[poppink, thick, decorate, decoration={snake, amplitude=0.5mm, segment length=2mm}] (-0.6,0.3) -- (0.6,0.3);
    \draw[poppink, thick, decorate, decoration={snake, amplitude=0.5mm, segment length=2mm}] (-0.6,-0.3) -- (0.6,-0.3);
    \node[poppink] at (0, 0.8) {2 ARN (+) identiques};
    \node[poppink] at (0, -0.8) {(\textasciitilde 9,7 kb chacun)};
    
    % Enzymes : Transcriptase Inverse (p66/p51), Intégrase (p32), Protéase (p10)
    \node[popgreen, circle, draw, thick, minimum size=0.8cm] (RT) at (-1.2, 1.0) {TI};
    \node[popgreen, circle, draw, thick, minimum size=0.8cm] (IN) at (1.2, 1.0) {IN};
    \node[popgreen, circle, draw, thick, minimum size=0.8cm] (PR) at (0, -1.5) {PR};
    
    \draw[popgreen, ->, thick] (RT) -- ++(-0.5,0.5) node[left, popgreen, align=center]{Transcriptase\\ Inverse (p66/p51)};
    \draw[popgreen, ->, thick] (IN) -- ++(0.5,0.5) node[right, popgreen] {Intégrase (p32)};
    \draw[popgreen, ->, thick] (PR) -- ++(0,-0.7) node[below, popgreen] {Protéase (p10)};
    
    % Nucléoprotéines p7/p9 (autour ARN)
    \node[poppurple, font=\tiny] at (-1.5, -0.5) {Nucléoprotéines p7/p9};
\end{tikzpicture}
\legende{Structure du virion VIH-1 (coupe transversale schématique). L'enveloppe lipidique porte les glycoprotéines gp120 (fixation) et gp41 (fusion). La capside protéique (p24) protège les deux brins d'ARN génomique associés aux nucléoprotéines et aux trois enzymes virales essentielles : Transcriptase Inverse (TI), Intégrase (IN), Protéase (PR).}
\end{popfigure}

\noindent Le virion VIH-1 est une particule sphérique d'environ 100 à 120\,nm de diamètre. De l'extérieur vers l'intérieur, on distingue :
\begin{enumerate}
    \item \textbf{L'enveloppe lipidique} : bicouche phospholipique dérivée de la membrane plasmique de la cellule hôte lors du bourgeonnement. Elle porte les \cle{glycoprotéines d'enveloppe} (Env), organisées en trimères (gp160 clivé en \textbf{gp120} et \textbf{gp41}). Le gp120 (sous-unité de surface) assure la reconnaissance des récepteurs cellulaires ; le gp41 (sous-unité transmembranaire) médie la fusion des membranes.
    \item \textbf{La matrice (protéine p17)} : protéine structurelle tapissant la face interne de l'enveloppe, assurant la stabilité du virion et le transport du complexe de pré-intégration vers le noyau.
    \item \textbf{La capside (protéine p24)} : coque protéique conique (symétrie icosaédrique déformée) formée par l'assemblage d'environ 1500 copies de p24. Elle protège le noyau ribonucléoprotéique et joue un rôle clé dans le trafic intracellulaire et l'évasion de l'immunité innée (détection par cGAS).
    \item \textbf{Le noyau ribonucléoprotéique} : contient \textbf{deux brins d'ARN monobrin de polarité positive} (ARNg, \textasciitilde 9,7\,kb chacun), non covalents mais appariés par leur région 5' (dimère). Chaque brin est enrobé de protéines nucléocapsidiques \textbf{p7/p9} (domaines zinc-doigt) qui protègent l'ARN et agissent comme chaperons lors de la transcription inverse. Trois enzymes virales y sont associées en quelques copies : la \textbf{Transcriptase Inverse (TI, p66/p51)}, l'\textbf{Intégrase (IN, p32)} et la \textbf{Protéase (PR, p10)}.
\end{enumerate}

\begin{attention}[Transcriptase Inverse $\neq$ ADN Polymérase cellulaire]
\emph{Erreur fréquente au Probatoire :} Confondre la transcriptase inverse virale et l'ADN polymérase cellulaire (réplication de l'ADN hôte).
\begin{itemize}
    \item \textbf{Transcriptase Inverse (TI)} : Enzyme \emph{virale}, apportée par le virion. Utilise un \emph{ARN} comme matrice pour synthétiser de l'\emph{ADN}. \textbf{N'a pas d'activité de relecture (proofreading 3'$\to$5' exonucléase)} $\rightarrow$ taux d'erreur élevé ($\sim 3 \times 10^{-5}$ mutations/base/cycle).
    \item \textbf{ADN Polymérase cellulaire} : Enzyme \emph{hôte}. Utilise de l'\emph{ADN} comme matrice pour synthétiser de l'\emph{ADN}. Possède une activité de relecture $\rightarrow$ fidélité élevée.
\end{itemize}
L'absence de relecture de la TI est la cause majeure de l'extrême variabilité génétique du VIH (quasi-espèces), source d'échappement immunitaire et de résistance aux antirétroviraux.
\end{attention}

\courssub{Cible cellulaire : le lymphocyte T CD4$^+$ et les co-récepteurs}

Le VIH ne se réplique productivement que dans les cellules exprimant le récepteur \cle{CD4} (glycoprotéine de 55\,kDa, membre de la superfamille des immunoglobulines), marqueur majeur des \cle{lymphocytes T auxiliaires (Th)}. Mais le CD4 seul ne suffit pas : l'entrée requiert un \cle{co-récepteur} chimiocinique à 7 domaines transmembranaires (famille GPCR).
\begin{itemize}
    \item \textbf{CCR5 (R5)} : Exprimé sur les lymphocytes T mémoire effecteurs, macrophages, cellules dendritiques. Utilisé par les souches \cle{viraux R5-tropiques} (majoritaires au stade primaire et asymptomatique, transmissibles sexuellement et materno-fœtalement).
    \item \textbf{CXCR4 (X4)} : Exprimé sur les lymphocytes T naïfs, certains Th. Utilisé par les souches \cle{X4-tropiques} (apparaissent souvent au stade Sida, plus cytopathiques, syncytium-induisantes).
    \item \textbf{Double tropisme (R5X4)} : Souches utilisant les deux.
\end{itemize}

\begin{propriete}[Réservoirs viraux et latence]
Le VIH ne se réplique activement que dans les \textbf{cellules CD4$^+$ activées} (facteurs de transcription NF-$\kappa$B, NFAT présents). Dans les \textbf{cellules CD4$^+$ au repos (mémoire centrale, cellules souches hématopoïétiques)}, le virus peut intégrer son ADN (ADN proviral) mais la transcription est bloquée : c'est la \cle{latence virale}. Ces cellules à longue durée de vie (mois à années) constituent les \cle{réservoirs viraux}. Elles sont invisibles au système immunitaire et insensibles aux ARV (qui ciblent la réplication active). C'est l'obstacle majeur à la \emph{guérison stérilisante} (éradication totale), d'où la nécessité d'un traitement à vie actuel.
\end{propriete}

\courssub{Cycle de réplication : les 7 étapes (FTITABM)}

Le cycle infectieux est un enchaînement précis de 7 étapes, chacune étant une cible thérapeutique potentielle. Pour le mémoriser, utilisez la méthode mnémotechnique suivante :

\begin{methode}[Mémoriser le cycle : FTITABM]
\emph{Les 7 lettres pour les 7 étapes :}
\begin{enumerate}
    \item \textbf{F}ixation / \textbf{F}usion (Entrée)
    \item \textbf{T}ranscription inverse (ARN $\to$ ADN)
    \item \textbf{I}ntégration (ADN proviral)
    \item \textbf{T}ranscription / Traduction (Expression génique)
    \item \textbf{A}ssemblage (Formation particule immature)
    \item \textbf{B}ourgeonnement (Sortie / Acquisition enveloppe)
    \item \textbf{M}aturation (Clivage Protéase $\to$ Virion infectieux)
\end{enumerate}
\end{methode}

\begin{popfigure}
\begin{tikzpicture}[scale=0.9, every node/.style={font=\small, align=center}, node distance=1.2cm and 1.5cm]
    % Styles
    \tikzset{
        stepbox/.style={rectangle, draw=popblue, thick, fill=popblueL, rounded corners, minimum height=1.2cm, minimum width=2.8cm},
        arrow/.style={->, >=Stealth, thick, popdark},
        loc/.style={rectangle, draw=poporange, thick, fill=poporangeL, rounded corners, minimum height=0.8cm, minimum width=2.5cm, font=\footnotesize}
    }
    
    % Localisation Cytoplasme (gauche) et Noyau (droite)
    \node[loc] (cyto) at (-4, -0.5) {CYTOPLASME};
    \node[loc] (nuc) at (4, -0.5) {NOYAU};
    
    % Étapes
    \node[stepbox, align=center] (E1) at (-4, 2.5){1. FIXATION / FUSION\\ gp120 $\to$ CD4/CCR5 $\to$ gp41 fusion};
    \node[stepbox, align=center] (E2) at (-4, 1.0){2. TRANSCRIPTION INVERSE\\ ARN $\xrightarrow{TI}$ ADN db (Provirus)};
    \node[stepbox, align=center] (E3) at (4, 1.0){3. INTÉGRATION\\ ADN viral $\xrightarrow{IN}$ ADN hôte};
    \node[stepbox, align=center] (E4) at (4, -0.5){4. TRANSCRIPTION / TRADUCTION\\ ARNm $\to$ Protéines (Gag, Pol, Env...)};
    \node[stepbox, align=center] (E5) at (-4, -2.0){5. ASSEMBLAGE\\ Pr55Gag + ARN + Enzymes $\to$ Particule immature};
    \node[stepbox, align=center] (E6) at (-4, -3.5){6. BOURGEONNEMENT\\ Budding $\to$ Acquisition enveloppe};
    \node[stepbox, align=center] (E7) at (4, -3.5){7. MATURATION\\ Protéase clive Pr55Gag $\to$ Virion mature infectieux};
    
    % Flèches
    \draw[arrow] (E1) -- (E2);
    \draw[arrow] (E2) -- node[midway, above, sloped, font=\tiny] {Transport noyau (p17, IN)} (E3);
    \draw[arrow] (E3) -- (E4);
    \draw[arrow] (E4) -- node[midway, above, sloped, font=\tiny] {Export ARNm (Rev/CRM1)} (E5);
    \draw[arrow] (E5) -- (E6);
    \draw[arrow] (E6) -- node[midway, below, sloped, font=\tiny] {Maturation extracellulaire / surface} (E7);
    
    % Légendes localisation (remplace les accolades TikZ non supportées)
    \draw[popmuted, thick] (-5.8, 3.2) -- (-5.8, -4.2);
    \node[popmuted, rotate=90, anchor=south] at (-6.3, -0.5) {Majoritairement Cytoplasme};
    \draw[popmuted, thick] (5.8, 1.7) -- (5.8, -4.2);
    \node[popmuted, rotate=-90, anchor=north] at (6.3, -1.2) {Noyau (Intégration, Transcription)};
\end{tikzpicture}
\legende{Diagramme de flux du cycle de réplication du VIH-1 (7 étapes FTITABM). Les étapes 2 (Transcription inverse) et 5-6 se déroulent dans le cytoplasme ; les étapes 3 (Intégration) et 4 (Transcription) nécessitent le noyau. L'étape 7 (Maturation) a lieu au moment ou après le bourgeonnement.}
\end{popfigure}

\begin{experience}[Détail des étapes du cycle viral (Investigation guidée)]
\emph{Observation :} Au microscope électronique, on voit des virions s'attacher à la membrane de lymphocytes, puis des particules immatures bourgeonner, devenir matures et infecter de nouvelles cellules.
\emph{Hypothèse :} Le virus suit un programme ordonné : entrée $\to$ copie ADN $\to$ intégration $\to$ production composants $\to$ assemblage $\to$ sortie $\to$ maturation.
\emph{Données biochimiques :} Inhibiteurs spécifiques bloquent chaque étape (ex: Enfuvirtide blocage fusion, Ténofovir blocage TI, Raltégravir blocage Intégrase, Darunavir blocage Protéase).
\emph{Interprétation :} Chaque étape est validée comme cible thérapeutique.
\emph{Conclusion :} Le cycle se déroule en 7 étapes successives :
\begin{enumerate}
    \item \textbf{Fixation/Fusion} : gp120 se lie au CD4 $\to$ changement conformationnel $\to$ liaison au co-récepteur (CCR5/CXCR4) $\to$ insertion du peptide de fusion gp41 $\to$ fusion des membranes virale et cellulaire $\to$ libération du noyau dans le cytoplasme.
    \item \textbf{Transcription Inverse} (Cytoplasme) : La TI synthétise l'ADN simple brin (brin fort) sur matrice ARN $\to$ activité RNase H dégrade l'ARN matrice $\to$ synthèse du brin faible $\to$ formation de l'ADN double brin linéaire flanqué des LTR (Long Terminal Repeats).
    \item \textbf{Intégration} (Noyau) : Le complexe de pré-intégration (ADN viral + IN + p17 + Vpr) traverse le pore nucléaire $\to$ L'Intégrase effectue le \emph{processing} (coupure 2 nt CA en 3') puis le \emph{strand transfer} (jonction covalente ADN viral-ADN hôte) $\to$ Réparation par l'hôte $\to$ \textbf{ADN proviral} intégré (définitif).
    \item \textbf{Transcription/Traduction} (Noyau/Cytoplasme) : Promoteur LTR 5' $\to$ ARNm (épissés pour Tat, Rev, Nef ; non épissés pour Gag, Pol, Env) $\to$ Export nucléaire (Rev/CRM1) $\to$ Traduction sur ribosomes cytoplasmiques (polyprotéines Pr55Gag, Pr160Gag-Pol, gp160).
    \item \textbf{Assemblage} (Membrane plasmique) : Pr55Gag se multimerise à la face interne de la membrane (domaine MA) $\to$ Recrutement de l'ARN génomique (signal d'empaquetage $\Psi$ via domaine NC) + Pr160Gag-Pol + trimères Env (gp120/gp41) $\to$ Formation de la particule immature.
    \item \textbf{Bourgeonnement} : Échappement via la voie ESCRT (Endosomal Sorting Complex Required for Transport) $\to$ Acquisition de l'enveloppe lipidique cellulaire portant les gp120/gp41.
    \item \textbf{Maturation} (Extracellulaire / surface) : La Protéase virale (PR) s'auto-active par clivage autocatalytique $\to$ Clivage précis et ordonné de Pr55Gag en MA(p17), CA(p24), NC(p7), p6 et de Pr160Gag-Pol en RT, IN, PR $\to$ Réorganisation structurelle (condensation capside conique) $\to$ \textbf{Virion mature, infectieux}.
\end{enumerate}
\end{experience}

\courssub{Variabilité génétique : le moteur de l'échappement}

\begin{exemplebox}[Exemple chiffré : La fabrique à mutants]
\begin{itemize}
    \item Production virale quotidienne : \textbf{$10^9$ à $10^{10}$ virions/jour} chez un patient non traité.
    \item Taux de mutation de la TI : \textbf{$\sim 3 \times 10^{-5}$ mutations par base par cycle de réplication}.
    \item Taille du génome : $\sim 10^4$ bases.
    \item \textbf{Calcul :} Chaque base du génome mute environ $3 \times 10^{-5} \times 10^{10} = 3 \times 10^5$ fois par jour. \textbf{Tous les mutants ponctuels possibles existent chaque jour dans le patient.}
\end{itemize}
$\rightarrow$ Population virale = \cle{quasi-espèces} (nuages de variants génétiquement proches). C'est la base de l'échappement aux anticorps neutralisants, aux réponses CTL (restriction HLA) et à la sélection de résistances aux ARV en monothérapie ou observance insuffisante.
\end{exemplebox}

\begin{popfigure}
\begin{tikzpicture}[scale=0.8, every node/.style={font=\small}, grow'=right, level distance=3.5cm, sibling distance=1.2cm,
    edge from parent/.style={draw=popdark, thick, ->, >=Stealth},
    level 1/.style={sibling distance=3cm},
    level 2/.style={sibling distance=1.5cm}]
    \node[circle, draw=popblue, thick, fill=popblueL, minimum size=0.8cm] {VIH-1}
    child { node[circle, draw=poporange, thick, fill=poporangeL, minimum size=0.7cm] {Groupe M (Majoritaire)}
        child { node[rectangle, draw=popgreen, fill=popgreenL, rounded corners] {Sous-types A, B, C, D... K} }
        child { node[rectangle, draw=popgreen, fill=popgreenL, rounded corners, align=center]{Formes Recombinantes (CRF)\\(ex: CRF02\_AG au Cameroun)} }
    }
    child { node[circle, draw=poppurple, thick, fill=poppurpleL, minimum size=0.7cm] {Groupe N (Non-M, Non-O)} }
    child { node[circle, draw=poppink, thick, fill=poppinkL, minimum size=0.7cm] {Groupe O (Outlier)} }
    child { node[circle, draw=popgold, thick, fill=popgoldL, minimum size=0.7cm] {Groupe P (Pendulaire)} };
    
    \node[align=left, font=\footnotesize, popmuted] at (8, -3) {
        \textbf{Note :} Le Groupe M représente $>99\%$ des infections mondiales.\\
        Au Cameroun : Prédominance du CRF02\_AG ($\sim 50\%$), puis A, G, D, CRF36\_cpx.\\
        Le VIH-2 (non montré) forme des groupes A-H, endémique Afrique de l'Ouest.
    };
\end{tikzpicture}
\legende{Arbre phylogénétique simplifié du VIH-1. Quatre groupes (M, N, O, P) issus de transmissions zoonotiques indépendantes (chimpanzé $\to$ homme pour M, N ; gorille $\to$ homme pour O, P). Le Groupe M (Major) a explosé pandémiquement et se divise en 9 sous-types purs (A-D, F-H, J, K) et de nombreuses Formes Recombinantes Circulantes (CRF). Au Cameroun, le CRF02\_AG est majoritaire.}
\end{popfigure}

\courssub{Conséquences cytopathiques et dépletion progressive des CD4}

L'infection par le VIH entraîne la mort des cellules CD4$^+$ par plusieurs mécanismes non exclusifs :
\begin{enumerate}
    \item \textbf{Lyse virale directe} : Accumulation de protéines virales, bourgeonnement massif $\to$ rupture membrane.
    \item \textbf{Formation de syncytia} : gp120/gp41 exprimés à la surface de la cellule infectée se lient au CD4 de cellules voisines non infectées $\to$ fusion cell-cellule $\to$ cellules géantes multinucléées (syncytia) qui meurent par apoptose/necrose. Fréquent avec souches X4.
    \item \textbf{Apoptose (mort programmée)} : Mécanisme majeur \emph{in vivo}.
        \begin{itemize}
            \item \emph{Voie intrinsèque} : Protéines virales (Vpr, Nef, Tat) perturbent mitochondries $\to$ Cytochrome c $\to$ Caspases.
            \item \emph{Voie extrinsèque} : Fas/FasL, TRAIL.
            \item \emph{Mort des "bystanders" (témoins)} : Cellules CD4$^+$ non infectées meurent par signalisation aberrante (gp120 soluble, activation immune chronique, inflammation).
        \end{itemize}
    \item \textbf{Destruction de l'architecture lymphoïde} : Fibrose des centres germinatifs (dépôt collagène) $\to$ perte du microenvironnement nécessaire à la régénération des cellules T naïves et à la maturation des réponses anticorps $\to$ \textbf{immunodéficience structurelle irréversible} si traitement tardif.
\end{enumerate}

\noindent Résultat : \textbf{Dépletion progressive, continue et irréversible (sans ARV) du pool de lymphocytes T CD4$^+$}, passant de $\sim 800-1200$/mm³ (normal) $< 200$/mm³ (seuil Sida déclaré, risque d'infections opportunistes graves).

\courssub{Le VIH-2 : particularités épidémiologiques et cliniques}

\begin{savaistu}[Le VIH-2 : le "petit frère" ouest-africain]
Découvert en 1986 au Sénégal (équipe Montagnier/Barre-Sinoussi + chercheurs sénégalais). Endémique en Afrique de l'Ouest (Côte d'Ivoire, Guinée-Bissau, Sénégal), rare au Cameroun ($<1\%$ des cas).
\begin{itemize}
    \item \textbf{Moins transmissible} : Charge virale plasmatique plus faible naturellement ; transmission sexuelle et materno-fœtale moins efficace.
    \item \textbf{Mmoins pathogène} : Longue période asymptomatique (souvent $>20$ ans sans traitement) ; progression vers Sida plus lente.
    \item \textbf{Résistance naturelle aux INNRTI} (Névirapine, Éfavirenz, Etravirine) : Poche de liaison de la TI différente. \textbf{Ne jamais utiliser d'INNRTI en monothérapie ou bithérapie pour VIH-2.}
    \item \textbf{Sensible aux IP} (Darunavir, Lopinavir, Atazanavir) et \textbf{II} (Raltégravir, Dolutégravir). Ténofovir + Lamivudine (ou Emtricitabine) + IP/r ou II = schéma de référence.
    \item \textbf{Dépistage} : Tests combinés VIH-1/VIH-2 (ELISA 4ème génération, TDR) détectent les deux. Confirmation par Western Blot différenciant ou PCR spécifique.
\end{itemize}
Au Cameroun, tout test positif VIH-2 ou indéterminé (WB VIH-1 négatif/faible) doit être adressé au Centre Pasteur du Cameroun ou laboratoire de référence pour typage PCR.
\end{savaistu}

\begin{exoresolu}[Interprétation d'un schéma de cycle viral]
\textbf{Énoncé :} Le document ci-contre (schéma simplifié) représente le cycle du VIH dans un lymphocyte T CD4$^+$. 
\begin{enumerate}
    \item Nommez les structures virales A, B, C, D, E.
    \item Identifiez les étapes 1 à 7 (numéros sur flèches).
    \item L'enzyme F permet le passage de l'étape 2 à l'étape 3. Nommez-la et précisez son rôle biochimique exact.
    \item L'enzyme G agit à l'étape 7. Nommez-la. Quel serait l'effet d'un inhibiteur de cette enzyme sur la structure du virion produit ?
\end{enumerate}
\textit{(Schéma mental : Virion entrant $\xrightarrow{1}$ cytoplasme $\xrightarrow{2, F}$ ADN $\xrightarrow{3}$ noyau $\xrightarrow{4}$ ARNm/protéines $\xrightarrow{5}$ assemblage $\xrightarrow{6}$ bourgeonnement $\xrightarrow{7, G}$ virion mature)}

\tcblower
\textbf{Solution détaillée :}
\begin{enumerate}
    \item \textbf{A} = gp120 (glycoprotéine de surface, fixation CD4) ; \textbf{B} = gp41 (glycoprotéine transmembranaire, fusion) ; \textbf{C} = Capside p24 ; \textbf{D} = ARN génomique (2 brins +) ; \textbf{E} = Transcriptase Inverse (TI) / Intégrase / Protéase (enzymes internes).
    \item \textbf{1} = Fixation (gp120-CD4) + Fusion (gp41) ; \textbf{2} = Transcription inverse (ARN $\to$ ADN db) ; \textbf{3} = Intégration (ADN proviral dans chromosome hôte) ; \textbf{4} = Transcription (ADN $\to$ ARNm) + Traduction (ARNm $\to$ Protéines) ; \textbf{5} = Assemblage (Pr55Gag + ARN + Enzymes à la membrane) ; \textbf{6} = Bourgeonnement (acquisition enveloppe) ; \textbf{7} = Maturation (clivage protéolytique).
    \item \textbf{F} = \textbf{Intégrase (IN)}. Rôle : Catalyse la coupure des deux brins d'ADN viral aux extrémités 3' (processing : retrait 2 nucléotides CA) puis le \emph{strand transfer} : attaque nucléophile des hydroxyls 3' viraux sur l'ADN hôte $\to$ jonction covalente ADN viral-ADN hôte (intégration concertée des deux brins).
    \item \textbf{G} = \textbf{Protéase (PR)}. Effet d'un inhibiteur de la Protéase (IP) : Blocage du clivage des polyprotéines précurseurs (Pr55Gag, Pr160Gag-Pol). Les virions libérés sont \textbf{immatures} : capside non condensée (aspect "bétonnière" ou sphérique floue au MET au lieu du cône dense), non infectieux, incapables d'effectuer une nouvelle transcription inverse lors du cycle suivant.
\end{enumerate}
\end{exoresolu}

\begin{aretenir}[Bilan : Le VIH, une machine à évoluer]
\begin{itemize}
    \item \textbf{Structure} : Rétrovirus (Lentivirus), enveloppe gp120/gp41, capside p24, 2 ARN (+), 3 enzymes (TI, IN, PR).
    \item \textbf{Cible} : Lymphocyte T CD4$^+$ (récepteur CD4 + co-récepteur CCR5 ou CXCR4). Réservoirs = cellules CD4$^+$ au repos (ADN proviral latent).
    \item \textbf{Cycle (FTITABM)} : 7 étapes ciblées par 6 classes d'ARV (IE, INNRTI, IP, II, IF, IC).
    \item \textbf{Variabilité} : Erreurs TI (pas de relecture) $\to$ $10^9-10^{10}$ virions/j $\to$ quasi-espèces $\to$ échappement immunitaire + résistance ARV (nécessité trithérapie).
    \item \textbf{Pathogénie} : Mort CD4$^+$ (lyse, syncytia, apoptose bystander) + fibrose organes lymphoïdes $\to$ immunodéficience progressive.
    \item \textbf{VIH-2} : Moins transmissible, moins pathogène, \textbf{Résistant naturel aux INNRTI} (attention prescription !).
\end{itemize}
\end{aretenir}

\coursec{Modes de transmission et prévention}

\noindent Dans la section précédente, nous avons vu que le \cle{VIH} est un rétrovirus fragile hors de son hôte, dont le cycle infectieux cible préférentiellement les lymphocytes T \cle{CD4+}. Cette fragilité extracellulaire et ce tropisme cellulaire expliquent pourquoi le virus ne se transmet pas par n'importe quel contact, mais nécessite un \textbf{véhicule biologique} riche en cellules cibles ou en virions libres, et une \textbf{porte d'entrée} vers le compartiment sanguin ou les muqueuses. C'est ce que nous allons analyser maintenant : quels liquides transmettent le virus ? Par quelles portes d'entrée ? Et comment bloquer efficacement ces passages ?

\courssub{Liquides biologiques : contaminants ou non contaminants}

\begin{experience}[Observation : Pourquoi certains liquides transmettent-ils le VIH et pas d'autres ?]
\textbf{Mise en situation.} Au centre de santé de district, un infirmier reçoit deux patients : l'un s'inquiète d'avoir partagé un verre avec un collègue séropositif ; l'autre a eu un rapport sexuel non protégé. Le premier ne bénéficie pas de PEP, le second oui.
\textbf{Question.} Quelle propriété intrinsèque des liquides biologiques explique cette différence de prise en charge ?
\textbf{Hypothèse.} Seuls les liquides contenant une \textbf{concentration virale suffisante} (charge virale élevée) et/ou des \textbf{cellules cibles infectées} (lymphocytes, macrophages) en quantité critique permettent l'infection.
\textbf{Données expérimentales (études de culture virale et épidémiologiques).}
\begin{itemize}
    \item \textbf{Sang, sperme, sécrétions vaginales, lait maternel, liquides séreux (pleural, péricardique, céphalo-rachidien) :} Culture virale positive fréquente ; forte concentration en ARN VIH et cellules infectées. \textbf{Contaminants.}
    \item \textbf{Salive, sueur, larmes, urine, selles :} Culture virale quasi systématiquement négative ; charge virale très faible ou indétectable ; présence d'inhibiteurs (anticorps, enzymes, pH). \textbf{Non contaminants} (sauf contamination visible par du sang).
\end{itemize}
\textbf{Interprétation.} La barrière muqueuse intacte résiste à de faibles doses virales. L'infection nécessite un \textbf{inoculum minimal} (dose infectante 50\% ou ID$_{50}$) que seuls les liquides « contaminants » apportent naturellement.
\textbf{Conclusion.} La classification des liquides en « contaminants » et « non contaminants » repose sur des données virologiques quantitatives, non sur le simple fait qu'ils soient des fluides corporels.
\end{experience}

\begin{definition}[Liquides biologiques contaminants pour le VIH]
Sont considérés \textbf{contaminants} (risque de transmission avéré) : le \textbf{sang} (et dérivés sanguins), le \textbf{sperme} et le \textbf{liquide pré-éjaculatoire}, les \textbf{sécrétions vaginales} et \textbf{cervicales}, le \textbf{lait maternel}, les \textbf{liquides séreux} (pleural, péricardique, péritonéal, céphalo-rachidien, synovial, amniotique).
Sont \textbf{non contaminants} (risque nul ou négligeable sans sang visible) : la \textbf{salive}, la \textbf{sueur}, les \textbf{larmes}, l'\textbf{urine}, les \textbf{selles}, les \textbf{crachats}, le \textbf{vomi}.
\end{definition}

\begin{popfigure}
\begin{center}
\begin{tabularx}{\linewidth}{|l|X|X|}
\hline
\rowcolor{popblueL}
\textbf{Liquide / Matériel} & \textbf{Niveau de risque} & \textbf{Situations d'exposition typiques} \\
\hline
Sang total, sérum, plasma & \textbf{Très élevé} & Transfusion, partage seringues (UDI), AES (piqûre aiguille), scarifications, accouchement \\
\hline
Sperme / Liquide pré-éjaculatoire & \textbf{Élevé} & Rapports vaginaux, anaux, bucco-génitaux (éjaculation dans la bouche) \\
\hline
Sécrétions vaginales / cervicales & \textbf{Élevé} & Rapports vaginaux (pénétrant), cunnilingus (risque faible mais documenté) \\
\hline
Lait maternel & \textbf{Modéré à Élevé} & Allaitement maternel (risque cumulé durée/exposition) \\
\hline
Liquides séreux (LCR, pleural, etc.) & \textbf{Élevé (théorique)} & Ponctions biomédicales, chirurgie, autopsie (rare en population générale) \\
\hline
Salive, sueur, larmes, urine, selles & \textbf{Nul / Négligeable} & Baiser social, partage couverts/verres, piscine, toilettes, morsure sans sang \\
\hline
Salive \textbf{contaminée par du sang} & \textbf{Faible à Modéré} & Baiser profond (french kiss) avec saignements gingivaux, morsure traumatique \\
\hline
\end{tabularx}
\end{center}
\legende{Classification des liquides biologiques selon leur risque de transmission du VIH. Seuls les liquides de la zone rouge (haut du tableau) ont une efficacité épidémiologique majeure.}
\end{popfigure}

\begin{attention}[Pièges fréquents : Les fausses voies de transmission]
\emph{« On attrape le sida par les moustiques / en serrant la main / aux toilettes / en partageant un repas / en embrassant sur la joue. »} \textbf{FAUX.}
\begin{itemize}
    \item \textbf{Moustique :} Le VIH ne se réplique pas dans le moustique (pas de cycle biologique) ; la proboscis n'injecte pas de sang du repas précédent (mécanique vs biologique) ; dose virale insuffisante.
    \item \textbf{Contacts quotidiens :} Peau intacte = barrière infranchissable. Liquides non contaminants = dose virale $< \text{ID}_{50}$.
    \item \textbf{Baiser social / partage couverts :} Salive non contaminante (inhibiteurs, pH, faible charge virale).
    \item \textbf{Piscine / toilettes :} Dilution massive, virus fragile (sensible au chlore, au séchage, aux UV).
\end{itemize}
\emph{Au Probatoire :} Ne jamais confondre « présence théorique de traces d'ADN viral (PCR positive) » et « infectiosité réelle (culture positive) ».
\end{attention}

\courssub{Les trois modes de transmission : données épidémiologiques camerounaises}

\textbf{2.2.1 Transmission sexuelle (voie majeure : $\approx$ 80\% des cas au Cameroun).}\par
C'est le contact entre muqueuses (pénienne, vaginale, rectale, buccale) et liquides contaminants (sperme, sécrétions vaginales, sang menstruel).

\begin{itemize}
    \item \textbf{Voies basses :}
    \begin{itemize}
        \item \textbf{Vaginal :} Risque par acte $\approx$ 0,04\% (insertif H$\to$F) à 0,08\% (réceptif F$\leftarrow$H). Muqueuse cervicale (zone de transformation) et rectale riches en cellules cibles (\cle{CD4+}, \cle{CCR5+}).
        \item \textbf{Anal (réceptif) :} Risque $\approx$ 1,4\% par acte (le plus élevé). Muqueuse rectale fine, unique couche épithéliale, vascularisation riche, traumatismes fréquents $\to$ porte d'entrée majeure.
    \end{itemize}
    \item \textbf{Voie haute (bucco-génitale) :} Fellation (réceptif) : risque très faible ($\ll$ 0,01\%) mais documenté (éjaculation dans la bouche + gingivites/aphtes). Cunnilingus : risque théorique (sécrétions vaginales + lésions buccales).
    \item \textbf{Co-facteurs majeurs augmentant le risque (x 2 à x 10) :}
    \begin{itemize}
        \item \textbf{IST ulcéreuses} (syphilis, herpès HSV-2, chancroïde) : brèche muqueuse + influx cellules \cle{CD4+}.
        \item \textbf{Charge virale élevée} (infection primaire : CV $> 10^6$ cp/mL ; ou échec thérapeutique) : risque $\times 10$ à $\times 20$.
        \item \textbf{Non-circoncision (homme) :} Prépuce = réservoir de cellules de Langerhans (cibles), humidité, micro-lésions. Circoncision médicale volontaire (CMV) réduit risque H$\leftarrow$F de $\approx$ 60\% (essais ANRS 1265, Kisumu, Rakai, Orange Farm).
        \item \textbf{Rapports secs / violents / multiples partenaires.}
    \end{itemize}
\end{itemize}

\begin{exemplebox}[Risque par acte sexuel non protégé (estimation méta-analyses)]
\begin{center}
\begin{tabular}{|l|c|c|}
\hline
\rowcolor{popblueL}
\textbf{Type d'acte (réceptif $\leftarrow$ insertif)} & \textbf{Risque par acte (moyenne)} & \textbf{Facteur multiplicateur si CV $> 10^5$ cp/mL} \\
\hline
Anal réceptif (H ou F) & $\sim$ 1,38\% (1/72) & $\times 10$ à $\times 20$ \\
\hline
Anal insertif (H) & $\sim$ 0,11\% (1/909) & $\times 10$ \\
\hline
Vaginal réceptif (F) & $\sim$ 0,08\% (1/1250) & $\times 10$ \\
\hline
Vaginal insertif (H) & $\sim$ 0,04\% (1/2500) & $\times 10$ \\
\hline
Bucco-génital (réceptif, éjaculation) & $< 0,01\%$ & Variable \\
\hline
\end{tabular}
\end{center}
\legende{Ordres de grandeur pour un acte unique. Le risque \textbf{cumulé} sur une relation durable non protégée devient très élevé (ex : 1 an $\approx$ 100 actes $\to$ risque $\approx$ 1 - (1-0,0008)$^{100}$ $\approx$ 7,7\% pour vaginal réceptif).}
\end{exemplebox}

\textbf{2.2.2 Transmission sanguine (voie efficace mais minoritaire en population générale).}\par
Le sang est le liquide le plus concentré en virus (cellules infectées + virions libres). Toute effraction cutanée/muqueuse avec du sang contaminant est à risque.
\begin{itemize}
    \item \textbf{Usagers de drogues par injection (UDI) :} Partage seringues/aiguilles/cuillères/filtres. Risque par partage $\approx$ 0,6\% à 2,4\%. \textbf{Très fort risque} (injection directe intravasculaire).
    \item \textbf{Accidents d'exposition au sang (AES) :} Personnel de santé (piqûre aiguille creuse, coupure matériel souillé). Risque moyen $\approx$ 0,3\% (1/300) par piqûre profonde avec sang contaminé. \textbf{Justifie la PEP (voir 2.3.3)}.
    \item \textbf{Transfusion / Produits sanguins :} Risque résiduel $\approx$ 1/500 000 à 1/1 000 000 (dépistage sérologique + PCR pool au Cameroun).
    \item \textbf{Pratiques traditionnelles / esthétiques :} Scarifications, excisions, tatouages, piercings, coiffure (rasoir partagé), circoncision traditionnelle. Risque réel si stérilisation défaillante.
\end{itemize}

\textbf{2.2.3 Transmission materno-fœtale (TMF) ou verticale.}\par
Sans intervention : risque global 15\% à 45\% (selon allaitement).
\begin{itemize}
    \item \textbf{In utero (transplacentaire) :} $\approx$ 5-10\%. Passage virus à travers barrière placentaire (micro-transfusions, infection trophoblaste). Favorisé par charge virale maternelle élevée, paludisme, IST.
    \item \textbf{Per-partum (accouchement) :} $\approx$ 10-20\%. Contact direct sang/sécrétions maternelles + muqueuses/peau lésée du nouveau-né. Rupture membranes prolongée $>$ 4h = facteur de risque.
    \item \textbf{Post-partum (allaitement) :} $\approx$ 5-20\% (risque cumulé sur 18-24 mois). Lait maternel contient virus libre et cellules infectées. Allaitement exclusif 6 mois puis sevrage progressif recommandé si ARV maternels efficaces.
\end{itemize}
\begin{propriete}[Efficacité de la Prise en Charge Mère-Enfant (PMTCT / PECME)]
Avec une \textbf{prise en charge optimale} (ARV maternelle dès 14 SA ou au diagnostic, charge virale indétectable à l'accouchement, prophylaxie nouveau-né adaptée au risque, allaitement exclusif 6 mois sous couverture ARV), le taux de transmission \textbf{descend $< 5\%$ (objectif OMS : $< 2\%$ pour élimination)}.
\textit{Au Probatoire :} Connaître les 4 piliers PMTCT : 1. Prévention infection femme, 2. Prévention grossesses non désirées PVVIH, 3. Blocage TMF (ARV), 4. Prise en charge mère/enfant/famille.
\end{propriete}

\courssub{La prévention combinée : une boîte à outils multidimensionnelle}

Aujourd'hui, aucune méthode unique ne suffit. L'OMS et le Cameroun (Plan Stratégique National) promeuvent la \textbf{prévention combinée} : association synergique de mesures \textbf{biomédicales}, \textbf{comportementales} et \textbf{structurelles}.

\textbf{2.3.1 Barrières mécaniques : Préservatifs.}\par
\begin{itemize}
    \item \textbf{Préservatif masculin (latex/polyisoprène) :} Efficacité $\approx$ 80-95\% en utilisation \textbf{réelle} (vs 98\% théorique). Protège contre VIH, IST, grossesse. \textbf{Seule méthode triple protection.}
    \item \textbf{Préservatif féminin (nitrile) :} Même efficacité, contrôle par la femme, peut être mis en place avant le rapport.
    \item \textbf{Lubrifiant à base d'eau / silicone :} Indispensable pour anal (réduit déchirures préservatif + muqueuse). \textbf{Interdit :} Vaseline, huiles, beurres (détruisent latex en minutes).
\end{itemize}

\textbf{2.3.2 Approches biomédicales révolutionnaires : PrEP, TasP, PEP, CMV.}\par

\begin{definition}[PrEP (Prophylaxie Pré-Exposition) et TasP (Treatment as Prevention)]
\begin{itemize}
    \item \textbf{PrEP :} Prise d'ARV par une \textbf{personne séronégative} à \textbf{haut risque} \textbf{avant} exposition. Schéma standard : \textbf{TDF/3TC} (Ténofovir/Lamivudine) ou \textbf{TDF/FTC} (Ténofovir/Emtricitabine) 1 cp/jour (continue) ou « à la demande » (2-1-1 : 2 cp 2-24h avant, 1 cp 24h après, 1 cp 48h après ; \textbf{validé HSH seulement}). Efficacité $> 90\%$ si observance $\ge$ 4 cp/sem.
    \item \textbf{TasP :} Personne \textbf{séropositive} sous ARV efficace (charge virale \textbf{indétectable} $< 50$ cp/mL depuis $> 6$ mois, observance stricte) \textbf{ne transmet pas le VIH} par voie sexuelle.
\end{itemize}
\end{definition}

\begin{propriete}[Le principe « U = U » : Undetectable = Untransmittable (Indétectable = Intransmissible)]
\textbf{Validé scientifiquement} par les études majeures : \textbf{HPTN 052} (couples sérodifférents, ARV précoce $\to$ réduction transmission 96\%) et \textbf{PARTNER 1 \& 2} (couples gays/hétéros, CV $< 200$ cp/mL, $\approx$ 130 000 actes sexuels non protégés \textbf{zéro} transmission phylogénétiquement liée).
\textbf{Conditions strictes d'application :}
\begin{enumerate}
    \item Observance thérapeutique \textbf{excellente} (pas d'oubli répété).
    \item Charge virale \textbf{indétectable} ($< 50$ cp/mL) confirmée sur \textbf{2 mesures espacées de 3-6 mois}.
    \item Durée de suppression virale \textbf{$> 6$ mois}.
    \item Absence d'IST concomitantes non traitées.
\end{enumerate}
\emph{Message clé :} « Une personne séropositive bien traitée n'est pas un danger, elle est une chance pour la prévention. » Lutte contre la \textbf{sérophobie} interne et externe.
\end{propriete}

\begin{methode}[Conduite à tenir face à une exposition accidentelle (AES ou rapport non protégé à risque) : La PEP]
\textbf{Objectif :} Empêcher l'installation de l'infection avant que le virus ne dissémine (fenêtre $\approx$ 48-72h post-inoculation).
\begin{enumerate}
    \item \textbf{Soins locaux immédiats :} Lavage abondant \textbf{eau + savon} (peau/muqueuse) ; rinçage \textbf{sérum physiologique / eau} (œil/bouche) ; \textbf{pas d'injection d'antiseptique} dans l'œil/urètre/vagin ; \textbf{pas de saignée/protraction} (aggravation lésion).
    \item \textbf{Urgence médicale :} Se rendre aux urgences / centre PEC \textbf{$< 4$h idéal, $< 48$h max, jusqu'à 72h} (au-delà : efficacité non prouvée).
    \item \textbf{Évaluation risque :} Type liquide, porte d'entrée, statut source (connu/inconnu, CV si connu), statut exposé (sérologie J0 obligatoire).
    \item \textbf{Prescription PEP (28 jours) :} Schéma recommandé OMS/Cameroun : \textbf{2 INI + 1 IP} (ex: \textbf{TDF/3TC} + \textbf{LPV/r} ou \textbf{ATV/r} ou \textbf{DTG}) \textbf{OU 3 ARV} (ex: \textbf{TDF/3TC/DTG}). \textit{Préférence actuelle : base DTG (tolérance/simplicité).}
    \item \textbf{Suivi biologique et clinique :} Sérologie VIH \textbf{J0, M1, M3, M6} (recherche Ag p24 + Ac si séroconversion tardive) ; Fonction rénale (créatinine) si TDF ; Hépatique si IP ; Grossesse (test $\beta$-HCG).
    \item \textbf{Conseil :} Préservatif pendant 3 mois (fenêtre sérologique) ; Dépistage IST ; Vaccination VHB si non immunisé ; Soutien psychologique.
\end{enumerate}
\end{methode}

\begin{popfigure}
\begin{center}
\begin{tikzpicture}[
    node distance=1.2cm and 1.5cm,
    box/.style={draw=popblue, thick, rounded corners, fill=popblue!10, minimum width=2.8cm, minimum height=1.2cm, align=center, font=\small},
    arrow/.style={-Stealth, thick, popdark},
    icon/.style={font=\Large, text=popblue}
]
% Nodes
\node[box, align=center] (prep){\textbf{PrEP}\\\textit{Prophylaxie Pré-Exposition}\\\textbf{TDF/3TC} ou \textbf{TDF/FTC}\\\textit{Personne séronégative à risque}};
\node[box, right=of prep, align=center] (tas){\textbf{TasP}\\\textit{Treatment as Prevention}\\\textit{CV Indétectable = Intransmissible}\\\textit{Personne séropositive traitée}};
\node[box, below=of prep, align=center] (pep){\textbf{PEP}\\\textit{Prophylaxie Post-Exposition}\\\textit{$< 72$h, 28 jours ARV}\\\textit{AES, viol, rupture préservatif}};
\node[box, right=of pep, align=center] (condom){\textbf{Préservatif M/F}\\\textit{Barrière mécanique}\\\textit{Seule protection triple}\\\textit{VIH + IST + Grossesse}};
\node[box, below=of pep, align=center] (circ){\textbf{Circoncision Médicale}\\\textit{Réduction risque H$\leftarrow$F $\sim$60\%}\\\textit{Acte chirurgical unique}\\\textit{Complément, pas substitut}};
\node[box, right=of circ, align=center] (dep){\textbf{Dépistage + Traitement IST}\\\textit{Entrée dans la cascade}\\\textit{Couple, Populations clés}\\\textit{Autotest, TROD, Labo}};

% Arrows / Connections (conceptual cascade)
\draw[arrow] (prep) -- node[above, font=\tiny, text=poppurple] {Complémentarité} (tas);
\draw[arrow] (prep) -- node[left, font=\tiny, text=popgreen] {Avant exposition} (pep);
\draw[arrow] (tas) -- node[right, font=\tiny, text=popgreen] {Source contrôlée} (pep);
\draw[arrow] (condom) -- node[below, font=\tiny, text=poporange] {Barrière universelle} (pep);
\draw[arrow] (circ) -- node[below, font=\tiny, text=poporange] {Structurel} (dep);
\draw[arrow, dashed] (dep) -- node[above, font=\tiny, text=poppurple] {Dépistage = Porte d'entrée} (prep);
\draw[arrow, dashed] (dep) -- (tas);

% Central concept
\node[draw=popgold, thick, rounded corners, fill=popgold!15, fit=(prep)(tas)(pep)(condom)(circ)(dep), inner sep=0.5cm, label={[font=\bfseries\color{popdark}, yshift=1.2cm]above:PRÉVENTION COMBINÉE}] {};
\end{tikzpicture}
\end{center}
\legende{Infographie : La cascade de prévention combinée. Chaque outil a sa place (avant/après exposition, personne négative/positive, barrière/médicament). Leur association maximise l'impact populationnel.}
\end{popfigure}

\begin{popfigure}
\begin{center}
\begin{tikzpicture}
\begin{axis}[
    ybar,
    bar width=0.6cm,
    width=\linewidth,
    height=8cm,
    enlarge y limits={upper, value=0.2},
    ylabel={Réduction relative du risque (\%)},
    ymin=0, ymax=110,
    xtick=data,
    xticklabels={Préservatif (usage réel), PrEP (observance optimale), TasP (CV indétectable), Circoncision (\ensuremath{H\leftarrow F}), PEP (initiée $<$ 4h)},
    x tick label style={rotate=30, anchor=east, font=\small},
    ymajorgrids=true,
    grid style=dashed,
    nodes near coords={\pgfmathprintnumber\pgfplotspointmeta\%},
    nodes near coords style={font=\small\bfseries, anchor=south},
    title style={font=\bfseries, text=popdark},
    title={Efficacité comparative des stratégies de prévention biomédicale et mécanique},
    legend style={at={(0.5,-0.25)}, anchor=north, legend columns=-1, font=\small},
    popcurve/.style={draw=popblue, fill=popblue!30},
]
\addplot[popcurve] coordinates {
    (1, 85) 
    (2, 95) 
    (3, 100) 
    (4, 60) 
    (5, 80)
};
\legend{Efficacité estimée (méta-analyses / essais cliniques)}
\end{axis}
\end{tikzpicture}
\end{center}
\legende{Graphique : Efficacité comparative (réduction relative du risque). TasP (U=U) = 100\% (zéro transmission documentée). PrEP > 90\% si observance. Préservatif 80-95\% selon régularité. Circoncision 60\% (protection partielle, homme seulement). PEP $\approx$ 80\% si précoce (< 4h) et observance 28j.}
\end{popfigure}

\begin{savaistu}[La PrEP au Cameroun : accès gratuit pour les populations clés]
Depuis 2017, le Cameroun déploie la PrEP dans le cadre du Plan Stratégique National (PSN 2018-2022, 2024-2028). Elle est \textbf{gratuite} dans les Centres de Prise en Charge (PEC) et certains centres de santé de district pour les \textbf{populations clés} définies par l'OMS et le MINESEC/MSP :
\begin{itemize}
    \item \textbf{HSH} (Hommes ayant des rapports Sexuels avec des Hommes).
    \item \textbf{TS} (Travailleuses du Sexe).
    \item \textbf{UDI} (Usagers de Drogues par Injection).
    \item \textbf{Partenaires séronégatifs de PVVIH} (couples sérodifférents en désir d'enfant ou en attente TasP effective).
    \item \textbf{Personnes transgenres}, détenus, migrants.
\end{itemize}
Le schéma « à la demande » (2-1-1) est autorisé pour les HSH. La PrEP ne dispense \textbf{pas} du préservatif (pas de protection IST/Grossesse). Un suivi rénal (créatinine) et sérologique (VIH, VHB) trimestriel est obligatoire.
\end{savaistu}

\textbf{2.3.3 Approches comportementales et structurelles : au-delà de ABC.}\par
L'approche \textbf{ABC} (\textbf{A}bstinence / retarder le debut, \textbf{B}e faithful / fidélité mutuelle, \textbf{C}ondom) reste la base éducative, mais elle est \textbf{insuffisante seule} (fidélité non vérifiée, abstinence non tenue, préservatif non utilisé systématiquement). Elle est complétée par :
\begin{itemize}
    \item \textbf{Dépistage couple / conjoint :} Entrée commune dans la prévention (PrEP pour le négatif, TasP pour le positif).
    \item \textbf{Traitement précoce des IST :} Réduit co-facteurs ulcéreux/inflammatoires.
    \item \textbf{Réduction du nombre de partenaires / partenariats concomitants :} Diminue la taille du réseau de transmission.
    \item \textbf{Réduction des vulnérabilités structurelles :} Lutte contre violences basées sur genre (VBG), maintien filles à l'école, autonomisation économique femmes, dépénalisation populations clés (réduit stigmatisation = meilleur accès soins).
\end{itemize}

\courssub{Exercice résolu : Analyse d'une situation de prévention}

\begin{exoresolu}[Choix de stratégie préventive pour un couple sérodifférent]
\textbf{Énoncé.} Paul (28 ans) et Grâce (25 ans) sont en couple stable depuis 1 an. Paul est séropositif, sous ARV (\textbf{TDF/3TC/DTG}) depuis 18 mois. Sa dernière charge virale (M-1) est $< 20$ cp/mL. Grâce est séronégative (test J0 négatif). Ils souhaitent avoir un enfant naturellement. Ils n'utilisent plus de préservatif depuis 6 mois (confiance mutuelle). Grâce vient consulter pour savoir si elle doit prendre la PrEP.
\textbf{Question.} En tant que conseiller en santé, que répondez-vous à Grâce ? Justifiez en citant les concepts de TasP, U=U, PrEP, fenêtre de conception.
\textbf{Solution détaillée.}
\begin{enumerate}
    \item \textbf{Analyse de la situation de Paul (Source) :} ARV depuis 18 mois $>$ 6 mois. CV $< 20$ cp/mL (indétectable) confirmée. Observance supposée bonne (CV durable). \textbf{Condition U=U remplie.} Paul \textbf{ne transmet pas} le VIH par voie sexuelle (TasP validé).
    \item \textbf{Risque pour Grâce (Cible) :} Risque sexuel \textbf{nul} (U=U). Risque résiduel théorique uniquement si rebond viral imprévisible (blip) ou arrêt ARV.
    \item \textbf{Indication PrEP pour Grâce ?} \textbf{Non indiquée en routine.} La PrEP s'adresse aux personnes à \textbf{risque substantiel continu}. Ici, le risque est éliminé par la TasP efficace du partenaire. La prescription exposerait Grâce à des effets secondaires (rénaux, osseux avec TDF) et une charge pilulaire inutile.
    \item \textbf{Conduite à tenir recommandée :}
    \begin{itemize}
        \item \textbf{Rassurer le couple} sur le concept U=U (message scientifique fort, lutte contre anxiété).
        \item \textbf{Maintenir l'observance stricte de Paul} (piliers de la TasP). Contrôle CV tous 6 mois.
        \item \textbf{Projet grossesse :} Rapports non protégés ciblés sur \textbf{période fertile} (ovulation) pour maximiser chances conception, minimiser expositions (même si risque nul, hygiène de vie).
        \item \textbf{Dépistage IST couple} avant tentative (co-facteurs).
        \item \textbf{Suivi sérologique Grâce :} Test VIH à 3 mois après conception confirmée (routine grossesse), puis selon protocole PMTCT si grossesse.
    \end{itemize}
    \item \textbf{Nuance :} Si Paul avait une CV détectable ($> 50$ cp/mL) ou observance incertaine $\to$ PrEP pour Grâce \textbf{indiquée} pendant la période de conception + jusqu'à confirmation CV indétectable durable chez Paul.
\end{enumerate}
\tcblower
\textbf{Conclusion.} La prévention combinée n'est pas l'addition de tous les outils pour tous, mais le \textbf{choix rationnel, individualisé et évolutif} de l'outil le plus adapté à la situation (statut sérologique, charge virale, projet de vie, observance). Ici, la \textbf{TasP (U=U) est l'outil unique nécessaire et suffisant}.
\end{exoresolu}

\begin{exercice}[Application : Évaluation de risques et conduite à tenir]
\textbf{Cas 1.} Infirmier, piqûre aiguille creuse après prise de sang sur patient VIH+ (CV inconnue, non traité). Délai : 2h. \textbf{Quelle conduite ?} (Soins locaux, PEP ?, Suivi ?).
\textbf{Cas 2.} Jeune fille 19 ans, rapport vaginal non protégé il y a 36h avec partenaire statut inconnu. Elle demande la « pilule du lendemain VIH ». \textbf{Quelle réponse ?} (PEP ?, Délai ?, Examens J0 ?).
\textbf{Cas 3.} Couple HSH, l'un séropositif (CV 1200 cp/mL, ARV depuis 3 mois), l'autre séronégatif. Ils utilisent préservatif mais rupture hier. \textbf{PEP pour le partenaire négatif ?} (Justifier : CV source détectable = risque réel).
\textbf{Cas 4.} Femme enceinte 28 SA, découverte séropositivité (CV 45 000 cp/mL). Initie ARV (\textbf{TDF/3TC/DTG}) aujourd'hui. \textbf{Risque TMF estimé ?} Mesures urgentes pour l'accouchement et le nouveau-né ? (Prophylaxie nourrisson, CV contrôle M1, mode accouchement ?).
\end{exercice}

\begin{corrige}[Exercice n° Application]
\textbf{Cas 1.} Soins locaux immédiats (eau/savon). \textbf{PEP INDIQUÉE} (source VIH+, porte d'entrée percutanée profonde = risque $\approx$ 0,3\%). Initiation $< 4$h (idéal). Schéma 28j (ex: \textbf{TDF/3TC} + \textbf{DTG}). Sérologie J0, M1, M3, M6. Arrêt PEP si source testé négatif in fine.
\textbf{Cas 2.} \textbf{PEP INDIQUÉE} (délai 36h $<$ 72h). Rapport vaginal réceptif = risque $\approx$ 0,08\% (x10 si source CV haute). Initiation urgence. Sérologie J0 (VIH, VHB, VHC, Syphilis), $\beta$-HCG. Contraception d'urgence si non désirée. Suivi M1, M3, M6.
\textbf{Cas 3.} \textbf{PEP INDIQUÉE.} CV source = 1200 cp/mL (détectable) $\to$ TasP \textbf{non validée} (CV $>$ 50 cp/mL, durée $<$ 6 mois). Rupture préservatif = exposition réelle. Risque anal réceptif $\approx$ 1,4\% (base) $\times$ facteur CV. PEP 28j urgente.
\textbf{Cas 4.} Risque TMF sans intervention $\approx$ 25-35\% (in utero + per-partum + allaitement). \textbf{Urgences :} ARV mère immédiate (cible CV $< 50$ cp/mL à l'accouchement). Contrôle CV à M1 (4 semaines). Si CV $< 50$ cp/mL (ou seuil national, ex $< 1000$ cp/mL) à 36 SA $\to$ voie basse possible. Si CV $> 50$ cp/mL (ou seuil) $\to$ césarienne programmée 38 SA. \textbf{Nouveau-né :} Prophylaxie adaptée au risque maternel à l'accouchement (ex: \textbf{NVP} 6-12 sem si mère VL indétectable ; \textbf{AZT+NVP} 6-12 sem si risque élevé). Allaitement exclusif 6 mois sous couverture ARV mère (CV indétectable) puis sevrage progressif. Sérologie enfant J0, M1, M3, M6, M18 (test virologique \textbf{PCR ADN/ARN} précoce si possible).
\end{corrige}

\coursec{Évolution naturelle de l'infection et marqueurs biologiques}

Dans la section précédente, nous avons vu comment le VIH se transmet et comment on peut s'en protéger. Mais une question demeure : que devient une personne infectée si elle n'est pas traitée ? À Yaoundé comme ailleurs, on entend souvent : « Il a appris sa séropositivité il y a dix ans, il n'a jamais rien eu, et voilà qu'il tombe malade d'un coup. » Cette observation banale cache en réalité une histoire biologique précise, en trois actes, que les médecins savent suivre grâce à deux marqueurs de laboratoire : les \cle{lymphocytes T CD4} et la \cle{charge virale}. C'est cette évolution naturelle de l'infection que nous allons maintenant décortiquer.

\courssub{Une évolution en trois phases : l'observation de départ}

\textbf{Situation-problème.} Dans un centre de santé de Douala, on suit trois patients séropositifs non traités. Le premier, testé trois semaines après une prise de risque, se plaint d'une forte fièvre avec ganglions. Le deuxième, infecté depuis six ans, se porte « comme un charme ». Le troisième, infecté depuis douze ans, maigrit, tousse depuis des mois et présente des infections à répétition. Comment un même virus peut-il produire des tableaux si différents ?

\textbf{Hypothèse.} L'infection n'est pas un état fixe : c'est un processus évolutif, dont l'issue dépend de l'équilibre entre la réplication du virus et les défenses immunitaires.

\textbf{Données.} Le suivi de cohortes de patients non traités (études historiques des années 1980-1990, avant les ARV) a permis de mesurer, mois après mois, deux grandeurs : le nombre de lymphocytes T CD4 par mm$^3$ de sang et la concentration d'ARN viral dans le plasma (la charge virale, en copies/mL). Les courbes obtenues sont remarquablement reproductibles d'un patient à l'autre.

\textbf{Conclusion.} L'infection à VIH non traitée évolue typiquement en trois phases : une \cle{phase aiguë} (ou primo-infection), une \cle{phase chronique asymptomatique} longue, puis une \cle{phase de Sida déclaré}. Étudions-les en détail.

\courssub{Phase 1 : l'infection aiguë (primo-infection)}

\begin{definition}[Infection aiguë ou primo-infection]
La \cle{primo-infection} est la première phase de l'infection à VIH, qui s'étend sur les \textbf{2 à 4 semaines} suivant la contamination. Elle correspond à une réplication virale massive et incontrôlée, avant que le système immunitaire ne parvienne à contenir partiellement le virus.
\end{definition}

Que se passe-t-il concrètement ? Le virus, entré dans l'organisme, infecte préférentiellement les lymphocytes T CD4 et se multiplie de façon explosive. La \cle{charge virale} atteint un \textbf{pic considérable, souvent supérieur à $10^6$ voire $10^7$ copies/mL} de plasma. Simultanément, le nombre de CD4 \textbf{chute brutalement}, parfois sous les valeurs normales.

Dans 50 à 90 \% des cas, cette phase s'accompagne d'un \cle{syndrome pseudo-grippal} (dit aussi mononucléosique) : fièvre, ganglions (adénopathies), éruption cutanée, pharyngite, fatigue, douleurs musculaires. Le piège, c'est que ces signes ressemblent à une grippe ou à un paludisme banal : ils passent inaperçus, et pourtant c'est à ce moment que le patient est \textbf{extrêmement contagieux}, car sa charge virale est maximale.

\begin{definition}[Séroconversion et fenêtre sérologique]
La \cle{séroconversion} est le passage du statut sérologique de négatif à positif : c'est l'apparition, dans le sang, d'\textbf{anticorps anti-VIH} détectables par les tests (ELISA ou tests rapides).\\
La \cle{fenêtre sérologique} est la période comprise entre l'infection et la séroconversion. Sa durée dépend du test utilisé :
\begin{itemize}
\item Pour les tests de \textbf{4ème génération} (recherche simultanée de l'antigène p24 et des anticorps) : \textbf{2 à 3 semaines}.
\item Pour les \textbf{tests rapides} (recherche d'anticorps seulement) : \textbf{4 à 8 semaines}, au maximum 3 mois.
\end{itemize}
Pendant cette fenêtre, les tests de dépistage (qui recherchent les anticorps) sont \textbf{négatifs}, alors que la personne est déjà infectée et contagieuse : seuls l'antigène p24 et l'ARN viral sont détectables.
\end{definition}

\begin{attention}[Le piège de la fenêtre sérologique]
Un test négatif réalisé quelques jours après une prise de risque \textbf{ne prouve rien}. Il faut refaire le test 6 semaines à 3 mois après l'exposition (selon le type de test). C'est une erreur fréquente, au Probatoire comme dans la vie : un test négatif trop précoce est un \emph{faux rassurant}.
\end{attention}

\courssub{Phase 2 : la phase chronique (latence clinique)}

Après la tempête initiale, le système immunitaire (notamment les lymphocytes T CD8 cytotoxiques) reprend partiellement le contrôle. La charge virale redescend et se stabilise à un niveau propre à chaque individu : c'est le \cle{point de fixation} (ou \emph{set point}), atteint environ 6 mois après l'infection.

\begin{definition}[Phase chronique asymptomatique]
La \cle{phase chronique} est une période de \textbf{latence clinique} durant laquelle le patient ne présente aucun symptôme apparent, alors que le virus continue de se répliquer. Sa durée moyenne est de \textbf{8 à 10 ans sans traitement}. Pendant cette phase, un équilibre dynamique s'installe : la production virale quotidienne est à peu près égale à son élimination (clairance) :
\[
\text{Production virale} \approx \text{Clairance virale}
\]
Mais cet équilibre est trompeur : chaque jour, des milliards de virions sont produits et détruits, et les CD4 s'épuisent lentement, à raison d'environ \textbf{50 à 80 cellules/mm$^3$ par an}.
\end{definition}

C'est la phase la plus dangereuse \emph{sur le plan épidémiologique} : la personne se sent en pleine santé, ignore souvent sa séropositivité, et peut transmettre le virus pendant des années. D'où l'importance capitale du dépistage volontaire, que nous étudierons dans la section suivante.

\courssub{Phase 3 : le Sida déclaré}

Quand le stock de CD4 devient insuffisant, l'équilibre se rompt : la charge virale remonte, les défenses s'effondrent. C'est l'entrée dans la phase d'accélération, puis le \cle{Sida} (Syndrome d'ImmunoDéficience Acquise).

\begin{definition}[Sida]
Le \cle{Sida} est le stade terminal de l'infection à VIH. Il est défini biologiquement par un taux de \textbf{CD4 inférieur à 200 cellules/mm$^3$} (ou un pourcentage de CD4 inférieur à 14 \%), et/ou cliniquement par la survenue de \textbf{maladies définissant le Sida} : infections opportunistes graves et cancers.
\end{definition}

L'immunodépression sévère ouvre la porte à des agents infectieux normalement inoffensifs ou bien contrôlés :

\begin{itemize}
\item \textbf{Infections opportunistes} : pneumocystose (pneumonie à \emph{Pneumocystis jirovecii}), toxoplasmose cérébrale, tuberculose (très fréquente au Cameroun, première cause de mortalité des PVVIH en Afrique), candidose œsophagienne, méningite à cryptocoques (cryptococcose) ;
\item \textbf{Maladies néoplasiques (cancers)} : maladie de Kaposi (tumeurs cutanées violacées), lymphomes, cancer du col de l'utérus.
\end{itemize}

\begin{attention}[VIH et Sida : ne pas confondre]
Erreur classique au Probatoire : confondre \textbf{VIH} (le virus, l'infection) et \textbf{Sida} (le stade clinique terminal de cette infection). Être séropositif ne signifie pas « avoir le Sida ». Et surtout : on ne meurt pas \emph{du} VIH, mais \emph{des maladies opportunistes} qui profitent de l'immunodépression. Le VIH est le saboteur ; les infections opportunistes sont les exécutants.
\end{attention}

\begin{propriete}[Pronostic avec et sans traitement]
Sans traitement, la médiane de survie après l'entrée dans le stade Sida est d'environ \textbf{1 à 2 ans}. En revanche, avec des antirétroviraux (ARV) pris précocement et correctement, l'espérance de vie d'une personne vivant avec le VIH devient \textbf{quasi normale}, et la progression vers le Sida est bloquée. C'est toute la différence entre l'histoire naturelle de l'infection et l'histoire modifiée par le traitement.
\end{propriete}

\courssub{Les deux marqueurs biologiques du suivi}

Pour suivre l'évolution de l'infection, le biologiste dispose de deux indicateurs complémentaires, que tu dois savoir interpréter :

\begin{aretenir}[Les deux marqueurs de suivi du VIH]
\begin{itemize}
\item Les \cle{lymphocytes T CD4} (en cellules/mm$^3$ et en \%) : ils mesurent la \textbf{compétence immunitaire}, c'est-à-dire l'état des défenses. Leur seuil conditionne les décisions : mise sous traitement, prophylaxies des infections opportunistes (par exemple cotrimoxazole si CD4 $\leq$ 350/mm$^3$). Valeur normale : environ 500 à 1 500 cellules/mm$^3$.
\item La \cle{charge virale} VIH-1 ARN (en copies/mL, souvent exprimée en $\log_{10}$) : elle mesure l'\textbf{activité de la réplication virale}. C'est le marqueur de référence pour juger l'efficacité des ARV : sous traitement efficace, elle devient indétectable ($<$ 50 copies/mL).
\end{itemize}
\end{aretenir}

\textbf{Relation entre les deux marqueurs.} Charge virale et CD4 sont \textbf{inversement corrélés} : plus la charge virale est élevée (réplication intense), plus le déclin des CD4 est rapide, et plus l'évolution vers le Sida est précoce. Le niveau du point de fixation (set point) est donc un excellent prédicteur de la vitesse de progression : un set point à $10^5$ copies/mL annonce une évolution bien plus rapide qu'un set point à $10^3$ copies/mL.

\begin{popfigure}
\begin{center}
\begin{tikzpicture}[scale=1.0]
% Axes
\draw[->,thick] (0,0) -- (11.2,0) node[below left=2pt and -30pt]{\footnotesize Temps (années)};
\draw[->,thick] (0,0) -- (0,6.6) node[rotate=90,above left=14pt and -55pt]{\footnotesize CD4 (cellules/mm$^3$)};
\draw[->,thick,poporange] (11,0) -- (11,6.6) node[rotate=90,above right=14pt and -60pt,poporange]{\footnotesize $\log_{10}$ Charge virale (cp/mL)};
% Zones des phases
\fill[popblueL,opacity=0.45] (0,0) rectangle (1.4,6.4);
\fill[popgreenL,opacity=0.4] (1.4,0) rectangle (8.6,6.4);
\fill[poppinkL,opacity=0.5] (8.6,0) rectangle (11,6.4);
\node[popdark] at (0.7,6.15) {\footnotesize \textbf{Phase aiguë}};
\node[popdark] at (5,6.15) {\footnotesize \textbf{Phase chronique (8--10 ans)}};
\node[popdark] at (9.8,6.15) {\footnotesize \textbf{Sida}};
% Echelle CD4 gauche
\foreach \y/\lab in {1/200,3/600,5/1000} {\draw (-0.08,\y) -- (0.08,\y); \node[left,font=\footnotesize] at (-0.1,\y){\lab};}
% Echelle CV droite
\foreach \y/\lab in {1/2,3/4,5/6} {\draw[poporange] (10.92,\y) -- (11.08,\y); \node[right,font=\footnotesize,poporange] at (11.1,\y){\lab};}
% Seuil 200 CD4
\draw[dashed,popink,thick] (0,1) -- (11,1) node[pos=0.45,below,popink,font=\footnotesize]{Seuil Sida : CD4 $<$ 200/mm$^3$};
% Courbe CD4 (popblue)
\draw[very thick,popblue] plot[smooth] coordinates {(0,4.6) (0.5,2.6) (1.2,3.6) (2,3.4) (4,2.9) (6,2.3) (8,1.6) (9,0.9) (10.5,0.5)};
\node[popblue,font=\footnotesize] at (6.6,3.1) {\textbf{CD4}};
% Courbe CV (poporange)
\draw[very thick,poporange] plot[smooth] coordinates {(0,0.3) (0.45,5.9) (1.1,3.6) (2,3.2) (4,3.2) (6,3.4) (8,4.0) (9,4.8) (10.5,5.6)};
\node[poporange,font=\footnotesize] at (6.6,4.15) {\textbf{Charge virale}};
% Annotations
\fill[poporange] (0.45,5.9) circle (2.2pt); \node[poporange,font=\footnotesize,align=center] at (1.7,5.7) {Pic CV\\ $>10^6$--$10^7$ cp/mL};
\fill[popink] (1.15,0.6) circle (2pt); \node[popink,font=\footnotesize,align=center] at (2.3,0.55) {Séroconversion\\ (Ac anti-VIH)};
\fill[popdark] (1.9,3.2) circle (2.2pt); \node[popdark,font=\footnotesize,align=center] at (3.6,2.55) {Point de fixation\\ (\emph{set point})};
\end{tikzpicture}
\end{center}
\legende{Évolution naturelle de l'infection à VIH sans traitement. En bleu : déclin progressif des lymphocytes T CD4 (échelle de gauche). En orange : charge virale en échelle logarithmique (échelle de droite). Noter le pic de charge virale et la chute des CD4 pendant la phase aiguë, la stabilisation relative (set point) pendant la phase chronique, puis la rupture d'équilibre conduisant au Sida (CD4 $<$ 200/mm$^3$).}
\end{popfigure}

\courssub{Le bilan immunitaire et la classification de l'infection}

En pratique, le suivi biologique repose sur un \cle{bilan immunitaire} dont voici la logique :

\begin{popfigure}
\begin{center}
\begin{tikzpicture}[
node distance=0.55cm and 1.1cm,
box/.style={draw=popblue,thick,rounded corners,fill=popblueL,align=center,font=\small,inner sep=6pt},
res/.style={draw=poporange,thick,rounded corners,fill=poporangeL,align=center,font=\small,inner sep=6pt},
arr/.style={-{Stealth[length=3mm]},thick,popdark}]
\node[box, align=center] (nfs){Prélèvement sanguin\\ Formule sanguine complète};
\node[box,right=of nfs, align=center] (num){Numération des sous-populations\\ lymphocytaires (immunophénotypage)};
\node[res,below right=0.6cm and 1.1cm of num] (cd4) {T CD4 : ex. 350/mm$^3$};
\node[res,above right=0.6cm and 1.1cm of num] (cd8) {T CD8 : ex. 900/mm$^3$};
\node[box,right=3.2cm of num,fill=popgreenL,draw=popgreen, align=center] (ratio){Ratio CD4/CD8\\ $350/900 \approx 0{,}39 < 1$};
\node[res,below=0.6cm of ratio,fill=poppinkL,draw=popink, align=center] (conc){Inversion du rapport\\ = signature de l'infection à VIH};
\draw[arr] (nfs) -- (num);
\draw[arr] (num.east) -- (cd4.west);
\draw[arr] (num.east) -- (cd8.west);
\draw[arr] (cd4.east) -- (ratio.south west);
\draw[arr] (cd8.east) -- (ratio.north west);
\draw[arr] (ratio) -- (conc);
\end{tikzpicture}
\end{center}
\legende{Démarche du bilan immunitaire. Chez un sujet sain, le rapport CD4/CD8 est supérieur à 1 (environ 1,5 à 2). L'infection à VIH épuise les CD4 et stimule les CD8 : le rapport s'inverse et devient inférieur à 1, signe d'immunodépression.}
\end{popfigure}

Sur le plan clinique, l'Organisation Mondiale de la Santé (OMS) classe l'infection en 4 stades, croisés avec le stade immunologique (taux de CD4) :

\begin{popfigure}
\begin{center}
\begin{tikzpicture}[
grow=right,level distance=3.6cm,
level 1/.style={sibling distance=2.6cm},
level 2/.style={sibling distance=1.25cm},
every node/.style={font=\footnotesize},
edge from parent/.style={draw=popmuted,thick}]
\node[draw=popdark,fill=popcream,rounded corners,thick,align=center]{\textbf{Classification}\\ \textbf{de l'infection VIH}}
child {node[draw=popblue,fill=popblueL,rounded corners,align=center]{Immunologique\\ (CD4/mm$^3$)}
  child {node[draw,fill=white,rounded corners]{$<$ 200 : sévère}}
  child {node[draw,fill=white,rounded corners]{200--349 : avancée}}
  child {node[draw,fill=white,rounded corners]{350--499 : légère}}
  child {node[draw,fill=white,rounded corners]{$\geq$ 500 : aucune}}}
child {node[draw=poporange,fill=poporangeL,rounded corners,align=center]{Clinique OMS\\ (stades 1 à 4)}
  child {node[draw,fill=white,rounded corners,align=center]{Stade 4 : Sida\\ (infections opportunistes, Kaposi)}}
  child {node[draw,fill=white,rounded corners,align=center]{Stade 3 : symptômes sévères\\ (amaigrissement, diarrhée chronique)}}
  child {node[draw,fill=white,rounded corners,align=center]{Stade 2 : symptômes modérés\\ (zona, infections récidivantes)}}
  child {node[draw,fill=white,rounded corners,align=center]{Stade 1 : asymptomatique}}};
\end{tikzpicture}
\end{center}
\legende{Arbre de décision : double classification de l'infection à VIH. La classification clinique de l'OMS (stades 1 à 4) repose sur les signes observés ; la classification immunologique repose sur le taux de CD4. Les deux se croisent pour guider la prise en charge.}
\end{popfigure}

\begin{methode}[Interpréter un bilan VIH]
\begin{enumerate}
\item \textbf{Regarder les CD4} (nombre et pourcentage) : cela donne le \emph{stade immunologique} et le risque d'infections opportunistes.
\item \textbf{Regarder la charge virale} (en copies/mL et en $\log_{10}$) : cela renseigne sur l'\emph{activité de la réplication} virale.
\item \textbf{Croiser les deux} :
\begin{itemize}
\item CD4 bas + CV élevée $\Rightarrow$ urgence : ARV immédiats + prophylaxies des infections opportunistes ;
\item CD4 élevés + CV indétectable sans traitement $\Rightarrow$ profil rare de « contrôleur d'élite ».
\end{itemize}
\end{enumerate}
\end{methode}

\begin{exemplebox}[Un cas concret au centre de santé]
Patient de 32 ans, dépisté VIH positif. Bilan : \textbf{CD4 = 350 cellules/mm$^3$ (18 \%)}, \textbf{charge virale = 50 000 copies/mL} (soit $\log_{10} \approx 4{,}7$). Cliniquement, il présente un zona étendu : \textbf{stade OMS 2}.\\
Décision médicale : démarrage \textbf{immédiat} des antirétroviraux (recommandation OMS \emph{Treat All} : traiter toute personne séropositive dès le diagnostic, quel que soit le taux de CD4) + prophylaxie par \textbf{cotrimoxazole} (indiquée si CD4 $\leq$ 350/mm$^3$). Objectif du suivi : charge virale indétectable en 3 à 6 mois et remontée progressive des CD4.
\end{exemplebox}

\courssub{Facteurs influençant la vitesse d'évolution}

Tous les patients n'évoluent pas à la même vitesse. Plusieurs facteurs modulent la progression :

\begin{itemize}
\item \textbf{L'âge} : les personnes infectées après 40-50 ans progressent généralement plus vite vers le Sida ;
\item \textbf{La génétique de l'hôte} : certains allèles HLA protecteurs (comme \emph{HLA-B*57}) ralentissent l'évolution ; la mutation \emph{CCR5-$\Delta 32$} à l'état homozygote rend les cellules quasi réfractaires à l'infection (le virus ne trouve plus son co-récepteur) ;
\item \textbf{Les co-infections} : hépatites B et C, tuberculose, paludisme (très pertinent au Cameroun : chaque accès palustre stimule transitoirement la réplication virale) ;
\item \textbf{La nutrition, le stress, l'accès aux soins} : la dénutrition accélère l'immunodépression, tandis qu'un suivi médical précoce change tout le pronostic.
\end{itemize}

\begin{savaistu}[Les « contrôleurs d'élite »]
Moins de 0,5 \% des personnes vivant avec le VIH sont des \cle{contrôleurs d'élite} : ils maintiennent une charge virale inférieure à 50 copies/mL \textbf{sans aucun traitement}, grâce à une réponse des lymphocytes T CD8 cytotoxiques particulièrement efficace et/ou à des molécules HLA protectrices. Leur étude passionne les chercheurs, car comprendre leur immunité pourrait ouvrir la voie à un vaccin ou à une guérison fonctionnelle.
\end{savaistu}

\begin{exoresolu}[Interpréter l'évolution d'un patient non traité]
\textbf{Énoncé.} Monsieur N., séropositif depuis 9 ans, n'a jamais pris d'ARV. Ses bilans annuels donnent : année 1 : CD4 = 720/mm$^3$, CV = 12 000 cp/mL ; année 5 : CD4 = 480/mm$^3$, CV = 30 000 cp/mL ; année 9 : CD4 = 160/mm$^3$, CV = 400 000 cp/mL. Il consulte pour une toux chronique et un amaigrissement.\\
1) Calculer la vitesse moyenne de déclin des CD4 entre l'année 1 et l'année 9.\\
2) À quel stade se trouve le patient à l'année 9 ? Justifier.\\
3) Expliquer l'évolution croisée des deux marqueurs.
\tcblower
\textbf{Solution détaillée.}\\
1) Vitesse de déclin :
\[
v = \frac{720 - 160}{9 - 1} = \frac{560}{8} = 70 \text{ cellules/mm}^3\text{/an}
\]
Cette valeur est tout à fait conforme au déclin moyen attendu sans traitement (50 à 80 cellules/mm$^3$/an).\\
2) À l'année 9, les CD4 sont \textbf{inférieurs à 200/mm$^3$} : le patient est au stade de \textbf{Sida déclaré} sur le plan immunologique. La toux chronique et l'amaigrissement évoquent en outre une infection opportuniste (tuberculose à éliminer en priorité au Cameroun), ce qui correspond au stade clinique OMS 3 ou 4.\\
3) Les deux marqueurs sont \textbf{inversement corrélés} : la charge virale passe de 12 000 à 400 000 cp/mL (réplication de plus en plus intense, rupture de l'équilibre du set point), tandis que les CD4 s'effondrent de 720 à 160/mm$^3$. La réplication virale détruit les CD4 ; la perte des CD4 libère la réplication virale : c'est un cercle vicieux que seuls les ARV peuvent briser.
\end{exoresolu}

\begin{aretenir}[L'essentiel de la section]
\begin{itemize}
\item L'infection à VIH non traitée évolue en \textbf{3 phases} : aiguë (pic de CV, chute des CD4, séroconversion), chronique (latence clinique de 8 à 10 ans, déclin lent des CD4), Sida (CD4 $<$ 200/mm$^3$, infections opportunistes et cancers).
\item La \textbf{séroconversion} marque l'apparition des anticorps ; avant elle, la \textbf{fenêtre sérologique} rend les tests négatifs malgré l'infection (durée variable selon le test : 2-3 semaines pour test 4ème génération, jusqu'à 3 mois pour test rapide).
\item Deux marqueurs guident le suivi : les \textbf{CD4} (état immunitaire) et la \textbf{charge virale} (réplication), inversement corrélés.
\item Sans ARV, le Sida est mortel en 1 à 2 ans ; avec ARV précoces, l'espérance de vie est quasi normale. D'où l'importance du dépistage, objet de la section suivante.
\end{itemize}
\end{aretenir}

\coursec{Dépistage : principes, stratégies et algorithmes}

Dans la section précédente, nous avons suivi le destin d'une infection à VIH laissée à elle-même : après la séroconversion, le virus s'installe, les CD4 déclinent lentement et le Sida peut survenir des années plus tard, souvent sans que la personne ne se sente malade. Cette longue phase asymptomatique pose une question de santé publique capitale : \emph{comment savoir que l'on est infecté avant que la maladie ne se déclare ?} La réponse tient en un mot : le \cle{dépistage}. C'est la porte d'entrée de toute la chaîne de prise en charge, et c'est aussi l'un des gestes de prévention les plus puissants, car une personne qui connaît son statut peut se protéger et protéger les autres.

\courssub{Pourquoi dépister ? Une situation concrète}

\textbf{Situation-problème.} Amina, 22 ans, étudiante à Yaoundé, se rend au centre de santé pour une fatigue persistante. Le médecin lui propose un test VIH. Elle hésite : « Je me sens bien, pourquoi faire ce test ? » Son amie, elle, a fait un test trois jours après un rapport non protégé : il était négatif, et elle s'est sentie rassurée. Ont-elles raison toutes les deux ?

Pour répondre, il faut comprendre deux choses : \emph{ce que mesure un test} et \emph{quand il devient fiable}. C'est l'objet de cette section.

\begin{definition}[Dépistage du VIH]
Le \cle{dépistage} du VIH est la recherche, chez une personne ne présentant pas nécessairement de symptômes, de signes biologiques de l'infection : les \cle{anticorps} anti-VIH-1/2 produits par l'organisme, l'\cle{antigène p24} (protéine de la capside virale) ou l'\cle{ARN viral} lui-même. Un dépistage positif n'est jamais un diagnostic définitif : il doit être \textbf{confirmé}.
\end{definition}

\courssub{Les principes techniques des tests}

\textbf{4.2.1. Que cherche-t-on dans le sang ?}

Après l'infection, trois types de marqueurs apparaissent successivement dans le sang :

\begin{itemize}
\item l'\cle{ARN viral} (le génome du virus), détectable le plus tôt, dès la réplication virale initiale ;
\item l'\cle{antigène p24} (Ag p24), protéine interne de la capside, qui apparaît ensuite, avant les anticorps ;
\item les \cle{anticorps} anti-VIH (Ac), d'abord de classe IgM puis de classe IgG, produits par les lymphocytes B : c'est la \cle{séroconversion}.
\end{itemize}

Selon le marqueur recherché, on distingue plusieurs générations de tests :

\begin{center}
\begin{tabularx}{\linewidth}{|l|X|X|}
\hline
\rowcolor{popblueL} \textbf{Type de test} & \textbf{Marqueur recherché} & \textbf{Exemples / usage} \\
\hline
Test de 3\textsuperscript{e} génération & Anticorps anti-VIH-1/2 seuls & Tests rapides (TRD) : Determine, Uni-Gold, SD Bioline \\
\hline
Test de 4\textsuperscript{e} génération (test combiné, Duo/Trio) & Anticorps anti-VIH-1/2 \textbf{+} antigène p24 & Laboratoire ; dépistage plus précoce \\
\hline
Test virologique (PCR) & ARN viral (génome) & Diagnostic précoce du nourrisson, sécurité du don du sang, confirmation de primo-infection \\
\hline
\end{tabularx}
\end{center}

\begin{definition}[Test Rapide de Diagnostic (TRD)]
Un \cle{TRD} est un test \cle{immunochromatographique} réalisé sur une goutte de sang total (piqûre au bout du doigt) ou sur fluide oral, donnant un résultat en 15 à 20 minutes, sans laboratoire. Dans des conditions idéales d'utilisation, sa \cle{sensibilité} et sa \cle{spécificité} dépassent 99~\%. Sa rapidité et sa simplicité permettent le dépistage en dehors des hôpitaux (associations, campagnes communautaires, auto-dépistage).
\end{definition}

\textbf{4.2.2. Tests de confirmation}

Un test de dépistage positif doit être vérifié par un test de \cle{confirmation} : historiquement le \cle{Western Blot} (ou immunoblot), qui identifie séparément les anticorps dirigés contre chaque protéine virale ; dans la pratique actuelle au Cameroun, on utilise un \textbf{deuxième test rapide de principe différent} (stratégie des tests séquentiels, voir 4.4).

\begin{propriete}[Règle d'or de l'interprétation]
Tout test de dépistage positif \textbf{doit} être confirmé par un deuxième test (et un troisième en cas de discordance) avant toute annonce de séropositivité. On retiendra : \emph{« Un seul test ne suffit pas pour dire positif ; un seul test suffit pour dire négatif (hors période de fenêtre) »}. Cette règle protège des erreurs d'annonce, aux conséquences psychologiques et sociales majeures.
\end{propriete}

\courssub{La fenêtre sérologique : le piège du test trop précoce}

\textbf{4.3.1. Observation et démarche}

\textbf{Observation.} Des personnes fraîchement infectées peuvent présenter un test négatif alors qu'elles sont déjà porteuses du virus et contagieuses. \textbf{Hypothèse.} Il existe un délai entre l'entrée du virus et l'apparition de marqueurs détectables. \textbf{Données.} Le suivi de patients dont la date de contamination est connue (accidents d'exposition au sang, par exemple) permet de dater l'apparition de chaque marqueur. \textbf{Interprétation et conclusion.} Chaque test possède une \cle{fenêtre} durant laquelle il peut être faussement négatif ; sa durée dépend du marqueur recherché.

\begin{definition}[Fenêtre sérologique et période d'éclipse]
La \cle{période d'éclipse} est le tout premier intervalle (environ 10 premiers jours) où \emph{aucun} marqueur n'est encore détectable dans le sang. La \cle{fenêtre sérologique} est, pour un test donné, la période entre la contamination et le moment où \emph{ce test} devient capable de détecter l'infection. Pendant cette fenêtre, le test peut être négatif alors que la personne est infectée (faux négatif). Durées indicatives :
\begin{itemize}
\item test ARN (PCR) : environ \textbf{10 jours} (fin de l'éclipse) ;
\item test combiné Ag/Ac de 4\textsuperscript{e} génération : environ \textbf{2 à 3 semaines} ;
\item test rapide à anticorps seuls (3\textsuperscript{e} génération) : environ \textbf{3 à 4 semaines}.
\end{itemize}
Par sécurité, on considère un test 4\textsuperscript{e} génération fiable à \textbf{6 semaines} et la certitude totale acquise à \textbf{3 mois} après la prise de risque.
\end{definition}

\begin{popfigure}
\begin{center}
\begin{tikzpicture}
\begin{axis}[
  width=\linewidth, height=7.5cm,
  xlabel={Temps après l'infection (jours)},
  ylabel={Concentration du marqueur (échelle arbitraire)},
  xmin=0, xmax=90, ymin=0, ymax=10,
  xtick={0,10,20,30,40,50,60,70,80,90},
  ytick=\empty,
  legend style={at={(0.98,0.55)},anchor=east,font=\small,draw=popline,fill=white},
  axis lines=left, thick
]
% Zones : Éclipse (0-10), Fenêtre variable selon test
\addplot[draw=none,fill=popmuted!20] coordinates {(0,0) (10,0) (10,10) (0,10)} \closedcycle;
\node[font=\small\itshape, text=popdark, align=center] at (axis cs:5,9.3) {Période\\ d'éclipse};
% Seuil détection PCR ~10j
\draw[dashed, popink, thick] (axis cs:10,0) -- (axis cs:10,10);
\node[font=\tiny, popink, anchor=south west] at (axis cs:10.5,9.5) {Seuil PCR};
% Seuil détection Ag p24 ~18j
\draw[dashed, poporange, thick] (axis cs:18,0) -- (axis cs:18,10);
\node[font=\tiny, poporange, anchor=south west] at (axis cs:18.5,8.5) {Seuil Ag p24 (4\textsuperscript{e} gén.)};
% Seuil détection Ac (TRD) ~28j
\draw[dashed, popgreen, thick] (axis cs:28,0) -- (axis cs:28,10);
\node[font=\tiny, popgreen, anchor=south west] at (axis cs:28.5,7.5) {Seuil Ac (3\textsuperscript{e} gén.)};
% Courbes
% ARN viral
\addplot[popcurve, color=popink, domain=5:90, samples=120] {9.5/(1+exp(-(x-12)/3))};
\addlegendentry{ARN viral (PCR : $\sim$ 10 j)}
% Ag p24 (pic puis baisse)
\addplot[popcurve, color=poporange, domain=10:90, samples=120] {7/(1+exp(-(x-17)/2.5)) - 7/(1+exp(-(x-35)/4)) + 0.2};
\addlegendentry{Ag p24 (4\textsuperscript{e} gén. : $\sim$ 2--3 sem)}
% IgM
\addplot[popcurve, color=poppurple, domain=15:90, samples=120] {5.5/(1+exp(-(x-24)/3)) - 5/(1+exp(-(x-45)/6))};
\addlegendentry{Ac IgM}
% IgG
\addplot[popcurve, color=popgreen, domain=18:90, samples=120] {8/(1+exp(-(x-30)/4))};
\addlegendentry{Ac IgG (TRD 3\textsuperscript{e} gén. : $\sim$ 3--4 sem)}
\end{axis}
\end{tikzpicture}
\end{center}
\legende{Cinétique d'apparition des marqueurs de l'infection à VIH après la contamination. La période d'éclipse (zone grise) précède toute détection. L'ARN viral apparaît le premier (PCR positive vers 10 jours), puis l'antigène p24 (tests combinés de 4\textsuperscript{e} génération positifs vers 2--3 semaines), puis les anticorps IgM puis IgG (tests rapides de 3\textsuperscript{e} génération positifs vers 3--4 semaines). Les traits verticaux pointillés matérialisent les seuils de détection approximatifs pour chaque type de test.}
\end{popfigure}

\begin{attention}[Le test fait trop tôt : erreur fréquente et dangereuse]
Faire un test \textbf{quelques jours} après une prise de risque (par exemple 3 jours) et conclure à tort « je suis négatif ». Pendant la fenêtre, le test est muet alors que le virus est déjà présent et transmissible (charge virale très élevée). \textbf{Conduite correcte :} respecter le délai (test 4\textsuperscript{e} génération à 6 semaines, contrôle de certitude à 3 mois) et, en cas d'exposition récente à haut risque (rapport non protégé, accident d'exposition au sang), se rendre \textbf{en urgence} (moins de 48--72 h) dans une structure de santé pour une éventuelle \cle{prophylaxie post-exposition (PPE)}.
\end{attention}

\courssub{L'algorithme national de dépistage au Cameroun : la stratégie des 3 tests}

Pour garantir qu'aucun résultat positif ne soit annoncé par erreur, le Programme National de Lutte contre le Sida (PNLS) utilise une \cle{stratégie séquentielle} de trois tests rapides de principes différents (différents antigènes cibles ou différentes technologies) :

\begin{enumerate}
\item \textbf{Test 1} (très sensible, ex. Determine\textsuperscript{TM} HIV-1/2, Alere\textsuperscript{TM}) : s'il est \textbf{négatif}, le résultat rendu est \textbf{négatif} (hors fenêtre sérologique). S'il est \textbf{positif}, on poursuit.
\item \textbf{Test 2} (très spécifique, ex. Uni-Gold\textsuperscript{TM} HIV, Stat-Pak\textsuperscript{TM}) : s'il est \textbf{positif}, la séropositivité est \textbf{confirmée}. S'il est \textbf{négatif}, il y a \textbf{discordance}.
\item \textbf{Test 3} (« \emph{tie-breaker} », ex. SD Bioline\textsuperscript{TM} HIV-1/2 3.0) : il départage. S'il est positif, le résultat est \textbf{positif confirmé} ; s'il est négatif, le résultat est \textbf{négatif} (ou à contrôler après la fenêtre si risque récent).
\end{enumerate}

\begin{popfigure}
\begin{center}
\begin{tikzpicture}[
  font=\small,
  node distance=0.6cm and 1.2cm,
  test/.style={rectangle, rounded corners=3pt, draw=popblue, fill=popblueL, thick, text width=3.5cm, align=center, minimum height=1.1cm},
  res/.style={rectangle, rounded corners=3pt, draw=popgreen, fill=popgreenL, thick, text width=3.2cm, align=center, minimum height=0.9cm},
  disc/.style={rectangle, rounded corners=3pt, draw=poporange, fill=poporangeL, thick, text width=3.2cm, align=center, minimum height=0.9cm},
  fleche/.style={-{Stealth[length=2.5mm]}, thick, popdark}
]
\node[test, align=center] (t1){\textbf{Test 1 (Screening)}\\ Determine / Alere\\(Très sensible)};
\node[res, right=of t1, align=center] (neg1){\textbf{Résultat NÉGATIF}\\ (Hors fenêtre sérologique)};
\node[test, below=of t1, align=center] (t2){\textbf{Test 2 (Confirmation)}\\ Uni-Gold / Stat-Pak\\(Très spécifique)};
\node[res, right=of t2, align=center] (pos1){\textbf{POSITIF CONFIRMÉ}\\ Orientation vers la PEC};
\node[disc, below=of t2, align=center] (d){\textbf{DISCORDANCE}\\ Test 1 (+) / Test 2 (--)};
\node[test, below=of d, align=center] (t3){\textbf{Test 3 (Tie-breaker)}\\ SD Bioline};
\node[res, right=of t3, yshift=0.6cm, align=center] (pos2){\textbf{POSITIF CONFIRMÉ}\\ Orientation vers la PEC};
\node[res, right=of t3, yshift=-0.6cm, align=center] (neg2){\textbf{NÉGATIF}\\ (Contrôle après fenêtre si risque récent)};

\draw[fleche] (t1) -- node[above, font=\footnotesize]{(--)} (neg1);
\draw[fleche] (t1) -- node[left, font=\footnotesize]{(+)} (t2);
\draw[fleche] (t2) -- node[above, font=\footnotesize]{(+)} (pos1);
\draw[fleche] (t2) -- node[left, font=\footnotesize]{(--)} (d);
\draw[fleche] (d) -- (t3);
\draw[fleche] (t3.east) -- node[above, font=\footnotesize, yshift=1pt]{(+)} (pos2.west);
\draw[fleche] (t3.east) -- node[below, font=\footnotesize, yshift=-1pt]{(--)} (neg2.west);
\end{tikzpicture}
\end{center}
\legende{Algorithme national de dépistage du VIH au Cameroun (stratégie des trois tests séquentiels). Le premier test, très sensible, élimine les négatifs ; le deuxième, très spécifique, confirme les positifs ; le troisième tranche les discordances. Aucune séropositivité n'est annoncée sur un seul test.}
\end{popfigure}

\courssub{Les modalités de dépistage : aller vers chacun}

Le dépistage n'est efficace que s'il atteint les personnes concernées. Plusieurs stratégies complémentaires coexistent au Cameroun :

\begin{itemize}
\item le \cle{Dépistage Volontaire} (DV) : la personne vient d'elle-même dans un centre de dépistage (centres associatifs, hôpitaux, unités mobiles) ;
\item le \cle{Dépistage à l'Initiative du Prestataire} (DIP ou PICT) : le soignant propose systématiquement le test à tout consultant, car toute consultation est une occasion de dépistage ;
\item l'\cle{auto-dépistage} (VIH Self-Test) : la personne réalise elle-même le test chez elle, avec un kit sur fluide oral ou sang capillaire ;
\item le \cle{dépistage communautaire} et l'\emph{index testing} (dépistage des partenaires et enfants d'une personne diagnostiquée) ;
\item le \cle{dépistage prénatal} : proposé systématiquement à toute femme enceinte, il est la clé de la prévention de la transmission mère-enfant (PTME/PMTCT).
\end{itemize}

\begin{savaistu}[L'auto-dépistage au Cameroun]
L'auto-dépistage du VIH (VIHST) est légal et promu au Cameroun depuis 2019. Des kits comme OraQuick (fluide oral, sans piqûre) ou Atomo (goutte de sang) sont disponibles en pharmacies et auprès d'ONG. C'est un outil précieux pour les personnes qui craignent le regard des autres ou qui vivent loin des centres de santé. \textbf{Attention} : un auto-test positif doit \emph{toujours} être confirmé dans une structure de santé selon l'algorithme national, et un auto-test négatif reste soumis à la règle de la fenêtre sérologique.
\end{savaistu}

\courssub{Le conseil pré-test et post-test : un test n'est jamais un acte banal}

Le dépistage s'accompagne obligatoirement d'un \cle{conseil} structuré :

\begin{itemize}
\item \textbf{Avant le test (pré-test)} : information claire sur le VIH, sur la signification des résultats et sur la fenêtre sérologique ; \cle{consentement éclairé} (le test n'est jamais forcé, sauf cas très encadrés comme le don du sang) ; rappel de la \cle{confidentialité} absolue.
\item \textbf{Après le test (post-test)} : si le résultat est \textbf{négatif}, conseils de prévention (préservatif, dépistage régulier, PrEP pour les personnes très exposées) et rappel de la fenêtre si risque récent ; s'il est \textbf{positif confirmé}, accompagnement et \cle{orientation immédiate vers la prise en charge} (PEC) pour mise sous antirétroviraux.
\end{itemize}

\begin{methode}[Annoncer un résultat positif : protocole SPIKES adapté]
\begin{enumerate}
\item \textbf{Setting} (cadre) : pièce calme, confidentialité garantie, temps disponible.
\item \textbf{Perception} : évaluer ce que la personne sait et craint (« Que savez-vous du VIH aujourd'hui ? »).
\item \textbf{Invitation} : vérifier qu'elle est prête à recevoir le résultat.
\item \textbf{Knowledge} (annonce) : phrase claire et directe, sans détour : « Le test est positif : vous êtes porteur du VIH. »
\item \textbf{Emotions} : accueillir la réaction (silence, larmes, colère) sans juger ni précipiter.
\item \textbf{Strategy / Summary} (stratégie et résumé) : redonner espoir avec des faits (sous ARV, on vit longtemps et on ne transmet plus le virus quand la charge virale est indétectable = I=I), organiser l'orientation vers la PEC, fixer le prochain rendez-vous.
\end{enumerate}
\end{methode}

\courssub{Indicateurs : la cascade 95-95-95 et l'exemple du Cameroun}

Pour mesurer l'efficacité de la riposte nationale, l'ONUSIDA a fixé pour 2025 la cible \cle{95-95-95} : \textbf{95~\%} des personnes vivant avec le VIH (PVVIH) connaissent leur statut ; \textbf{95~\%} des personnes diagnostiquées sont sous traitement antirétroviral ; \textbf{95~\%} des personnes sous traitement ont une charge virale supprimée (indétectable, donc non transmissible : concept \cle{I=I}, Indétectable = Intransmissible).

\begin{popfigure}
\begin{center}
\begin{tikzpicture}[
  font=\small,
  barre/.style={rectangle, rounded corners=3pt, thick, minimum height=1.35cm, align=center, text=white},
  cible/.style={font=\footnotesize\itshape, text=popdark}
]
% Barre 1 : connaissance du statut ~85%
\fill[popmuted!30, rounded corners=3pt] (0,0) rectangle (9.5,1.35);
\fill[popblue, rounded corners=3pt] (0,0) rectangle (8.1,1.35);
\node[text=white, font=\small\bfseries] at (4.05,0.675) {1\textsuperscript{er} 95 : connaissent leur statut : $\sim$ 85~\%};
\node[cible, anchor=west] at (9.6,0.675) {cible : 95~\%};
% Barre 2 : sous ARV ~90%
\fill[popmuted!30, rounded corners=3pt] (0,-1.8) rectangle (9.5,-0.45);
\fill[poporange, rounded corners=3pt] (0,-1.8) rectangle (8.55,-0.45);
\node[text=white, font=\small\bfseries] at (4.3,-1.125) {2\textsuperscript{e} 95 : sous ARV : $\sim$ 90~\%};
\node[cible, anchor=west] at (9.6,-1.125) {cible : 95~\%};
% Barre 3 : charge virale supprimée ~85%
\fill[popmuted!30, rounded corners=3pt] (0,-3.6) rectangle (9.5,-2.25);
\fill[popgreen, rounded corners=3pt] (0,-3.6) rectangle (8.1,-2.25);
\node[text=white, font=\small\bfseries] at (4.05,-2.925) {3\textsuperscript{e} 95 : charge virale supprimée : $\sim$ 85~\%};
\node[cible, anchor=west] at (9.6,-2.925) {cible : 95~\%};
% Flèches gap
\draw[-{Stealth[length=2.5mm]}, very thick, poppink] (8.1,1.6) -- (9.5,1.6) node[midway, above, font=\footnotesize\bfseries, text=poppink]{gap 1};
\draw[-{Stealth[length=2.5mm]}, very thick, poppink] (8.55,-0.15) -- (9.5,-0.15) node[midway, above, font=\footnotesize\bfseries, text=poppink]{gap 2};
\draw[-{Stealth[length=2.5mm]}, very thick, poppink] (8.1,-1.95) -- (9.5,-1.95) node[midway, above, font=\footnotesize\bfseries, text=poppink]{gap 3};
\end{tikzpicture}
\end{center}
\legende{Infographie de la cascade de dépistage et de prise en charge au Cameroun (estimations récentes, ordres de grandeur 85--90--85) comparée aux cibles ONUSIDA 95-95-95. Le premier « gap » (personnes qui ignorent leur séropositivité) est le plus préoccupant, car il constitue le réservoir silencieux de la transmission.}
\end{popfigure}

\begin{exemplebox}[Le poids du premier « gap » au Cameroun]
On estime qu'environ \num{500000} personnes vivaient avec le VIH au Cameroun en 2023. Si le taux de connaissance du statut est d'environ 85~\%, alors environ $500\,000 \times (1 - 0,85) = 75\,000$ personnes \textbf{ignorent leur séropositivité}. Non traitées, elles peuvent transmettre le virus sans le savoir : c'est le réservoir silencieux de l'épidémie. D'où l'importance capitale des stratégies qui vont « chercher » les personnes : DIP, auto-dépistage, dépistage communautaire et index testing.
\end{exemplebox}

\begin{exoresolu}[Interpréter un parcours de dépistage]
\textbf{Énoncé.} Junior a eu un rapport non protégé il y a 10 jours. Inquiet, il fait un test rapide (3\textsuperscript{e} génération) : négatif. Trois mois plus tard, dans le cadre d'une campagne de dépistage, le Test 1 (Determine) est positif, le Test 2 (Uni-Gold) est négatif, puis le Test 3 (SD Bioline) est positif. 1) Le premier test négatif prouvait-il l'absence d'infection ? 2) Quel est le statut final de Junior ? 3) Quelle conduite doit suivre l'équipe ?
\tcblower
\textbf{Solution détaillée.}
\begin{enumerate}
\item \textbf{Non.} À 10 jours, un test rapide de 3\textsuperscript{e} génération (anticorps seuls, fenêtre de 3 à 4 semaines) est encore dans sa fenêtre sérologique : les anticorps ne sont pas encore détectables. Seule une PCR (ARN viral) aurait pu être positive à ce stade. Le résultat négatif ne permettait donc aucune conclusion ; il fallait refaire un test à 6 semaines (test 4\textsuperscript{e} génération) puis à 3 mois.
\item Selon l'algorithme national : Test 1 positif $\rightarrow$ Test 2 négatif = \textbf{discordance} $\rightarrow$ le Test 3 (tie-breaker) tranche : il est positif, donc le résultat est \textbf{VIH positif confirmé}.
\item L'équipe doit annoncer le résultat selon un protocole adapté (cadre confidentiel, annonce claire, accueil des émotions), puis \textbf{orienter immédiatement} Junior vers une unité de prise en charge pour mise sous antirétroviraux, avec conseils de prévention pour ses partenaires (index testing possible) et soutien psychosocial. On lui expliquera qu'avec un traitement bien suivi, sa charge virale peut devenir indétectable : il vivra longtemps et ne transmettra plus le virus (I=I).
\end{enumerate}
\end{exoresolu}

\begin{aretenir}[L'essentiel de la section]
\begin{itemize}
\item Les tests recherchent l'ARN viral (PCR, $\sim$ 10 jours), l'antigène p24 + anticorps (4\textsuperscript{e} génération, $\sim$ 2--3 semaines) ou les anticorps seuls (TRD 3\textsuperscript{e} génération, $\sim$ 3--4 semaines).
\item La \cle{fenêtre sérologique} interdit toute conclusion sur un test fait trop tôt : fiabilité à 6 semaines (test 4\textsuperscript{e} gén.), certitude à 3 mois.
\item Algorithme national : 3 tests séquentiels ; \emph{un test ne suffit pas pour dire positif, un test suffit pour dire négatif (hors fenêtre)}.
\item Le dépistage se décline en DV, DIP, auto-dépistage, dépistage communautaire et prénatal, toujours encadré par le conseil, le consentement et la confidentialité.
\item Objectif ONUSIDA 2025 : la cascade 95-95-95 (et I=I) ; au Cameroun, le premier gap (personnes ignorant leur statut) reste le défi majeur.
\end{itemize}
\end{aretenir}

\coursec{Prise en charge : Traitements Antirétroviraux (ARV) et suivi}

Après le diagnostic de séropositivité (section 4), l'annonce du résultat au patient marque le début d'une nouvelle étape : la \textbf{prise en charge globale}. Au Cameroun, comme dans la majorité des pays à revenu faible ou intermédiaire, cette prise en charge suit les recommandations consolidées de l'OMS (2021-2024) et les directives nationales du Programme National de Lutte contre le Sida (PNLS). Elle repose sur trois piliers indissociables : l'initiation rapide d'un traitement antirétroviral (ARV) efficace, un suivi biologique et clinique régulier, et un accompagnement psychosocial pour garantir l'observance à vie. L'objectif ultime n'est plus seulement de prolonger la vie, mais d'atteindre une \textbf{suppression virale durable} permettant au patient de vivre en bonne santé et de ne pas transmettre le virus (\textbf{TasP} : Treatment as Prevention).

\courssub{Objectifs de la thérapie antirétrovirale (TARV)}

\begin{definition}[Objectifs thérapeutiques de la TARV selon l'OMS et le PNLS Cameroun]
La thérapie antirétrovirale vise quatre objectifs cliniques et virologiques majeurs, évalués lors du suivi :
\begin{enumerate}
    \item \textbf{Suppression virale durable :} Obtenir et maintenir une charge virale (CV) plasmatique \textbf{inférieure à \SI{50}{copies/mL}} (seuil de détection des techniques standards) de façon persistante. C'est le marqueur principal de l'efficacité du traitement.
    \item \textbf{Restauration immunitaire :} Obtenir une élévation significative et durable du taux de lymphocytes T CD4+ (idéalement $> \SI{500}{\per\micro\liter}$) pour prévenir les infections opportunistes et les cancers liés au VIH.
    \item \textbf{Réduction de la morbidité et de la mortalité :} Faire régresser les signes cliniques (stade OMS), prévenir les infections opportunistes (TB, cryptococcose, etc.) et réduire la mortalité liée au Sida.
    \item \textbf{Prévention de la transmission (TasP) :} Un patient sous ARV efficace, avec CV indétectable depuis au moins 6 mois et bonne observance, \textbf{ne transmet pas le VIH} par voie sexuelle (\cle{I = I} : Indétectable = Intransmissible). Cela inclut la prévention de la transmission mère-enfant (PTME).
    \item \textbf{Amélioration de la qualité de vie :} Normalisation de la vie sociale, professionnelle et reproductive.
\end{enumerate}
\end{definition}

\begin{attention}[Piège conceptuel : "Guérison" vs "Contrôle"]
Les ARV \textbf{ne guérissent pas} l'infection par le VIH. Ils bloquent la réplication virale mais n'éliminent pas le \textbf{réservoir viral} (lymphocytes CD4+ mémoire à longue durée de vie, cellules du système nerveux central, macrophages tissulaires). L'arrêt du traitement entraîne inévitablement une \textbf{reprise virale (rebond)} en 2 à 4 semaines, avec risque de sélection de mutants résistants. La TARV est donc un engagement \textbf{à vie}, sauf dans le cadre de protocoles de recherche sur la rémission (voir \textit{Le savais-tu ?}).
\end{attention}

\courssub{Cibles thérapeutiques : Les classes d'ARV et leur mécanisme d'action}

Le cycle de réplication du VIH (vu en section 1) offre plusieurs cibles enzymatiques ou structurelles. Une trithérapie efficace associe au moins \textbf{trois molécules} issues d'au moins \textbf{deux classes différentes} pour créer une \textbf{barrière génétique} élevée (difficulté pour le virus de muter simultanément sur plusieurs cibles).

\begin{experience}[Démarche d'investigation : Identifier la cible d'un ARV]
\textbf{Observation :} En culture cellulaire, on ajoute différents ARV à des cellules CD4+ infectées par le VIH. On dose ensuite l'ADN viral total, l'ADN viral intégré (provirus) et l'ARN viral dans le surnageant.
\textbf{Résultats :}
\begin{itemize}
    \item Molécule A : Blocage de la formation de l'ADN viral total. L'ARN viral n'est pas produit.
    \item Molécule B : ADN viral total formé, mais \textbf{pas d'ADN intégré}. L'ARN viral n'est pas produit.
    \item Molécule C : ADN viral intégré présent, ARN viral produit, mais \textbf{virions libérés immatures, non infectieux}.
    \item Molécule D : Blocage de la fixation du virus sur la cellule (pas d'entrée).
\end{itemize}
\textbf{Interprétation :}
\begin{itemize}
    \item A cible la \textbf{Transcriptase Inverse (TI)} (INTI ou INNTI).
    \item B cible l'\textbf{Intégrase} (IRTI).
    \item C cible la \textbf{Protéase} (IP).
    \item D cible l'\textbf{Entrée} (IEF : CCR5 ou gp120).
\end{itemize}
\textbf{Conclusion :} Chaque classe d'ARV bloque une étape précise et irréversible du cycle viral. L'association de plusieurs classes empêche le virus de contourner le blocage par mutation simple.
\end{experience}

\begin{popfigure}
\begin{tikzpicture}[
    node distance=1.2cm and 1.5cm,
    box/.style={rectangle, draw=popdark, thick, fill=popcream, rounded corners, minimum height=1.2cm, minimum width=2.8cm, align=center, font=\small},
    arrow/.style={-Stealth, thick, popblue},
    inhib/.style={-Stealth, thick, poporange, dashed},
    target/.style={ellipse, draw=popgreen, thick, fill=popgreenL, minimum height=0.8cm, minimum width=2.2cm, align=center, font=\scriptsize}
]
% Cycle viral steps
\node[box, align=center] (entree){1. Entrée\\Fixation gp120/CD4/CCR5\\Fusion};
\node[box, right=of entree, align=center] (ti){2. Transcription Inverse\\ARN viral $\to$ ADN double brin\\Enzyme : Transcriptase Inverse};
\node[box, right=of ti, align=center] (integr){3. Intégration\\ADN viral $\to$ ADN hôte\\Enzyme : Intégrase};
\node[box, right=of integr, align=center] (transc){4. Transcription/Traduction\\ARNm $\to$ Protéines virales};
\node[box, right=of transc, align=center] (assemb){5. Assemblage/Maturation\\Polyprotéines Gag-Pol\\Enzyme : Protéase};

% Arrows between steps
\draw[arrow] (entree) -- (ti);
\draw[arrow] (ti) -- (integr);
\draw[arrow] (integr) -- (transc);
\draw[arrow] (transc) -- (assemb);

% Inhibitors (below)
\node[target, below=0.8cm of entree] (cib1) {IEF : MVC, Fostemsavir};
\node[target, below=0.8cm of ti, align=center] (cib2){INTI : TDF, 3TC, FTC, ABC, AZT\\INNTI : EFV, NVP, RPV, DOR};
\node[target, below=0.8cm of integr] (cib3) {IRTI : DTG, RAL, BIC};
\node[target, below=0.8cm of assemb] (cib4) {IP : LPV/r, DRV/r, ATV/r};

% Inhibition arrows
\draw[inhib] (cib1) -- (entree) node[midway, right, font=\tiny, poporange] {Bloque fixation/fusion};
\draw[inhib] (cib2) -- (ti) node[midway, right, font=\tiny, poporange] {Terminaison chaîne / Allostérique};
\draw[inhib] (cib3) -- (integr) node[midway, right, font=\tiny, poporange] {Bloque intégration provirus};
\draw[inhib] (cib4) -- (assemb) node[midway, right, font=\tiny, poporange] {Bloque clivage maturation};
\end{tikzpicture}
\legende{Cibles des différentes classes d'antirétroviraux sur le cycle de réplication du VIH. Les flèches pleines bleues montrent le cycle viral ; les flèches pointillées oranges montrent le point de blocage de chaque classe thérapeutique. INTI = Inhibiteurs Nucléosidiques de la Transcriptase Inverse ; INNTI = Inhibiteurs Non Nucléosidiques de la Transcriptase Inverse ; IRTI = Inhibiteurs de l'Intégrase ; IP = Inhibiteurs de la Protéase ; IEF = Inhibiteurs d'Entrée/Fusion.}
\end{popfigure}

\begin{propriete}[Classes thérapeutiques ARV disponibles au Cameroun (Guide national 2024)]
\begin{center}
\begin{tabularx}{\linewidth}{|l|X|X|X|}\hline
\rowcolor{popblueL}
\textbf{Classe (Abrév.)} & \textbf{Cible / Mécanisme} & \textbf{Molécules principales (Cameroun)} & \textbf{Particularités / Toxicités majeures} \\ \hline
\textbf{INTI} (NRTI) & \textbf{Transcriptase Inverse} : Analogues nucléosidiques $\to$ terminaison de chaîne ADN (compétition substrat). & \textbf{TDF} (Ténofovir), \textbf{3TC} (Lamivudine), \textbf{FTC} (Emtricitabine), \textbf{ABC} (Abacavir), \textbf{AZT} (Zidovudine). & TDF : Rénal (tubulopathie), Os (ostéoporose). AZT : Anémie, Neutropénie (moelle). ABC : Risque hypersensibilité (test HLA-B*5701 idéal). \\ \hline
\textbf{INNTI} (NNRTI) & \textbf{Transcriptase Inverse} : Fixation site allostérique $\to$ changement conformationnel $\to$ inhibition non compétitive. & \textbf{EFV} (Éfavirenz 400/600 mg), \textbf{NVP} (Névirapine), \textbf{RPV} (Rilpivirine), \textbf{DOR} (Doravirine). & EFV : Neuropsychiatrique (rêves, dépression), tératogène (1er trimestre). NVP : Hépatotoxicité, rash (lead-in 14j). RPV : Nécessite acidité gastrique, CV $< \SI{100000}{copies/mL}$. \\ \hline
\textbf{IRTI} (INSTI) & \textbf{Intégrase} : Bloque le transfert de brin (Strand Transfer) $\to$ empêche intégration provirus. & \textbf{DTG} (Dolutégravir), \textbf{RAL} (Raltégravir), \textbf{BIC} (Bictegravir). & \textbf{DTG : Haute barrière génétique, pilule petite, 1 prise/j, peu d'interactions. \cle{Standard 1ère ligne OMS}.} RAL : 2 prises/j, barre plus basse. \\ \hline
\textbf{IP} (PI) & \textbf{Protéase} : Bloque clivage polyprotéines Gag-Pol $\to$ virions immatures, non infectieux. & \textbf{LPV/r} (Lopinavir/ritonavir), \textbf{DRV/r} (Darunavir/ritonavir), \textbf{ATV/r} (Atazanavir/ritonavir). & \textbf{Boostés par Ritonavir (r) ou Cobicistat (c)}. Digestifs (diarrhée), Dyslipidémie, Diabète, Interactions médicamenteuses majeures (CYP3A4). \\ \hline
\textbf{IEF} & \textbf{Entrée} : MVC (Maraviroc) antagoniste CCR5. Fostemsavir (inhibiteur d'attachement gp120). & \textbf{MVC} (rare, 3ème ligne), \textbf{Fostemsavir} (nouveau, 3ème ligne). & MVC : Nécessite test de tropisme viral (R5 uniquement). Coût élevé. \\ \hline
\end{tabularx}
\end{center}
\end{propriete}

\courssub{Lignes de traitement au Cameroun (Guidelines 2024) : Schéma thérapeutique standardisé}

Le Cameroun a adopté l'approche \textbf{"Treat All"} (Traitement pour tous) depuis 2016 : tout patient VIH+, quel que soit son stade clinique ou son taux de CD4, doit être mis sous ARV le plus tôt possible (idéalement J+7 après confirmation, voire J+0 si stade avancé).

\begin{popfigure}
\begin{tikzpicture}[
    node distance=0.8cm and 1cm,
    line/.style={draw=popdark, thick, rounded corners, fill=popcream, minimum height=1.1cm, minimum width=3.5cm, align=center, font=\small},
    linealt/.style={line, fill=popblueL},
    linealtb/.style={line, fill=poporangeL},
    linealtc/.style={line, fill=popgreenL},
    arrow/.style={-Stealth, thick, popdark},
    brace/.style={decorate, decoration={brace, amplitude=5pt}, thick, popdark}
]
% 1ere ligne
\node[linealt, align=center] (l1a){\textbf{1ère LIGNE STANDARD (Adulte/Adolescent $>$\SI{30}{kg})}\\\textbf{TLD : TDF + 3TC + DTG} (1 cp/j)};
\node[line, below=of l1a] (l1b) {Alternative 1 : \textbf{TLE400} (TDF+3TC+EFV 400mg)};
\node[line, below=of l1b] (l1c) {Alternative 2 (Insuff. Rénale) : \textbf{ABC + 3TC + DTG}};
\node[line, below=of l1c] (l1d) {Alternative 3 (Femme enceinte 1er trim. si DTG refusé/contre-indiqué*) : \textbf{TLE400}};

% 2eme ligne
\node[linealtb, right=3cm of l1a, yshift=-2cm] (l2a) {\textbf{2ème LIGNE (Échec virologique 1ère ligne)}};
\node[line, below=of l2a, align=center] (l2b){Si 1ère ligne \textbf{INNTI/INTI} (ex: TLE) :\\ \textbf{AZT + 3TC + LPV/r} (ou DRV/r)};
\node[line, below=of l2b, align=center] (l2c){Si 1ère ligne \textbf{DTG} (ex: TLD) :\\ \textbf{DTG + 2 INTI neufs} (ex: AZT+3TC ou TDF+3TC si CV $<$\SI{1000}{copies/mL}) \\ \textit{+ Génotypage obligatoire}};

% 3eme ligne
\node[linealtc, right=3cm of l2a, yshift=-1.5cm] (l3a) {\textbf{3ème LIGNE (Échec 2ème ligne)}};
\node[line, below=of l3a] (l3b) {Comité d'expertise national (CEN)};
\node[line, below=of l3b, align=center] (l3c){Nouvelles molécules :\\ \textbf{Lenacapavir} (Capside, 1 inj/6 mois)\\ \textbf{Fostemsavir} (Attachement gp120)\\ \textbf{DTG} (si jamais utilisé)\\ \textbf{DRV/r} optimisé};

% Pediatrie / Grossesse
\node[line, below=1.5cm of l1d] (ped) {\textbf{PÉDIATRIE ($<$\SI{30}{kg} / $<$\SI{10}{ans})}};
\node[line, below=of ped] (ped1) {1ère ligne : \textbf{ABC + 3TC + DTG} (dispersible 10/5/5 mg)};
\node[line, below=of ped1] (ped2) {Nourrisson $<$\SI{4}{sem} / $<$\SI{3}{kg} : AZT+3TC+LPV/r (sirop)};

\node[line, below=1.5cm of ped2] (gross) {\textbf{FEMME ENCEINTE / ALLAITEMENT}};
\node[line, below=of gross, align=center] (gross1){DTG \textbf{recommandé} (OMS 2024) \textbf{pour toutes} (y.c. 1er trim.)\\ \textit{avec contraception efficace si désir de grossesse non planifié}};
\node[line, below=of gross1] (gross2) {Si DTG contre-indiqué/refusé : TLE400 (EFV 400mg)};

% Arrows
\draw[arrow] (l1a.east) -- ++(1,0) |- (l2a.west) node[midway, above, font=\tiny, poporange] {Échec virologique CV $>$\SI{1000}{copies/mL}};
\draw[arrow] (l2a.east) -- ++(1,0) |- (l3a.west) node[midway, above, font=\tiny, poporange] {Échec 2ème ligne};
\end{tikzpicture}
\legende{Schéma thérapeutique standardisé au Cameroun (Guidelines PNLS/OMS 2024). TLD = Ténofovir/Lamivudine/Dolutégravir (pilule unique, 1 prise/jour). *Note : Les données de surveillance (Botswana, Brésil) ont levé la contre-indication formelle du DTG au 1er trimestre (risque défaut tube neural non confirmé significatif), mais la contraception efficace reste recommandée chez la femme en âge de procréer sous DTG.}
\end{popfigure}

\begin{methode}[Choix de la 1ère ligne : Pourquoi TLD (TDF+3TC+DTG) est le standard ?]
\begin{enumerate}
    \item \textbf{Efficacité virologique supérieure :} DTG (Intégrase 2ème gén.) a une \textbf{haute barrière génétique} (mutations multiples nécessaires pour résistance) vs EFV/NVP (barrière basse, 1 mutation = résistance).
    \item \textbf{Tolérance / Observance :} 1 comprimé / 1 prise / jour, petite taille, peu d'effets neuropsychiatriques (vs EFV), pas de lead-in (vs NVP).
    \item \textbf{Robustesse :} Efficace même si CV initiale élevée ($> \SI{100000}{copies/mL}$) ou CD4 bas (vs RPV).
    \item \textbf{Coût :} Génériques préqualifiés OMS $\approx$ \SI{60-70}{USD/an/patient} (vs \SI{200-300}{USD} pour 2ème ligne IP).
    \item \textbf{Population large :} Adulte, ado, femme enceinte, co-infection TB (ajustement dose DTG \SI{50}{mg} 2x/j avec Rifampicine).
\end{enumerate}
\end{methode}

\courssub{Suivi biologique et clinique : Détecter l'échec, prévenir la résistance}

Le suivi ne vise pas seulement à "voir si ça marche", mais à \textbf{détecter précocement l'échec virologique} pour switcher avant l'accumulation de mutations de résistance complexes.

\begin{methode}[Calendrier de suivi biologique standard sous ARV (Adulte stable, Cameroun)]
\begin{enumerate}
    \item \textbf{Mois 0 (Bilan pré-thérapeutique - BPT) :} Numération-Formule-Sanguine (NFS), \textbf{CD4}, \textbf{Charge Virale (CV) de base}, Créatinine (DFG si TDF), Glycémie, Hépatites B/C (AgHBs, AcAnti-HCV), Grossesse (femme), Recherche BK (TB).
    \item \textbf{Mois 3 (M3) :} CV (objectif : chute $> 1$ Log ou $< \SI{1000}{copies/mL}$), Créatinine (si TDF).
    \item \textbf{Mois 6 (M6) :} \textbf{CV (critère principal suppression : $< \SI{50}{copies/mL}$)}, CD4, Créatinine (si TDF).
    \item \textbf{Mois 12 (M12) :} CV, CD4, Lipidogramme (si IP), Glycémie, Hémoglobine (si AZT).
    \item \textbf{Ensuite (Patient stable, CV $< \SI{50}{copies/mL}$) :} \textbf{CV + CD4 tous les 12 mois}. NFS/Lipides/Glycémie selon tolérance.
    \item \textbf{En cas d'échec suspecté (CV $> \SI{1000}{copies/mL}$) :} \textbf{Contrôle CV à 3 mois (M+3)} après renforcement observance. Si CV persistante $> \SI{1000}{copies/mL}$ $\to$ \textbf{Génotypage (Séquençage région \textit{pol})} pour guider la 2ème ligne.
\end{enumerate}
\end{methode}

\begin{popfigure}
\begin{tikzpicture}
\begin{axis}[
    width=\linewidth, height=8cm,
    xlabel={Temps sous ARV efficaces (semaines / mois)},
    ylabel={Log$_{10}$ Charge Virale (cp/mL) / CD4 (cellules/$\mu$L)},
    grid=major, grid style={popline},
    legend style={at={(0.98,0.98)}, anchor=north east, fill=popcream, draw=popdark, font=\small},
    xmin=0, xmax=52,
    ymin=1, ymax=6.5,
    xtick={0,2,4,8,12,24,36,52},
    xticklabels={0, 2 sem, 1 mois, 2 mois, 3 mois, 6 mois, 9 mois, 12 mois},
    ytick={1,2,3,4,5,6},
    yticklabels={$10^1$, $10^2$, $10^3$, $10^4$, $10^5$, $10^6$},
    scaled y ticks=false,
    yticklabel style={/pgf/number format/fixed, /pgf/number format/precision=0},
]
% Charge Virale Biphasique
\addplot[popcurve, popblue, very thick, domain=0:52, samples=200] 
    ({x}, 
     {5.5 - 1.5*tanh((x-1)/1.5) - 1.0*tanh((x-10)/15)} 
    ) 
    node[pos=0.15, above left, font=\tiny, popblue] {Phase 1 : Chute rapide (1-2 Log/2 sem)}
    node[pos=0.7, below right, font=\tiny, popblue] {Phase 2 : Chute lente $\to$ $<50$};
    
% Seuil indétectable
\addplot[poporange, dashed, thick] coordinates {(0,1.7) (52,1.7)} node[pos=0.9, right, font=\tiny, poporange] {Seuil détection (\SI{50}{copies/mL})};

% CD4 (échelle secondaire approximative pour affichage)
\addplot[popcurve, popgreen, very thick, domain=0:52, samples=200] 
    ({x}, 
     {2.0 + 1.5*tanh((x-2)/3) + 1.0*tanh((x-20)/20)} 
    ) 
    node[pos=0.2, below, font=\tiny, popgreen] {Phase 1 : Redistribution}
    node[pos=0.8, above, font=\tiny, popgreen] {Phase 2 : Thymopoïèse / Prolifération};

\legend{Charge Virale (Log$_{10}$), Seuil indétectable, CD4 (échelle Log arbitraire pour illustration)}
\end{axis}
\end{tikzpicture}
\legende{Kinétique type de la réponse virologique et immunitaire sous ARV efficaces (ex: TLD). \textbf{Phase 1 (0-2 à 4 sem) :} Chute exponentielle rapide de la CV (clairance virions libres, demi-vie $\approx$ 6h) et remontée rapide des CD4 (redistribution depuis tissus/ganglions). \textbf{Phase 2 (mois - années) :} Chute lente de la CV résiduelle (cellules infectées à longue durée de vie, macrophages) jusqu'à indétectable ($<\SI{50}{copies/mL}$). Remontée CD4 lente et progressive (néogenèse thymique, prolifération périphérique), souvent incomplète si CD4 initial bas ($<$\SI{200}{\per\micro\liter}).}
\end{popfigure}

\begin{definition}[Critères d'échec virologique (OMS / Cameroun 2024)]
L'échec virologique est défini par une \textbf{Charge Virale $> \SI{1000}{copies/mL}$} à \textbf{deux reprises espacées de 3 mois} (M+3 après détection initiale), \textbf{après au moins 6 mois de traitement}, \textbf{sous réserve d'une bonne observance confirmée} (entretien d'éducation thérapeutique, comptage pilules, dosage plasmatique ARV si doute).
\begin{itemize}
    \item \textbf{CV \SIrange{50}{999}{copies/mL} :} "Viremie bas niveau" (Low-level viremia). \textbf{Pas d'échec}. Renforcer observance, contrôler à 3-6 mois. Rarement évolutif vers résistance sous DTG.
    \item \textbf{CV $> \SI{1000}{copies/mL}$ confirmée :} \textbf{Échec virologique avéré}. Risque majeur de sélection de résistances. $\to$ Demande \textbf{Génotypage (Test de résistance)} $\to$ Passage en 2ème ligne adapté.
\end{itemize}
\end{definition}

\begin{exoresolu}[Interprétation d'un suivi virologique]
\textbf{Cas :} Homme 35 ans, sous TLD depuis 14 mois. CV M0 = \SI{120000}{copies/mL}. CV M6 = \SI{45}{copies/mL} (Indétectable). CV M12 = \SI{1200}{copies/mL}. Le patient avoue avoir "oublié souvent" ses prises le mois dernier suite à un deuil.
\textbf{Question :} Quelle conduite tenir ?
\textbf{Solution :}
\begin{enumerate}
    \item \textbf{Ne pas crier à l'échec tout de suite.} La CV M12 est $> \SI{1000}{copies/mL}$ mais le patient rapporte une observance défaillante récente (cause identifiée, réversible).
    \item \textbf{Action immédiate :} Conseil d'observance intensif (éducation thérapeutique, boîte à pilules, soutien psychosocial). Reprise stricte du traitement.
    \item \textbf{Contrôle :} \textbf{CV de contrôle à M+3 (soit M15)}.
    \item \textbf{Scénarios :}
        \begin{itemize}
            \item CV M15 $< \SI{50}{copies/mL}$ : \textbf{"Blip" ou échec lié à l'observance}. Revenir au suivi annuel. Succès.
            \item CV M15 $> \SI{1000}{copies/mL}$ \textbf{malgré observance retrouvée} : \textbf{Échec virologique confirmé}. Demander Génotypage. Préparer switch 2ème ligne (DTG + 2 INTI neufs, ex: AZT+3TC).
        \end{itemize}
\end{enumerate}
\textbf{Enseignement :} Une seule CV $> \SI{1000}{copies/mL}$ ne suffit pas pour déclarer l'échec. La confirmation à 3 mois \textbf{sous bonne observance} est la règle d'or pour éviter les changements de ligne inutiles et coûteux.
\tcblower
\textbf{Réponse détaillée :} La conduite à tenir est le renforcement de l'observance suivi d'un contrôle à 3 mois. Le passage en 2ème ligne n'est décidé qu'en cas de persistance de la virémie $> \SI{1000}{copies/mL}$ malgré une observance prouvée $> 95\%$.
\end{exoresolu}

\courssub{Observance : Le pilier de la réussite thérapeutique}

\begin{definition}[Observance thérapeutique (Adhérence)]
C'est le \textbf{respect strict de la prescription médicale} : \textbf{Bon médicament} + \textbf{Bonne dose} + \textbf{Bon moment} + \textbf{Bon régime alimentaire (si applicable)} + \textbf{Durée continue}.
Le \textbf{seuil critique} pour la suppression virale durable sous trithérapie standard est \textbf{$> 95\%$ des prises effectuées} (soit $\le$ 3 doses oubliées / mois pour une bithérapie/jour, ou 1 dose/jour pour TLD). Sous DTG (haute barrière), une marge légèrement plus large existe, mais l'objectif reste 100\%.
\end{definition}

\begin{aretenir}[Facteurs favorisant / freinant l'observance au Cameroun]
\begin{center}
\begin{tabularx}{\linewidth}{|X|X|}\hline
\cellcolor{popgreenL} \textbf{Facteurs FAVORISANTS (Leviers)} & \cellcolor{poporangeL} \textbf{Facteurs FREINANTS (Barrières)} \\ \hline
\begin{itemize}\item Traitement simplifié : 1 cp/jour (TLD), petite taille.
\item Gratuité totale des ARV et examens (Fonds Mondial, État, PEPFAR).
\item Groupes de pairs (GAP, GAS), Mentors "Paires Éducatrices".
\item Outils : Boîtes à pilules (hebdo), Alarmes téléphone, Applications (Medisafe).
\item Éducation thérapeutique continue (ETP) : Comprendre \textbf{POURQUOI} (CV, CD4, TasP).
\item Intégration services : Dépistage TB, Grossesse, Planning familial au même lieu.
\item Soutien familial / conjoint (si statut partagé).
\end{itemize} & 
\begin{itemize}\item Stigmatisation / Discrimination (peur d'être vu à la pharmacie/CS).
\item Effets indésirables précoces (nausées EFV/DTG 1ères sem, rêves) non gérés.
\item Alcool / Drogues / Santé mentale non prise en charge (dépression).
\item Instabilité vie : Déplacements, perte emploi, conflits conjugaux.
\item Ruptures de stock (stock-out) ARV ou réactifs CV (système de santé).
\item Croyances : Médecine traditionnelle, "Guérison divine" $\to$ Arrêt ARV.
\item Fatigue thérapeutique ("Pill burden" psychologique après années).
\end{itemize} \\ \hline
\end{tabularx}
\end{center}
\end{aretenir}

\begin{attention}[Erreur fréquente au Probatoire et en vie réelle : "J'arrête, je me sens bien / CV indétectable"]
\textbf{Affirmation fausse :} "Puisque ma charge virale est indétectable et que je n'ai plus de symptômes, je peux arrêter les médicaments."
\textbf{Réalité scientifique :} L'arrêt expose à :
\begin{enumerate}
    \item \textbf{Rebond viral rapide :} CV remonte à $> 10^5$ cp/mL en 2-4 semaines (réactivation réservoirs).
    \item \textbf{Sélection de résistances :} Pendant la phase de remontée, le virus réplique sous pression résiduelle du médicament (demi-vie longue EFV, DTG) $\to$ sélection mutants résistants $\to$ \textbf{Échec définitif de la 1ère ligne}.
    \item \textbf{Chute CD4 brutale :} Syndrome de reconstitution immunitaire inflammatoire (SRII) à la reprise, ou aggravation clinique.
    \item \textbf{Risque transmission :} Perte du bénéfice TasP.
\end{enumerate}
\textbf{Message clé :} \textbf{ARV = À VIE.} Tout arrêt doit être une décision médicale collégiale (essai de rémission encadré).
\end{attention}

\courssub{Gestion des effets indésirables et toxicités : Anticiper pour ne pas arrêter}

\begin{exemplebox}[Conduite à tenir face aux effets indésirables majeurs (Exemples Cameroun)]
\begin{itemize}
    \item \textbf{DTG (Insomnie, prise de poids, hyperglycémie) :} Prendre le \textbf{matin} (pas le soir). Surveiller poids/glycémie à M6, M12, puis annuel. Si intolérance neuropsychiatrique sévère $\to$ Switch vers DRV/r ou EFV 400mg (rare).
    \item \textbf{TDF (Rénal : Protéinurie, Créatinine $\uparrow$ ; Os : Ostéopénie) :} Surveiller Créatinine/DFG à M3, M6, M12. Si DFG $< \SI{50}{mL/min}$ ou protéinurie $\to$ \textbf{Switch TDF $\to$ ABC} (ou TAF si dispo, pas encore standard Cameroun).
    \item \textbf{AZT (Anémie, Neutropénie) :} NFS mensuelle 3 premiers mois, puis trimestrielle. Si Hb $< \SI{8}{g/dL}$ ou Neutrophiles $< \SI{750}{\per\cubic\milli\meter}$ $\to$ Switch AZT $\to$ TDF ou ABC.
    \item \textbf{EFV (Neuropsychiatrique : Vertiges, cauchemars, dépression, suicidaire) :} Prendre au coucher, éviter alcool. Si persistant $> 2-4$ sem ou sévère $\to$ Switch EFV $\to$ DTG (standard).
    \item \textbf{LPV/r (Diarrhée chronique, Dyslipidémie, Diabète) :} Hygiène diététique, lopéramide si besoin. Lipidogramme/Glycémie annuel. Si intolérance digestive majeure $\to$ Switch vers DRV/r (mieux toléré) ou DTG (si pas résistance INNTI).
\end{itemize}
\textbf{Règle d'or :} \textbf{Jamais arrêter un ARV seul.} Toujours \textbf{switcher} (substituer la molécule fautive) pour maintenir la pression virale triple.
\end{exemplebox}

\courssub{Aspects économiques, psychosociaux et riposte nationale}

\begin{exemplebox}[Coût et Accès : La gratuité effective au Cameroun]
\begin{itemize}
    \item \textbf{1ère ligne (TLD) :} Coût achat génériques $\approx$ \textbf{\SI{60-70}{USD/an/patient}} ($\approx$ \SI{35000-40000}{FCFA}).
    \item \textbf{2ème ligne (AZT+3TC+LPV/r ou DRV/r) :} Coût $\approx$ \textbf{\SI{200-300}{USD/an}} (3 à 5 fois plus cher).
    \item \textbf{3ème ligne (Nouvelles molécules : Lenacapavir, Fostemsavir) :} Coût $\gg$ \SI{1000}{USD/an} (accès via dons/accords volontaires, comité d'expert).
    \item \textbf{Financement :} \textbf{Gratuité totale pour le patient} (ARV + Examens de suivi CV/CD4/Biologie) grâce à la conjugaison : \textbf{Fonds Mondial} (principal bailleur), \textbf{Budget de l'État} (ligne budgétaire PNLS), \textbf{PEPFAR} (USAID/CDC - appui technique, achat réactifs, renfort RH), \textbf{Partenaires ONG}.
    \item \textbf{Enjeu :} Pérenniser le financement domestique (Budget État) face à la baisse des dons internationaux. Lutte contre les ruptures de stock (chaîne d'approvisionnement CENAME $\to$ Régions $\to$ Districts $\to$ CS).
\end{itemize}
\end{exemplebox}

\begin{savaistu}[Recherche : Vers la rémission et le vaccin ?]
\begin{itemize}
    \item \textbf{Guérison fonctionnelle (Contrôle post-traitement) :} Cas rares (Cohorte \textbf{Visconti} France, \textbf{Mississippi Baby} USA, \textbf{Patient de Berlin/Londres/Düsseldorf} - greffe moelle $\Delta$CCR5). Points communs : \textbf{Traitement très précoce} (phase aiguë, Fiebig I-II), \textbf{Durée ARV longue} (années), \textbf{Arrêt encadré}. Le virus reste présent (réservoir) mais contrôlé par le système immunitaire sans traitement. \textbf{Reste exceptionnel ($<1\%$)}.
    \item \textbf{Vaccin préventif :} Échecs successifs (STEP, HVTN 505, Uhambo, Mosaico). Pistes actuelles : \textbf{Vaccins ARNm} (Moderna, IAVI - essais phase I), \textbf{Immunogènes "germline-targeting"} pour induire anticorps largement neutralisants (\textbf{bNAbs}).
    \item \textbf{PrEP injectable (Cabotégravir LA - CAB-LA) :} Injection tous les 2 mois. Efficacité $> 90\%$ (HPTN 083/084). Approuvé OMS/US/UE. \textbf{Arrivée progressive au Cameroun} (projets démonstration, coûts, chaîne de froid).
    \item \textbf{Lenacapavir (Sunlenca) :} Inhibiteur capside (nouvelle cible). Injection sous-cutanée \textbf{tous les 6 mois}. Efficacité majeure en PrEP (PURPOSE 1/2) et 3ème ligne. Espoir majeur pour simplifier vie patients.
\end{itemize}
\end{savaistu}

\begin{exercice}[Cas clinique intégré : Suivi et décision thérapeutique]
\textbf{Patiente :} Mme M., 28 ans, vendeuse au Marché Central de Douala. Découverte séropositivité VIH-1 lors CPN (Grossesse 14 SA). CD4 = \SI{320}{\per\micro\liter}. CV = \SI{85000}{copies/mL}. Créatinine = \SI{8}{mg/L} (Normale). Hb = \SI{11}{g/dL}. AgHBs négatif. Pas de TB.
\begin{enumerate}
    \item Quelle 1ère ligne initier ? Justifier le choix au regard de la grossesse et des guidelines 2024.
    \item Planifier le suivi biologique jusqu'à M12 (dates, examens).
    \item À M6, CV = \SI{35}{copies/mL}, CD4 = \SI{480}{\per\micro\liter}. La patiente arrête le traitement à M8 ("fatiguée de prendre des médicaments, elle se sent bien"). Elle revient à M14. CV = \SI{45000}{copies/mL}. CD4 = \SI{210}{\per\micro\liter}. Interpréter. Conduite à tenir ?
    \item Quel est le risque pour l'enfant (né à terme, allaitement exclusif prévu 6 mois) si la mère n'avait pas repris le traitement ?
\end{enumerate}
\end{exercice}

\begin{corrige}[Exercice : Cas clinique Mme M.]
\begin{enumerate}
    \item \textbf{1ère ligne : TLD (TDF \SI{300}{mg} + 3TC \SI{300}{mg} + DTG \SI{50}{mg}) - 1 cp/jour.}
    \textbf{Justification :} Standard 1ère ligne OMS/Cameroun 2024. DTG recommandé \textbf{y compris au 1er trimestre} (données rassurantes sur défaut tube neural), à condition d'assurer une contraception efficace post-partum si désir d'espacement. Alternative TLE400 (EFV 400mg) si refus DTG ou contre-indication rare. TDF ok (Fonction rénale normale). ABC non nécessaire.
    \item \textbf{Suivi :}
    \begin{itemize}
        \item M0 (Initiation) : Déjà fait (BPT).
        \item M3 : CV + Créatinine (surveillance TDF).
        \item M6 : CV + CD4 + Créatinine. (Objectif CV $< \SI{50}{copies/mL}$).
        \item M12 (Post-partum ~6 mois) : CV + CD4 + Lipides + Glycémie + NFS.
    \end{itemize}
    \item \textbf{Interprétation M14 :} \textbf{Reprise virale massive (Rebond)} due à l'arrêt du traitement (Non-observance volontaire). CV remonte au niveau pré-traitement (réservoir réactivé). Chute CD4 (destruction cellules cibles). \textbf{Risque résistance :} Faible si arrêt brutal de \textbf{tous} les ARV simultanément (pas de monothérapie fonctionnelle), mais mutations archivées possibles si reprise tardive.
    \textbf{Conduite :} 1. Accueil non jugeant + Éducation thérapeutique (rappel I=I, vie normale). 2. \textbf{Reprise immédiate TLD} (même ligne, pas de résistance attendue). 3. Contrôle CV à M+1 (M15) puis M+3 (M17) pour confirmer re-suppression. 4. Soutien psychosocial (Groupe de pairs, boîte pilules).
    \item \textbf{Risque enfant :} \textbf{Très élevé.} Sans ARV maternel (CV $> \SI{45000}{copies/mL}$), risque transmission \textbf{in utero / partum / post-partum (allaitement)} $\approx$ 15-25\% sans prophylaxie, jusqu'à 35-45\% avec allaitement prolongé non traité. L'enfant doit recevoir \textbf{Prophylaxie ARV à la naissance (AZT+3TC+NVP si risque élevé)} pendant 6-12 semaines + Dépistage virologique précoce (PCR ADN VIH) à 6 sem, 9 mois, fin allaitement.
\end{enumerate}
\end{corrige}

\coursec{Aspects psychosociaux, droits et riposte nationale}

Dans la section précédente, nous avons vu que les antirétroviraux (ARV) permettent aujourd'hui à une personne vivant avec le VIH (PVVIH) de mener une vie longue et productive, à condition d'être dépistée tôt, mise sous traitement et maintenue en observance. Mais la biologie ne suffit pas. Une jeune femme de Douala qui craint d'être chassée de chez elle si son statut est découvert ne se rendra jamais au centre de dépistage ; un patient stigmatisé à l'hôpital interrompra son traitement. La lutte contre le VIH/Sida est donc aussi un combat \emph{social, juridique et politique}. C'est l'objet de cette dernière section : comprendre les obstacles psychosociaux, connaître les droits des patients, et décrire l'organisation de la riposte nationale camerounaise.

\courssub{Stigmatisation et discrimination : l'ennemi social de la riposte}

\textbf{Situation concrète.} Dans un quartier de Yaoundé, un commerçant apprend qu'il est séropositif. Il cache ses comprimés, arrête de se rendre à l'hôpital de jour, et finit par interrompre son traitement de peur d'être reconnu. Deux ans plus tard, il développe un Sida déclaré. Ce n'est pas le virus qui a gagné du terrain : c'est la \emph{peur du regard des autres}.

\begin{definition}[Stigmatisation et discrimination]
La \cle{stigmatisation} (concept du sociologue Erving Goffman, 1963) est un attribut profondément discréditant qui réduit la personne « entière et habituelle » à une personne « souillée, écartée ». Appliquée au VIH, elle se manifeste par la peur, le rejet (famille, travail, école, services de soins) et les rumeurs morales. La \cle{discrimination} est le passage à l'acte : un traitement différentiel et injuste fondé sur le statut sérologique réel ou perçu (licenciement, exclusion scolaire, refus de soins). On distingue aussi l'\cle{auto-stigmatisation} : la personne intériorise les jugements négatifs, perd l'estime d'elle-même et s'isole d'elle-même.
\end{definition}

Les conséquences sanitaires de la stigmatisation sont directement mesurables et contrecarrent chaque étape de la cascade de soins :

\begin{itemize}
\item \textbf{Non-dépistage} : la peur du résultat et de la divulgation dissuade de faire le test ; l'infection est découverte tardivement, au stade Sida.
\item \textbf{Statut caché} : la personne dissimule sa maladie, y compris à son partenaire, ce qui favorise les transmissions.
\item \textbf{Non-observance} : prendre un comprimé quotidien devient difficile quand on cache son traitement ; d'où des échecs thérapeutiques et des résistances.
\item \textbf{Dépression et isolement} : la détresse psychologique aggrave la non-observance et dégrade la qualité de vie.
\end{itemize}

\textbf{Les leviers de lutte.} Le Cameroun s'est doté d'un cadre juridique et communicationnel :
\begin{itemize}
\item la \textbf{loi n° 2007/004} portant protection des personnes vivant avec le VIH, qui interdit toute discrimination fondée sur le statut sérologique ;
\item les campagnes de \textbf{sensibilisation} (médias, écoles, églises et mosquées, leaders d'opinion, témoignages publics de PVVIH) qui « normalisent » le dépistage et le traitement ;
\item la formation du personnel de santé à l'accueil sans jugement.
\end{itemize}

\begin{attention}[Ce qui est illégal au Cameroun]
Erreur fréquente — et faute grave : croire que le VIH justifie une rupture de contrat de travail, une exclusion scolaire ou un refus de soins chirurgicaux. C'est \textbf{illégal} : la loi n° 2007/004, le Code du Travail et les décrets d'application interdisent explicitement ces discriminations. Un employeur, un directeur d'école ou un chirurgien qui agirait ainsi s'expose à des poursuites. Le statut sérologique relève du secret médical et ne peut être exigé à l'embauche ou à l'inscription scolaire.
\end{attention}

\courssub{Populations clés : une vulnérabilité concentrée}

\textbf{Observation épidémiologique.} Les enquêtes biologiques et comportementales (IBBS) montrent que l'épidémie camerounaise est \emph{concentrée} : certains groupes présentent des prévalences 5 à 10 fois supérieures à celle de la population générale.

\begin{definition}[Populations clés]
Les \cle{populations clés} (KP, \emph{Key Populations}) sont des groupes à risque accru d'exposition au VIH en raison de leurs pratiques ou de leur situation sociale : les HSH (hommes ayant des rapports sexuels avec des hommes), les travailleuses du sexe (TS), les usagers de drogues par injection (UDI), les personnes transgenres et les détenus.
\end{definition}

\begin{popfigure}
\begin{center}
\begin{tikzpicture}
\begin{axis}[
    ybar, bar width=0.55cm,
    width=0.95\linewidth, height=6.2cm,
    ymin=0, ymax=30,
    ylabel={Prévalence VIH (\%)},
    symbolic x coords={{Pop. générale}, HSH, TS, UDI, Détenus},
    xtick=data,
    x tick label style={font=\small, rotate=12, anchor=east},
    ymajorgrids, grid style={popline!40},
    every axis/.append style={font=\small},
]
\addplot[fill=popblue, draw=popdark] coordinates {({Pop. générale},2.7) (HSH,20.5) (TS,24.3) (UDI,12.0) (Détenus,6.0)};
\end{axis}
\end{tikzpicture}
\end{center}
\legende{Prévalence du VIH par groupe de population au Cameroun (ordres de grandeur issus des enquêtes IBBS et EDS récentes). Les populations clés (HSH, TS, UDI) présentent des prévalences 5 à 10 fois supérieures à celle de la population générale.}
\end{popfigure}

\textbf{Pourquoi ces groupes sont-ils si exposés ?} Aux facteurs biologiques (rapports anaux réceptifs plus à risque, partage de seringues) s'ajoutent des \emph{barrières structurelles} : criminalisation de certaines pratiques, stigmatisation violente, accès difficile aux soins (peur d'être jugé ou dénoncé au centre de santé). Résultat : ces populations se dépistent moins, sont mises sous ARV plus tard et constituent des réservoirs épidémiques qui alimentent la transmission vers la population générale.

\begin{methode}[L'approche « Key Populations Friendly Services »]
Pour lever ces barrières, la riposte camerounaise déploie des services adaptés :
\begin{enumerate}
\item Former le personnel de santé à l'accueil \textbf{sans jugement} et à la confidentialité renforcée.
\item Créer des \textbf{plages horaires ou des sites dédiés} (cliniques communautaires, unités mobiles).
\item Associer des \textbf{pairs éducateurs} issus des communautés elles-mêmes, qui inspirent confiance.
\item Offrir un paquet complet : dépistage, préservatifs et lubrifiants, PrEP, traitement des IST, réduction des risques pour les UDI (programmes d'échange de seringues là où ils existent).
\end{enumerate}
\end{methode}

\courssub{Genre et VIH : la vulnérabilité des femmes et des filles}

\textbf{Situation concrète.} Une fille de 16 ans, mariée précocement à un homme plus âgé, ne peut ni négocier l'usage du préservatif ni refuser les rapports. Sa vulnérabilité n'est pas un choix : elle est \emph{construite socialement}.

Les femmes et les filles sont plus exposées au VIH pour des raisons à la fois \textbf{biologiques} (la muqueuse vaginale offre une grande surface d'exposition ; le risque de transmission homme $\to$ femme est supérieur au risque inverse) et \textbf{sociales} : violences basées sur le genre (VBG), dépendance économique, mariages précoces, inégalité d'accès à l'éducation et à l'information.

\textbf{La PTME et l'option B+.} Sans intervention, une mère séropositive non traitée transmet le VIH à son enfant dans 15 à 45 \% des cas (grossesse, accouchement, allaitement). La \cle{PTME} (prévention de la transmission mère-enfant, en anglais PMTCT) a révolutionné ce pronostic avec l'\textbf{option B+} : \emph{toute} femme enceinte ou allaitante séropositive est mise sous ARV \textbf{à vie}, quel que soit son taux de CD4. Sous traitement efficace, le risque de transmission tombe \textbf{sous 2 \%}.

\begin{methode}[La notification assistée du partenaire (NAP / APN)]
Quand une personne est dépistée positive, ses partenaires et ses enfants doivent être proposés au test — sans violence ni rupture. Protocole en 4 étapes :
\begin{enumerate}
\item \textbf{Conseiller} le patient sur l'importance d'informer son ou ses partenaire(s) (bénéfice santé pour tous).
\item \textbf{Proposer des options} : le patient informe lui-même (avec soutien du conseiller) ; le personnel de santé informe de façon anonyme et confidentielle ; ou un « contrat » : le patient s'engage à le dire sous X jours, sinon le personnel le fait.
\item \textbf{Dépister} le(s) partenaire(s) et les enfants.
\item \textbf{Lier vers la prise en charge} : ARV si positif, PrEP si négatif mais exposé (couple sérodifférent).
\end{enumerate}
\end{methode}

\begin{propriete}[Principe « Ne pas nuire » (\emph{Do No Harm})]
Toute intervention liée au VIH (dépistage, notification, divulgation) doit évaluer et atténuer au préalable les risques de violence, de rupture conjugale ou d'exclusion sociale que la révélation du statut peut déclencher, en particulier pour les femmes. La divulgation n'est jamais forcée : elle est préparée, accompagnée et sécurisée.
\end{propriete}

\begin{popfigure}
\begin{center}
\begin{tikzpicture}[scale=1.1]
% Cercles concentriques
\fill[popblue!10] (0,0) circle (3.6);
\fill[popgreen!12] (0,0) circle (2.7);
\fill[poporange!15] (0,0) circle (1.8);
\fill[popgold!25] (0,0) circle (0.9);
\draw[popdark, thick] (0,0) circle (3.6);
\draw[popdark, thick] (0,0) circle (2.7);
\draw[popdark, thick] (0,0) circle (1.8);
\draw[popdark, thick] (0,0) circle (0.9);
\node[font=\small\bfseries, text=popdark] at (0,0) {Individu};
\node[font=\scriptsize, text=popdark] at (0,-0.4) {(âge, sexe,} ;
\node[font=\scriptsize, text=popdark] at (0,-0.65) {comportements)};
\node[font=\scriptsize\bfseries, text=popdark] at (0,1.3) {Relationnel};
\node[font=\scriptsize, text=popdark, align=center] at (0,-1.3) {couple, famille,\\pairs, VBG};
\node[font=\scriptsize\bfseries, text=popdark] at (0,2.2) {Communautaire};
\node[font=\scriptsize, text=popdark, align=center] at (0,-2.2) {école, travail,\\quartier, soins};
\node[font=\scriptsize\bfseries, text=popdark] at (0,3.1) {Sociétal / Structurel};
\node[font=\scriptsize, text=popdark, align=center] at (0,-3.1) {lois, normes sociales,\\pauvreté, criminalisation};
\end{tikzpicture}
\end{center}
\legende{Modèle socio-écologique des déterminants de la vulnérabilité au VIH : le risque individuel est imbriqué dans des niveaux relationnel, communautaire et structurel (lois, normes, pauvreté). Une riposte efficace doit agir sur tous les niveaux.}
\end{popfigure}

\courssub{La riposte nationale camerounaise : organisation et financement}

\textbf{Du premier cas au plan stratégique.} La riposte camerounaise s'est construite par étapes, du constat initial à une stratégie multisectorielle coordonnée.

\begin{popfigure}
\begin{center}
\begin{tikzpicture}[scale=1]
\draw[popdark, very thick, -{Stealth}] (0,0) -- (13.6,0);
\foreach \x/\an/\txt in {
  0.4/1985/{1\textsuperscript{er} cas de Sida\\notifié au Cameroun},
  2.6/2000/{Création du CNLS},
  4.8/2007/{Loi 2007/004 de\\protection des PVVIH},
  7.0/2017/{Gratuité des ARV\\et du dépistage},
  9.2/2019/{Auto-dépistage\\autorisé},
  11.4/2022/{TLD en 1\textsuperscript{re} ligne\\(Dolutégravir)},
  13.2/2024/{PSN 2024--2028}
}{
  \fill[poporange] (\x,0) circle (0.09);
  \draw[popdark] (\x,0) -- (\x,0.35);
  \node[above, font=\scriptsize\bfseries, text=popblue] at (\x,0.35) {\an};
  \node[below, font=\scriptsize, text=popdark, align=center] at (\x,-0.15) {\txt};
}
\end{tikzpicture}
\end{center}
\legende{Chronologie simplifiée des grands jalons de la riposte nationale contre le VIH/Sida au Cameroun (1985--2024). TLD : Ténofovir + Lamivudine + Dolutégravir, trithérapie de référence en première ligne.}
\end{popfigure}

\begin{definition}[CNLS et PSN]
Le \cle{CNLS} (Comité National de Lutte contre le Sida), créé en 2000 et placé sous l'autorité de la Présidence de la République, coordonne la riposte multisectorielle (santé, éducation, justice, communication, associations). Le \cle{PSN} (Plan Stratégique National) 2024--2028 fixe les objectifs : atteindre les « 95--95--95 » (95 \% des PVVIH connaissent leur statut, 95 \% d'entre elles sous ARV, 95 \% de celles-ci avec charge virale indétectable) et éliminer la transmission mère-enfant.
\end{definition}

\textbf{Financement et gratuité.} La riposte repose sur quatre sources : le budget de l'État, le Fonds Mondial de lutte contre le Sida, la Tuberculose et le Paludisme, le PEPFAR (plan d'urgence américain) et d'autres partenaires (ONUSIDA, OMS). Grâce à ce financement, le \textbf{dépistage, les ARV et les examens de suivi (CD4, charge virale, bilans biologiques) sont gratuits} dans les structures publiques et conventionnées.

\begin{exemplebox}[Un investissement rentable]
L'effort national de lutte contre le VIH représente environ \textbf{150 à 200 milliards de FCFA par an}. Ce coût est en réalité un investissement : selon la Banque Mondiale, chaque dollar investi dans la riposte rapporte environ \textbf{6 à 10 dollars} en productivité préservée, en coûts hospitaliers évités et en orphelins évités. Traiter tôt coûte infiniment moins cher que soigner un Sida déclaré.
\end{exemplebox}

\begin{savaistu}[La différenciation de la prise en charge (DSD)]
Le Cameroun a adopté la stratégie DSD (\emph{Differentiated Service Delivery}) pour les patients stables : \textbf{groupes de soutien communautaire} (CAG : un membre du groupe récupère les ARV pour tous), \textbf{distribution multi-mois} (MMD : 3 à 6 mois d'ARV en une seule visite) et \textbf{cliniques « fast-track »} (file rapide pour les patients stables). Résultat : hôpitaux décongestionnés, meilleure rétention dans les soins, et les patients ne perdent plus une journée de travail chaque mois à l'hôpital.
\end{savaistu}

\courssub{Les droits du patient : un bouclier juridique et éthique}

La riposte efficace repose sur la confiance, et la confiance repose sur des droits garantis :

\begin{itemize}
\item \textbf{Confidentialité médicale absolue} : le statut sérologique est couvert par le secret médical ; aucun personnel de santé ne peut le révéler (employeur, famille, assurance) sans consentement.
\item \textbf{Consentement éclairé} : le test VIH est \emph{volontaire} ; il ne peut être obligatoire, sauf exceptions légales (don du sang, où le dépistage systématique protège le receveur, et certaines filières militaires).
\item \textbf{Non-discrimination} : accès égal aux soins, à l'emploi et à l'éducation, garanti par la loi 2007/004.
\item \textbf{Droit à l'information} : le patient a accès à son dossier médical et à des explications claires sur son traitement.
\end{itemize}

\begin{exoresolu}[Analyser une situation de discrimination]
\textbf{Énoncé.} Une infirmière séropositive, sous ARV avec charge virale indétectable, est licenciée par une clinique privée de Bafoussam qui a appris son statut par un collègue. Le directeur justifie : « c'est pour protéger les patients ». Analyse cette situation sur les plans scientifique, éthique et juridique, et propose deux actions.
\tcblower
\textbf{Solution détaillée.}
\begin{enumerate}
\item \emph{Plan scientifique} : une personne sous ARV avec charge virale indétectable ne transmet pas le virus sexuellement (I = I, « indétectable = intransmissible ») et ne présente \textbf{aucun risque} dans l'exercice des soins courants, les précautions standard (gants, gestion des aiguilles) s'appliquant à tout le monde. L'argument du directeur est donc scientifiquement faux.
\item \emph{Plan éthique} : le collègue a violé la confidentialité médicale ; la clinique ajoute une discrimination à cette violation, ce qui entretient la stigmatisation et décourage le dépistage des soignants.
\item \emph{Plan juridique} : le licenciement fondé sur le statut sérologique viole la loi n° 2007/004 et le Code du Travail ; il est nul et ouvre droit à réparation.
\item \emph{Actions possibles} : saisir l'inspection du travail et/ou les tribunaux pour licenciement abusif ; porter plainte contre la divulgation du secret médical ; s'appuyer sur une association de PVVIH et le CNLS pour l'accompagnement juridique et psychosocial.
\end{enumerate}
\textbf{Conclusion} : la connaissance des droits est un outil de santé publique — elle protège les personnes et rend la riposte possible.
\end{exoresolu}

\courssub{Nouvelles perspectives scientifiques}

La recherche avance vite ; quatre pistes transforment déjà ou transformeront la lutte :

\begin{itemize}
\item \textbf{PrEP injectable longue durée} : le \cle{Cabotégravir LA}, injecté \textbf{tous les 2 mois}, protège les personnes très exposées sans contrainte de comprimé quotidien — une arme majeure pour les populations clés.
\item \textbf{Anneau vaginal à la Dapivirine} : dispositif discret, contrôlé par la femme elle-même, libérant un ARV localement pendant un mois ; il répond directement au problème de la négociation du préservatif.
\item \textbf{Vaccins} : des candidats vaccins préventifs et thérapeutiques sont en essais cliniques, même si la grande variabilité du VIH complique leur mise au point.
\item \textbf{Stratégies d'éradication des réservoirs} : « \emph{Shock and Kill} » (réveiller le virus dormant dans les cellules réservoirs pour le détruire) et « \emph{Block and Lock} » (verrouiller durablement le provirus en état silencieux) visent une guérison fonctionnelle.
\end{itemize}

\begin{aretenir}[L'essentiel de la section]
\begin{itemize}
\item La \textbf{stigmatisation} (attribut discréditant) et la \textbf{discrimination} (acte injuste) bloquent le dépistage, l'observance et la rétention : elles sont combattues par la loi 2007/004 et la sensibilisation.
\item Les \textbf{populations clés} (HSH, TS, UDI, transgenres, détenus) ont des prévalences 5 à 10 fois supérieures ; elles exigent des services adaptés et sans jugement.
\item Les \textbf{femmes et filles} sont plus vulnérables (biologie + inégalités) ; l'option B+ protège mère et enfant, la notification assistée protège les partenaires dans le respect du principe « ne pas nuire ».
\item La riposte camerounaise est coordonnée par le \textbf{CNLS} (PSN 2024--2028), financée par l'État et ses partenaires, et garantit la \textbf{gratuité} du dépistage, des ARV et du suivi.
\item Les droits du patient — \textbf{confidentialité, consentement éclairé, non-discrimination, accès à l'information} — sont le socle juridique de la confiance.
\item Les innovations (PrEP injectable, anneau vaginal, stratégies d'éradication) ouvrent la voie vers la fin de l'épidémie.
\end{itemize}
\end{aretenir}

\newpage

\coursec{Activité d'intégration}

\courssub{Objectifs de l'activité}

Cette activité d'intégration te place dans la peau d'un membre de l'équipe de santé d'un hôpital camerounais. À partir d'un dossier patient anonymisé, tu devras mobiliser \emph{tous} les savoirs et savoir-faire de la leçon pour :

\begin{itemize}
\item expliquer et interpréter un résultat de dépistage du VIH (algorithme à trois tests) ;
\item interpréter des données biologiques simples (taux de \cle{CD4}, \cle{charge virale}) et situer le stade clinique et immunologique de l'infection ;
\item évaluer un risque de \cle{transmission mère-enfant} (TME) et l'effet de la prophylaxie ;
\item justifier le choix d'un traitement \cle{antirétroviral} (ARV) de première ligne ;
\item argumenter en faveur du \cle{dépistage volontaire}, de la prévention du partenaire et de la prise en charge précoce ;
\item identifier les enjeux psychosociaux, juridiques et communautaires de l'annonce d'une séropositivité.
\end{itemize}

\begin{aretenir}[Compétence visée]
Intégrer les connaissances sur le cycle viral, la transmission, l'évolution de l'infection, le dépistage et les ARV pour résoudre une situation-problème authentique de santé publique, en adoptant une attitude responsable, solidaire et respectueuse des droits des personnes vivant avec le VIH (PVVIH).
\end{aretenir}

\courssub{Situation}

\textbf{Dossier patient n\textsuperscript{o} 2024-117 — Hôpital de District de Bafoussam (Région de l'Ouest, Cameroun).}

Marie, 24 ans, couturière, se présente pour la première consultation prénatale (CPN1) de sa grossesse, estimée à 14 semaines d'aménorrhée (SA). Conformément à la politique nationale de \cle{PTME} (Prévention de la Transmission Mère-Enfant), un dépistage du VIH lui est proposé, avec son consentement éclairé. Elle accepte.

\textbf{Résultats de l'algorithme national de dépistage (3 tests) :}
\begin{itemize}
\item Test 1 (test rapide de première intention, très sensible) : \textbf{positif} ;
\item Test 2 (test de confirmation, très spécifique) : \textbf{positif} ;
\item Conclusion du laboratoire : \textbf{séropositivité VIH-1 confirmée} (le test 3 n'est requis qu'en cas de discordance T1+/T2$-$).
\end{itemize}

\textbf{Bilan biologique initial de Marie :}
\begin{itemize}
\item Taux de lymphocytes T CD4 : \SI{450}{cellules/mm^3} ;
\item Charge virale (CV) plasmatique : \SI{85 000}{copies/mL} ;
\item État clinique : aucune infection opportuniste, bon état général.
\end{itemize}

Paul, 26 ans, conjoint de Marie, commerçant, accepte à son tour le dépistage proposé au couple : son test est \textbf{négatif}. Le couple vit ensemble depuis 3 ans, sans utilisation systématique du préservatif. Marie craint la réaction de sa belle-famille et hésite à révéler son statut. Paul se demande s'il est « à l'abri » et s'il doit refaire un test.

\textbf{Données de référence utiles (fournies dans le dossier) :}
\begin{itemize}
\item Risque de transmission mère-enfant \emph{sans aucune intervention} : 25 \% à 35 \% (grossesse, accouchement, allaitement confondus) ;
\item Risque de TME \emph{avec prise en charge complète} (ARV de la mère dès la grossesse, accouchement encadré, prophylaxie du nouveau-né, conduite de l'allaitement selon les recommandations nationales) : ramené à \textbf{moins de 2 \%} ;
\item Classification immunologique OMS : CD4 $>$ 500/mm\textsuperscript{3} : absence d'immunodépression significative ; 350--499 : immunodépression légère ; 200--349 : modérée ; $<$ 200 : sévère ;
\item Stratégie nationale : traitement ARV proposé à \emph{toute} PVVIH dès le diagnostic, quels que soient les CD4 (« tester et traiter »), et systématiquement à toute femme enceinte séropositive (option B+) ;
\item Premières lignes disponibles : \textbf{TLD} (Ténofovir + Lamivudine + Dolutégravir) ou \textbf{TLE} (Ténofovir + Lamivudine + Éfavirenz).
\end{itemize}

\courssub{Consignes}

Par binôme, à l'aide du dossier et des ressources de la leçon, réponds aux questions suivantes, puis prépare une affiche de synthèse pour la restitution collective.

\begin{enumerate}
\item \textbf{Interpréter le dépistage de Marie.} Peut-on affirmer sa séropositivité ? Le test 3 était-il nécessaire ? Pourquoi ne peut-on pas évoquer une « fenêtre silencieuse » dans son cas ?
\item \textbf{Évaluer le risque pour l'enfant.} Sur 100 naissances de mères séropositives, combien d'enfants infectés attendrait-on sans intervention ? Avec l'option B+ ? Conclus sur l'intérêt de la PTME.
\item \textbf{Interpréter les marqueurs biologiques de Marie.} Situe son stade immunologique (CD4) et commente sa charge virale. Le traitement ARV est-il urgent ? Justifie.
\item \textbf{Choisir la première ligne ARV.} Entre TLD et TLE, quelle combinaison recommandes-tu pour Marie (grossesse + charge virale élevée) ? Justifie par les propriétés des molécules.
\item \textbf{Conseiller Paul.} Son test négatif signifie-t-il qu'il est définitivement non infecté ? Quelles mesures de prévention lui proposes-tu (préservatif, PrEP, TasP) ? Quand devra-t-il refaire un test ?
\item \textbf{Dimensions psychosociales.} Identifie deux risques psychosociaux pour le couple et propose deux actions concrètes (juridiques, médicales ou communautaires) pour y répondre.
\end{enumerate}

\courssub{Ressources à mobiliser}

\begin{itemize}
\item \textbf{Cycle du VIH} : le virus cible les lymphocytes T CD4 ; la charge virale reflète l'intensité de la réplication ; les ARV bloquent des étapes du cycle (transcription inverse, intégration) et font chuter la charge virale.
\item \textbf{Algorithme de dépistage} : T1 très sensible ; si T1+ $\Rightarrow$ T2 très spécifique ; T1+ et T2+ $\Rightarrow$ séropositivité confirmée ; discordance T1+/T2$-$ $\Rightarrow$ T3 d'arbitrage. La \cle{fenêtre silencieuse} (environ 3 semaines à 3 mois) est la période où une personne infectée peut encore tester négatif.
\item \textbf{Évolution naturelle} : primoinfection $\rightarrow$ séroconversion $\rightarrow$ phase asymptomatique (CD4 $>$ 500, puis déclin progressif) $\rightarrow$ Sida déclaré (CD4 $<$ 200 ou infections opportunites).
\item \textbf{Prévention du partenaire} : préservatif ; \cle{PrEP} (prophylaxie pré-exposition) pour le partenaire séronégatif ; \cle{TasP} (« Treatment as Prevention ») : une charge virale indétectable durable ($<$ 50 copies/mL) sous ARV rend la transmission sexuelle quasi nulle (« I = I », indétectable = intransmissible).
\item \textbf{Cadre juridique camerounais} : confidentialité du statut sérologique, consentement éclairé, non-discrimination des PVVIH.
\end{itemize}

\begin{popfigure}
\begin{tikzpicture}[
    node distance=1.2cm and 1.5cm,
    box/.style={draw=popblue, thick, rounded corners, fill=popblueL, minimum width=2.8cm, minimum height=1cm, align=center, font=\small},
    diamond/.style={draw=poporange, thick, aspect=2, fill=poporangeL, minimum width=3cm, minimum height=1.2cm, align=center, font=\small},
    arrow/.style={-Stealth, thick, popdark},
    start/.style={draw=popgreen, thick, rounded corners, fill=popgreenL, minimum width=2.5cm, minimum height=0.8cm, align=center, font=\small}
]
\node[start, align=center] (start){Dépistage VIH\\ (Consentement éclairé)};
\node[box, below=of start, align=center] (t1){Test 1 (T1)\\ \emph{Très sensible}};
\node[diamond, below=of t1] (dec1) {T1 positif ?};
\node[box, below left=of dec1, xshift=-1.5cm, align=center] (neg){Négatif\\ \textbf{Non infecté}\\(hors fenêtre)};
\node[box, below right=of dec1, xshift=1.5cm, align=center] (t2){Test 2 (T2)\\ \emph{Très spécifique}};
\node[diamond, below=of t2] (dec2) {T2 positif ?};
\node[box, below left=of dec2, xshift=-1.5cm, align=center] (t3){Test 3 (T3)\\ \emph{Arbitrage}};
\node[box, below right=of dec2, xshift=1.5cm, align=center] (pos){\textbf{Séropositivité confirmée}\\ (Prise en charge)};
\node[diamond, below=of t3] (dec3) {T3 positif ?};
\node[box, below left=of dec3, xshift=-1cm, align=center] (neg2){Négatif\\ \textbf{Non infecté}};
\node[box, below right=of dec3, xshift=1cm, align=center] (pos2){\textbf{Séropositivité confirmée}\\ (Prise en charge)};

\draw[arrow] (start) -- (t1);
\draw[arrow] (t1) -- (dec1);
\draw[arrow] (dec1) -- node[left, pos=0.5, font=\small] {Non} (neg);
\draw[arrow] (dec1) -- node[right, pos=0.5, font=\small] {Oui} (t2);
\draw[arrow] (t2) -- (dec2);
\draw[arrow] (dec2) -- node[left, pos=0.5, font=\small] {Non} (t3);
\draw[arrow] (dec2) -- node[right, pos=0.5, font=\small] {Oui} (pos);
\draw[arrow] (t3) -- (dec3);
\draw[arrow] (dec3) -- node[left, pos=0.5, font=\small] {Non} (neg2);
\draw[arrow] (dec3) -- node[right, pos=0.5, font=\small] {Oui} (pos2);
\end{tikzpicture}
\legende{Algorithme national de dépistage du VIH à trois tests (stratégie séquentielle). T1 = test de dépistage (haute sensibilité), T2 = test de confirmation (haute spécificité), T3 = test d'arbitrage en cas de discordance T1+/T2-.}
\end{popfigure}

\courssub{Démarche et corrigé commenté}

\begin{exoresolu}[Le parcours de soins de Marie et Paul : corrigé commenté]
Énoncé : répondre aux six consignes à partir du dossier de Bafoussam.
\tcblower
\textbf{1. Interprétation du dépistage de Marie.}

L'algorithme national repose sur l'enchaînement de tests complémentaires : le test 1, très sensible, ne laisse passer presque aucun vrai positif (peu de faux négatifs) ; le test 2, très spécifique, élimine les faux positifs (peu de faux positifs). Ici, T1+ \emph{et} T2+ : la séropositivité de Marie est \textbf{confirmée} sans ambiguïté. Le test 3 n'est exigé qu'en cas de discordance (T1+/T2$-$) ; il n'est donc pas nécessaire dans ce cas. La \cle{fenêtre silencieuse} (période de séroconversion, environ 3 semaines à 3 mois post-infection) correspond à une infection récente où les anticorps ne sont pas encore détectables : les tests restent négatifs. Or les tests de Marie sont \emph{positifs}, donc la question de la fenêtre ne se pose pas pour elle — son infection date d'au moins plusieurs semaines, probablement de la phase asymptomatique (elle est en bon état clinique, stade OMS 1).

\textbf{2. Risque de transmission mère-enfant.}

Sans intervention, le risque est de 25 \% à 35 \% : sur 100 naissances, on attendrait \textbf{25 à 35 enfants infectés}. Avec l'option B+ (ARV de la mère dès maintenant, accouchement encadré, prophylaxie du nouveau-né avec Névirapine ou Zidovudine selon protocole, allaitement maternel exclusif sous couverture ARV selon recommandations nationales), le risque tombe à moins de 2 \% : \textbf{moins de 2 enfants infectés sur 100}. Le calcul est éloquent :
\[ \text{Enfants protégés sur 100 naissances} \approx 30 - 2 = 28. \]
La PTME évite environ 28 infections sur 100 : c'est l'un des arguments les plus forts en faveur du dépistage prénatal systématique et de la mise sous traitement immédiate.

\textbf{3. Interprétation des marqueurs biologiques.}

CD4 = \SI{450}{cellules/mm^3} : Marie se situe dans la tranche 350--499, soit une \textbf{immunodépression légère} (stade 2 OMS) ; cliniquement asymptomatique (stade OMS 1). Sa charge virale de \SI{85 000}{copies/mL} est \textbf{élevée}, signe d'une réplication virale active, ce qui augmente à la fois le risque de progression de la maladie (baisse future des CD4) et le risque de transmission (verticale vers l'enfant et horizontale vers le partenaire). Selon la stratégie « tester et traiter » (adoptée par le Cameroun depuis 2017) et l'option B+, la mise sous ARV est \textbf{immédiate et urgente} : elle protège la mère (restauration immunitaire), l'enfant à naître (prévention TME) et, à terme, Paul (TasP).

\textbf{4. Choix de la première ligne : TLD.}

On retient le \textbf{TLD} (Ténofovir + Lamivudine + Dolutégravir). Justifications :
\begin{itemize}
\item Le dolutégravir (inhibiteur d'intégrase, INI) fait chuter la charge virale \emph{plus rapidement} que l'éfavirenz (INNTI) — un atout majeur ici avec une CV de \SI{85 000}{copies/mL} et une grossesse en cours, où il faut atteindre l'indétectabilité ($<$ 50 copies/mL) avant l'accouchement pour minimiser le risque résiduel de TME.
\item Le TLD présente une meilleure tolérance globale (moins d'effets neuropsychiques : insomnies, cauchemars, dépression fréquents avec l'éfavirenz) et une barrière génétique à la résistance plus élevée (sélection de mutants plus difficile).
\item Le TLE (TDF+3TC+EFV) reste une alternative valable (notamment si contre-indication au DTG, rare), mais moins adaptée à ce profil grossesse + CV élevée + urgence d'indétectabilité.
\end{itemize}

\textbf{5. Conseils à Paul.}

Un test négatif ne vaut que pour la période hors fenêtre silencieuse : comme le dernier rapport non protégé est récent (vie de couple active), Paul doit \textbf{refaire un test à 6 semaines (J45) puis à 3 mois (J90)} après la dernière exposition à risque pour lever le doute d'une infection en phase de fenêtre. En attendant, le couple doit utiliser systématiquement le \textbf{préservatif} (masculin ou féminin) à chaque rapport. Deux stratégies biomédicales complémentaires s'offrent à lui :
\begin{itemize}
\item La \textbf{PrEP} (prophylaxie pré-exposition) : prise quotidienne d'un comprimé ARV (généralement TDF+3TC) par le partenaire séronégatif, très efficace ($>$ 95 \% de protection si observance parfaite). Elle peut être initiée dès maintenant au centre de santé.
\item À moyen terme, le \textbf{TasP} (« Treatment as Prevention ») : quand Marie, sous TLD bien observé, aura une charge virale indétectable durable ($<$ 50 copies/mL pendant au moins 6 mois), le risque de transmission sexuelle deviendra quasi nul (concept « I = I » : Indétectable = Intransmissible), validé par les études PARTNER, HPTN 052, Opposites Attract.
\end{itemize}
Paul n'est donc pas « à l'abri » aujourd'hui, mais le couple dispose d'outils efficaces pour vivre sereinement et concevoir cet enfant en sécurité.

\textbf{6. Dimensions psychosociales.}

Deux risques majeurs :
\begin{enumerate}
\item La \textbf{stigmatisation et le rejet} (belle-famille, voisinage, milieu professionnel), pouvant conduire Marie au secret, à l'isolement et à l'interruption de son traitement (perte de vue).
\item La \textbf{rupture d'observance} liée à la détresse psychologique (choc du diagnostic, anxiété pour le bébé), aux tensions conjugales (accusations, violence) ou aux effets indésirables initiaux.
\end{enumerate}
Deux actions concrètes :
\begin{enumerate}
\item Orienter le couple vers un \textbf{groupe de soutien de PVVIH} (GSP) ou une association communautaire (ex. : RECAJ+, CAMNAFAW, présentes à Bafoussam) et vers l'\textbf{assistante sociale} de l'hôpital pour un accompagnement psychosocial continu (préparation à l'annonce, gestion du stress, droits sociaux).
\item Rappeler et faire appliquer le \textbf{cadre juridique camerounais} (Loi n° 2021/014 du 9 juillet 2021 relative à la prévention, la prise en charge et le contrôle du VIH/Sida) : le statut sérologique est couvert par le secret médical (art. 14), l'annonce à la famille n'appartient qu'à Marie (droit à la confidentialité), toute discrimination est prohibée (art. 16). Un accompagnement psychologique de couple (conseil pré et post-test conjoint) renforcera l'adhésion au traitement et la prévention au sein du couple.
\end{enumerate}
\end{exoresolu}

\courssub{Bilan de l'activité}

Cette étude de cas a montré que la lutte contre le VIH/Sida est une \emph{chaîne cohérente} où chaque maillon compte : le dépistage précoce (ici, grâce à la consultation prénatale) déclenche la confirmation biologique, l'évaluation immuno-virologique (CD4, charge virale), la mise sous traitement immédiat (TLD), la prévention de la transmission mère-enfant (option B+, qui divise le risque par plus de 15) et la protection du partenaire (préservatif, PrEP, TasP). Tu as également constaté que la réussite de la prise en charge ne dépend pas seulement des molécules : l'observance, le soutien psychosocial, le respect de la confidentialité et des droits des PVVIH sont des conditions tout aussi décisives. Retiens le message central : \textbf{dépister tôt, traiter tôt, accompagner toujours} — c'est ainsi que le Cameroun progresse vers les objectifs 95-95-95 (95 \% des PVVIH connaissent leur statut, 95 \% des diagnostiqués sont sous ARV, 95 \% des traités ont une CV indétectable) et vers une génération sans sida.

\newpage

\coursec{Méthodes et exercices résolus}

\coursec{Lutte contre le VIH/Sida}

\courssub{Le VIH : structure et cycle infectieux}

\begin{definition}[Virus de l'Immunodéficience Humaine (VIH)]
Le \cle{VIH} est un \cle{rétrovirus} (famille des \emph{Retroviridae}, genre \emph{Lentivirus}) à enveloppe, dont le génome est constitué de deux brins d'\cle{ARN} simple brin de polarité positive. Deux types existent : le \cle{VIH-1} (pandémique, plus virulent) et le \cle{VIH-2} (endémique en Afrique de l'Ouest, moins transmissible). Au Cameroun, le VIH-1 (groupes M, N, O) circule activement.
\end{definition}

\begin{popfigure}
\begin{tikzpicture}[scale=0.9, every node/.style={font=\small}]
% Envelope
\draw[popblue, thick, fill=popblue!10] (0,0) ellipse (2.5cm and 1.5cm);
% Glycoproteins
\foreach \angle/\r in {30/1.5, 60/1.5, 90/1.5, 120/1.5, 150/1.5, 210/1.5, 240/1.5, 270/1.5, 300/1.5, 330/1.5} {
    \draw[poporange, thick] (\angle:\r) -- (\angle:\r+0.4) node[above right, poporange] {\tiny gp120/gp41};
}
% Matrix (p17) - inner side of envelope
\draw[popdark, dashed] (0,0) ellipse (2.3cm and 1.3cm);
\node[popdark] at (2.5, -1.2) {\tiny Matrice (p17)};
% Capsid (p24)
\draw[poppurple, thick, fill=poppurple!10, rotate=15] (0,0) ellipse (1.2cm and 0.7cm);
\node[poppurple] at (0,0.2) {\tiny Capside (p24)};
% Viral core content
\draw[popgreen, thick] (-0.6,0) -- (-0.6, -0.3) node[below] {\tiny ARN viral (2 brins)};
\draw[popgreen, thick] (0.6,0) -- (0.6, -0.3) node[below] {\tiny ARN viral (2 brins)};
% Enzymes inside capsid
\node[popred] at (0, -0.6) {\tiny Transcriptase inverse (p66/p51)};
\node[popred] at (0, -0.85) {\tiny Intégrase (p32)};
\node[popred] at (0, -1.1) {\tiny Protéase (p10)};
% Legend anchors
\node[anchor=west, align=left] at (3.2, 1.5) {\textbf{Légende :}};
\node[anchor=west, align=left] at (3.2, 1.1) {\textcolor{popblue}{■} Enveloppe lipidique (origine cellulaire)};
\node[anchor=west, align=left] at (3.2, 0.7) {\textcolor{poporange}{■} Glycoprotéines \textbf{gp120} (fixation) \& gp41 (fusion)};
\node[anchor=west, align=left] at (3.2, 0.3) {\textcolor{popdark}{■} Protéine matrice \textbf{p17}};
\node[anchor=west, align=left] at (3.2, -0.1) {\textcolor{poppurple}{■} Capside protéique \textbf{p24} (antigène majeur)};
\node[anchor=west, align=left] at (3.2, -0.5) {\textcolor{popgreen}{■} Génome : 2 ARN (+) + \textbf{tRNA Lys} (amorce)};
\node[anchor=west, align=left] at (3.2, -0.9) {\textcolor{popred}{■} Enzymes virales embarquées};
\end{tikzpicture}
\legende{Structure du virion VIH-1 (coupe schématique). La gp120 assure la reconnaissance du récepteur CD4 ; la gp41 médie la fusion. La capside (p24) protège l'ARN et les enzymes. L'antigène p24 est dosé dans le sang pour le diagnostic précoce (test combiné Ag/Ab).}
\end{popfigure}

\begin{methode}[Décrire le VIH et ordonner les étapes du cycle infectieux]
\begin{enumerate}
\item \textbf{Décrire la structure} de l'extérieur vers l'intérieur : \cle{enveloppe} lipidique (origine membrane cellulaire hôte) portant les glycoprotéines \cle{gp120} (fixation) et \cle{gp41} (fusion) $\rightarrow$ protéine \cle{matrice (p17)} $\rightarrow$ \cle{capside} protéique (protéine \cle{p24}, antigène majeur dosé en routine) $\rightarrow$ deux molécules d'\cle{ARN} génomique (chacune liée à un \cle{tRNA Lys} servant d'amorce) et les trois enzymes virales : \cle{transcriptase inverse} (p66/p51), \cle{intégrase} (p32), \cle{protéase} (p10).
\item \textbf{Ordonner les étapes du cycle} dans le lymphocyte T CD4 (cible principale, mais aussi macrophages, cellules dendritiques) :
\begin{enumerate}
\item \textbf{Fixation (Adsorption)} : la \cle{gp120} se lie au récepteur \cle{CD4} puis à un \cle{corécepteur} (CCR5 ou CXCR4) $\rightarrow$ changement conformationnel de gp41.
\item \textbf{Pénétration (Fusion/Entrée)} : \cle{gp41} médie la fusion de l'enveloppe virale avec la membrane plasmique $\rightarrow$ libération du noyau capsidique dans le cytoplasme.
\item \textbf{Transcription inverse} : la \cle{transcriptase inverse} synthétise un ADN simple brin (ADNs) sur le matrice d'ARN, dégrade l'ARN (activité RNase H), puis synthétise le second brin $\rightarrow$ formation de l'\cle{ADN viral double brin} (pré-intégration). \emph{Erreur fréquente : confusion avec la transcription classique ADN$\rightarrow$ARN.}
\item \textbf{Intégration} : l'\cle{intégrase} insère l'ADN viral (appelé \cle{provirus}) dans le génome de la cellule hôte (chromosome). Le provirus est alors répliqué à chaque division cellulaire.
\item \textbf{Expression (Transcription/Traduction)} : activation cellulaire $\rightarrow$ transcription du provirus en ARN messagers (ARNm) et ARN génomiques $\rightarrow$ traduction en polyprotéines virales (Gag, Pol, Env) au niveau du réticulum endoplasmique/Golgi.
\item \textbf{Assemblage, Maturation et Bourgeonnement} : rassemblement des composants au niveau de la membrane plasmique $\rightarrow$ bourgeonnement (acquisition de l'enveloppe) $\rightarrow$ clivage des polyprotéines par la \cle{protéase virale} (maturation) $\rightarrow$ virions infectieux libérés.
\end{enumerate}
\item \textbf{Relier structure et fonction} : gp120/CD4 $\rightarrow$ tropisme ; Transcriptase inverse $\rightarrow$ cible des INRTI/INNTI ; Intégrase $\rightarrow$ cible des II ; Protéase $\rightarrow$ cible des IP. L'intégration explique la persistance (réservoirs) et l'incurabilité actuelle.
\item \textbf{Au brouillon} : vérifier l'ordre chronologique, le vocabulaire exact (transcription \emph{inverse} $\neq$ transcription), les noms d'enzymes.
\end{enumerate}
\end{methode}

\begin{experience}[Modélisation du cycle infectieux (démonstration guidée)]
\textbf{Objectif :} Visualiser l'ordre des étapes et les cibles thérapeutiques.
\textbf{Matériel :} Schémas découpés représentant les 6 étapes (fixation, fusion, transcription inverse, intégration, expression, bourgeonnement), noms des enzymes (transcriptase inverse, intégrase, protéase), noms des classes d'ARV (INRTI, INNTI, II, IP, IF).
\textbf{Protocole :}
\begin{enumerate}
\item Placer les 6 étapes dans l'ordre chronologique sur une frise.
\item Associer chaque enzyme à l'étape qu'elle catalyse.
\item Placer chaque classe d'ARV sur l'étape qu'elle bloque.
\end{enumerate}
\textbf{Observations/Interprétation attendues :}
\begin{itemize}
\item Fixation $\leftarrow$ Inhibiteurs de fixation (ex: Maraviroc - anti-CCR5, Enfuvirtide - fusion).
\item Transcription inverse $\leftarrow$ INRTI (ex: Zidovudine, Lamivudine, Tenofovir) + INNTI (ex: Névirapine, Éfavirenz, Dolutégravir \emph{attention : Dolutégravir est un II, corriger : Relpivirine, Doravirine}).
\item Intégration $\leftarrow$ II (ex: Raltégravir, Dolutégravir, Bictégravir).
\item Maturation (Protéase) $\leftarrow$ IP (ex: Lopinavir/r, Atazanavir/r, Darunavir/r).
\end{itemize}
\textbf{Conclusion :} La trithérapie standard (ex: TDF+3TC+DTG au Cameroun) cible au moins deux étapes distinctes (souvent Transcription inverse + Intégration) pour prévenir la résistance.
\end{experience}

\begin{exoresolu}[Le cycle du VIH dans le désordre]
Voici, dans le désordre, six étapes du cycle infectieux du VIH dans un lymphocyte T CD4 :
(a) intégration de l'ADN viral dans le chromosome de la cellule hôte ; (b) fixation de la gp120 sur le récepteur CD4 ; (c) assemblage et libération de nouveaux virions ; (d) synthèse d'ADN viral à partir de l'ARN viral par la transcriptase inverse ; (e) production de protéines virales à partir du provirus ; (f) fusion de l'enveloppe virale avec la membrane et entrée du contenu viral.
\begin{enumerate}
\item Remets ces étapes dans l'ordre chronologique.
\item Nomme l'enzyme virale responsable de l'étape (d) et précise sa particularité.
\item Explique pourquoi l'étape (a) rend l'infection difficile à éradiquer.
\end{enumerate}
\tcblower
\textbf{1. Ordre chronologique.}
(b) fixation de la gp120 sur le récepteur CD4 (et corécepteur CCR5/CXCR4) $\rightarrow$ (f) fusion de l'enveloppe et entrée du contenu viral (décapsidation partielle) $\rightarrow$ (d) transcription inverse (ARN $\xrightarrow{\text{Transcriptase inverse}}$ ADN double brin) $\rightarrow$ (a) intégration du provirus dans le génome hôte (Intégrase) $\rightarrow$ (e) expression des gènes viraux (transcription/traduction cellulaire) $\rightarrow$ (c) assemblage, maturation (Protéase) et libération des virions (bourgeonnement).

\textbf{2. L'enzyme de l'étape (d).}
Il s'agit de la \textbf{transcriptase inverse} (reverse transcriptase, p66/p51). Sa particularité : elle réalise la synthèse d'ADN à partir d'une matrice d'ARN (\textbf{flux inverse} : ARN $\rightarrow$ ADN), contredisant le dogme central classique (ADN $\rightarrow$ ARN $\rightarrow$ Protéine). Elle ne possède pas d'activité de relecture (proofreading) $\rightarrow$ taux de mutation élevé ($\sim 3 \times 10^{-5}$/base/cycle) $\rightarrow$ diversité génétique et résistance aux ARV. C'est la cible des INRTI et INNTI.

\textbf{3. Conséquence de l'intégration (a).}
Une fois l'ADN viral intégré sous forme de \textbf{provirus}, il devient partie intégrante du génome de la cellule hôte. Il peut rester \textbf{latent (silencieux)} pendant des années dans des \cle{réservoirs viraux} (lymphocytes T CD4 mémoire au repos, macrophages, cellules du SNC). Les ARV bloquent la production de \emph{nouveaux} virus mais n'éliminent pas le provirus intégré ni les cellules réservoirs. C'est pourquoi l'infection est \textbf{incurable} à l'heure actuelle et que le traitement doit être pris \textbf{à vie} (observance stricte).

\textbf{Conclusion} : la connaissance précise du cycle (fixation, entrée, transcription inverse, intégration, expression, bourgeonnement) permet de comprendre le mode d'action des 6 classes d'antirétroviraux et la persistance de l'infection malgré la suppression virale.
\end{exoresolu}

\courssub{Modes de transmission et prévention}

\begin{propriete}[Liquides biologiques infectants et voies d'entrée]
Le VIH est présent en quantité \textbf{infectante} uniquement dans : le \cle{sang} (et dérivés), le \cle{sperme}, les \cle{secrétions vaginales}, le \cle{lait maternel}. Les autres liquides (salive, larmes, sueur, urine, selles, vomissures) ne contiennent pas ou trop peu de virus pour transmettre l'infection, \textbf{sauf s'ils sont visiblement souillés par du sang}.
La transmission nécessite un \textbf{contact direct} entre un liquide infectant et une \cle{voie d'entrée} : muqueuses (génitales, anales, buccales si lésées), peau lésée (plaie, piqûre), injection intraveineuse/sous-cutanée, passage placentaire.
\end{propriete}

\begin{methode}[Analyser une situation à risque et proposer une prévention adaptée]
\begin{enumerate}
\item \textbf{Identifier le liquide biologique} mis en jeu dans la situation.
\item \textbf{Vérifier la condition de transmission} : contact direct liquide infectant $\rightarrow$ voie d'entrée (muqueuse, plaie, injection). Sans effraction cutanée ni contact muqueux $\rightarrow$ risque nul.
\item \textbf{Conclure sur le risque} : réel (rapport sexuel non protégé, partage seringue, transfusion non testée, transmission mère-enfant) ou nul (poignée de main, baiser social, piqûre moustique, repas partagé, piscine, toilettes).
\item \textbf{Proposer la prévention adaptée} (stratégie \textbf{combinée}) :
\begin{itemize}
\item \textbf{Via sexuelle} (80\% des cas au Cameroun) : \cle{Abstinence} (100\% efficace), \cle{Fidélité} réciproque entre partenaires \textbf{dépistés}, \cle{Préservatif} (masculin ou féminin) \textbf{correctement utilisé} à chaque rapport. \cle{PrEP} (Prophylaxie Pré-Exposition, ex: TDF/3TC) pour personnes à risque élevé. \cle{PEP} (Prophylaxie Post-Exposition, traitement d'urgence < 72h, idéalement < 4h) après rupture de préservatif ou viol.
\item \textbf{Via sanguine} : Seringues/matériel \textbf{à usage unique} (stérile), programmes d'échange de seringues (RDR - Réduction des Risques) pour usagers de drogues, \textbf{dépistage systématique} de tout donneur de sang (sérologie VIH Ag/Ab + Hépatites B/C + Syphilis), stérilisation matériel médical/chirurgical/dentaire/coiffure/tatouage.
\item \textbf{Via mère-enfant (PTME - Prévention Transmission Mère-Enfant)} : Dépistage systématique 1ère consultation prénatale (CPN) ; si séropositive $\rightarrow$ \textbf{ARV à vie (Option B+)} dès le diagnostic (TDF+3TC+DTG) ; accouchement encadré (césarienne si charge virale $> 1000$ copies/mL) ; \textbf{prophylaxie néonatale} (NVP ou AZT selon protocole national) 6 semaines ; \textbf{allaitement maternel exclusif} 6 mois sous couvert ARV maternel (charge virale indétectable) puis sevrage rapide, \textbf{OU} lait de substitution \textbf{SI conditions AFASS} (Acceptable, Faisable, Abordable, Durable, Sûr) réunies. Dépistage enfant à 6 semaines (PCR ADN), 9 mois, fin allaitement.
\end{itemize}
\item \textbf{Rédiger} en nommant précisément le mode de transmission (sexuel, sanguin, vertical) et le moyen de prévention.
\end{enumerate}
\end{methode}

\begin{exoresolu}[Risque réel ou idée reçue ? Contexte Douala]
Situations rapportées :
\begin{enumerate}
\item Aïcha partage chaque jour son repas avec sa voisine séropositive.
\item Bertin a eu un rapport sexuel non protégé avec une partenaire dont il ignore le statut sérologique.
\item Carine craint d'avoir été contaminée par les piqûres de moustiques.
\item Un toxicomane utilise la seringue d'un ami séropositif pour s'injecter une drogue.
\end{enumerate}
Pour chaque situation : risque réel/nul + justification + prévention si risque.
\tcblower
\textbf{Méthode} : Identifier liquide + voie d'entrée.

\textbf{1. Repas partagé (Aïcha)} : Liquide = salive (non infectante, sauf sang visible). Voie d'entrée = muqueuse buccale \textbf{intacte}. \textbf{Risque nul.} Prévention : aucune nécessaire. Lutter contre la stigmatisation (idées reçues).

\textbf{2. Rapport sexuel non protégé (Bertin)} : Liquides = sperme / sécrétions vaginales (infectants). Voie d'entrée = muqueuses génitales (pénétration) ou anales. \textbf{Risque réel (élevé).} Prévention : Préservatif correct/systématique ; Dépistage couple avant rapports non protégés ; PrEP si risque récurrent ; PEP d'urgence si < 72h.

\textbf{3. Piqûres de moustiques (Carine)} : Moustique n'injecte \textbf{pas} le sang du repas précédent (injecte sa salive anticoagulante). VIH ne se réplique \textbf{pas} dans le moustique (pas de récepteur CD4, digestion virale). \textbf{Risque nul.} (Rappel : moustique = vecteur \emph{Plasmodium} paludisme, arbovirus dengue/chikungunya).

\textbf{4. Partage de seringue} : Liquide = sang résiduel dans seringue/aiguille (très infectant). Voie d'entrée = injection directe dans circulation sanguine (voie parentérale). \textbf{Risque réel (très élevé, efficacité transmission $\sim$ 90\%).} Prévention : Matériel stérile à usage unique ; Programmes RDR (échange seringues) ; Traitement substitution (méthadone) ; PEP si accident exposition (AES) personnel soignant.

\textbf{Conclusion} : Seules situations 2 et 4 = risque réel. Situations 1 et 3 = idées reçues nuisibles (stigmatisation, fausse sécurité).
\end{exoresolu}

\courssub{Évolution naturelle de l'infection et marqueurs biologiques}

\begin{popfigure}
\begin{tikzpicture}
\begin{axis}[
    width=\linewidth, height=8cm,
    xlabel={Temps après infection (années)}, ylabel={Valeurs},
    xmin=0, xmax=12, ymin=0, ymax=1200,
    axis y line*=left,
    legend style={at={(0.5,1.02)}, anchor=south, legend columns=2, font=\footnotesize},
    grid=major, grid style={popline},
]
% Viral Load (log scale simulated on linear axis for simplicity, or use two axes)
% We'll plot log10(VL) scaled to fit, or use two y axes. Let's use two y axes for rigor.
\addplot[popcurve, domain=0:12, samples=200, thick]
    ({x}, {700*exp(-0.8*x) + 100 + 300*exp(0.15*(x-6))}); % Schematic shape
\addlegendentry{Charge virale (échelle log arbitraire)}

% CD4 count
\addplot[popgreen, thick, domain=0:12, samples=200] 
    ({x}, {800 - 50*x + 100*exp(-0.5*x) - 200*exp(0.1*(x-8))});
\addlegendentry{CD4 (cellules/mm$^3$)}

% Phases annotations
\node[popblue, anchor=south] at (axis cs:0.5, 900) {Phase primaire};
\node[poporange, anchor=south] at (axis cs:4, 900) {Phase asymptomatique};
\node[popred, anchor=south] at (axis cs:10, 900) {Phase Sida};

% Threshold lines
\draw[dashed, popred, thick] (axis cs:0, 200) -- (axis cs:12, 200) node[right] {Seuil Sida (200 CD4)};
\draw[dashed, popmuted] (axis cs:0, 500) -- (axis cs:12, 500) node[right] {Seuil normal bas (500 CD4)};
\end{axis}
\end{tikzpicture}
\legende{Évolution schématique de la charge virale (courbe bleue, échelle logarithmique) et des lymphocytes T CD4 (courbe verte, échelle linéaire) au cours de l'infection à VIH non traitée. Trois phases : 1) Primaire (pic viral, séroconversion), 2) Asymptomatique (équilibre instable, CD4 décroissent), 3) Sida (CD4 < 200/mm$^3$, reprise virale, infections opportunistes).}
\end{popfigure}

\begin{definition}[Phases de l'infection à VIH non traitée]
\begin{enumerate}
\item \textbf{Phase primaire (Primo-infection)} : 0 à 3 mois. Réplication virale massive $\rightarrow$ \cle{charge virale} très élevée ($> 10^6$ copies/mL). Syndrome pseudo-grippal/mononucléosique (fièvre, adénopathies, éruption, méningite aseptique) dans 50-70\% cas. Chute transitoire CD4. \textbf{Séroconversion} : apparition \cle{anticorps anti-VIH} (délai médian 3-4 semaines, max 3 mois = \cle{fenêtre sérologique}). Contagiosité maximale.
\item \textbf{Phase asymptomatique (Latence clinique)} : 2 à 10 ans (médiane 8-10 ans). Charge virale stabilisée (\cle{point de consigne viral}, ex: $10^4 - 10^5$ copies/mL). Déclin lent et continu des CD4 ($\sim$ 50-80 cellules/mm$^3$/an). Patient \textbf{séropositif, asymptomatique, mais contagieux}. Adénopathies généralisées persistantes (AGP) fréquentes.
\item \textbf{Phase Sida (Stade ultime)} : CD4 $< \textbf{200 cellules/mm}^3$ \textbf{ET/OU} survenue d'\cle{infections opportunistes} (IO) ou cancers définissant le Sida (stade C3 CDC). Charge virale remonte. IO bactériennes (tuberculose \emph{---} 1ère cause décès au Cameroun), fongiques (pneumocystose, cryptococcose, candidoses œsophagienne), virales (CMV, herpès zona disséminé), parasitaires (toxoplasmose cérébrale, cryptosporidiose). Cancers : sarcome de Kaposi, lymphomes, cancer col utérin.
\end{enumerate}
\end{definition}

\begin{methode}[Lire et interpréter un graphe charge virale / lymphocytes CD4]
\begin{enumerate}
\item \textbf{Identifier axes et unités} : Abscisse = temps (mois/années) ; Ordonnées = Charge virale (copies ARN/mL, échelle souvent log) et CD4 (cellules/mm$^3$). Normal CD4 : 500--1500/mm$^3$.
\item \textbf{Repérer les trois phases} (voir définition ci-dessus) : Pic viral + séroconversion $\rightarrow$ Palier viral + déclin CD4 $\rightarrow$ Rupture immunitaire (CD4 < 200) + reprise virale.
\item \textbf{Comparer les courbes} : Relation \textbf{inverse} : réplication virale $\rightarrow$ destruction CD4 (lyse cellulaire, apoptose, syncytia, élimination par CD8). Cercle vicieux.
\item \textbf{Exploiter valeurs numériques} : Citer valeurs lues (avec unités) pour chaque affirmation.
\item \textbf{Conclure} : Relier à définition Sida (CD4 < 200 + IO) et intérêt ARV (chute VL, remontée CD4, restauration immunité).
\end{enumerate}
\end{methode}

\begin{exoresolu}[Suivi biologique d'un patient non traité (10 ans)]
Données :
\begin{center}
\begin{tabularx}{\linewidth}{|l|X|X|X|X|X|}
\hline
\rowcolor{popblueL} Temps après infection & 1 mois & 6 mois & 3 ans & 7 ans & 10 ans \\
\hline
Charge virale (copies/mL) & \num{1000000} & \num{50000} & \num{60000} & \num{200000} & \num{800000} \\
\hline
CD4 (cellules/mm\textsuperscript{3}) & 400 & 700 & 550 & 300 & 120 \\
\hline
\end{tabularx}
\end{center}
\begin{enumerate}
\item Décrire l'évolution de la charge virale et du taux de CD4.
\item Identifier la phase à 10 ans (justifier).
\item Expliquer la relation entre les deux grandeurs.
\end{enumerate}
\tcblower
\textbf{1. Description.}
\textbf{Charge virale} : Pic à 1 mois (\num{1000000} copies/mL, primo-infection) $\rightarrow$ Chute à 6 mois (\num{50000}, réponse immunitaire innée/adaptative) $\rightarrow$ Palier "point de consigne" 3 ans (\num{60000}) $\rightarrow$ Remontée progressive 7 ans (\num{200000}) et 10 ans (\num{800000}, perte contrôle immunitaire).
\textbf{CD4} : Baisse initiale 1 mois (400) $\rightarrow$ Remontée partielle 6 mois (700, contrôle transitoire) $\rightarrow$ Déclin lent et régulier : 3 ans (550), 7 ans (300) $\rightarrow$ Effondrement 10 ans (\textbf{120 cellules/mm$^3$}).

\textbf{2. Phase à 10 ans.}
CD4 = 120 $< \textbf{200 cellules/mm$^3$}$. Seuil franchit $\rightarrow$ \textbf{Immunodéficience profonde} $\rightarrow$ \textbf{Stade Sida déclaré (C3)}. Risque immédiat d'infections opportunistes (TB, pneumocystose, méningite cryptococcique...) et cancers. Urgence thérapeutique ARV + prophylaxies IO (Cotrimoxazole, Fluconazole).

\textbf{3. Relation.}
Évolution \textbf{inverse} (corrélation négative). Mécanisme : VIH infecte spécifiquement lymphocytes T CD4 (gp120 $\rightarrow$ CD4/CCR5) $\rightarrow$ réplication $\rightarrow$ destruction cellulaire (lyse, apoptose, mort par syncytia, élimination CD8 cytotoxiques). Plus VL haute $\rightarrow$ plus destruction CD4 $\rightarrow$ moins contrôle immunitaire $\rightarrow$ plus réplication virale. Cercle vicieux.

\textbf{Conclusion} : Évolution naturelle vers Sida en $\sim$10 ans. ARV précoces (dès diagnostic, quelque soit CD4) brisent ce cercle : VL indétectable + CD4 $> 500$ = vie quasi normale.
\end{exoresolu}

\courssub{Dépistage : principes, stratégies et algorithmes}

\begin{definition}[Dépistage du VIH]
Recherche des \cle{anticorps anti-VIH} (Ac anti-gp120, gp41, p24) dans le sang total, sérum ou plasma. Au Cameroun : tests rapides (TR) de 4ème génération (détectent \textbf{Ag p24 + Ac}) ou ELISA (labo). \textbf{Principes éthiques (3C)} : \cle{Consentement} libre/éclairé, \cle{Confidentialité} absolue, \cle{Conseil} (pré-test et post-test). Gratuit dans les structures publiques et agréées (CTA, HDJ, ONG).
\end{definition}

\begin{propriete}[Algorithme national de dépistage (Cameroun, OMS)]
\begin{enumerate}
\item \textbf{Test de dépistage (T1)} : Test rapide sensible (ex: Determine, SD Bioline, First Response). 
\begin{itemize}
\item \textbf{Négatif} : Si \textbf{hors fenêtre sérologique} (dernier risque $> 3$ mois) $\rightarrow$ \textbf{Séronégatif}. Si risque récent ($< 6$ semaines) $\rightarrow$ \textbf{Refaire test à 6 semaines puis 3 mois}.
\item \textbf{Positif (Réactif)} : Ne \textbf{jamais} annoncer "séropositif" sur un seul test. $\rightarrow$ \textbf{Test de confirmation (T2)} obligatoire (autre marqueur, autre kit, ex: Stat-Pak, Uni-Gold).
\end{itemize}
\item \textbf{Test de confirmation (T2)} :
\begin{itemize}
\item \textbf{Positif} $\rightarrow$ \textbf{Séropositivité confirmée} $\rightarrow$ Annonce résultat + orientation prise en charge (bilan CD4/VL, TB, ITS, grossesse, ARV).
\item \textbf{Négatif} (discordance T1+/T2-) $\rightarrow$ \textbf{Test de départage (T3)} (ELISA labo ou 3ème TR) ou contrôle à 2-4 semaines.
\end{itemize}
\end{enumerate}
\textbf{Fenêtre sérologique} : Période post-contamination où Ag p24 et/ou Ac ne sont pas encore détectables. Durée : \textbf{2 à 8 semaines} (médiane 3-4 sem) pour tests 4ème géné (Ag/Ab). \textbf{Prudence :} test négatif $< 6$ sem post-risque = non concluant. Contrôle impératif à \textbf{3 mois}.
\end{propriete}

\begin{methode}[Raisonner sur un résultat de test et la fenêtre sérologique]
\begin{enumerate}
\item Identifier le type de test (Ag/Ab 4ème géné vs Ab seul 3ème géné).
\item Situer la date du dernier risque par rapport au test.
\item Appliquer l'algorithme : Négatif + hors fenêtre = rassurant. Négatif + dans fenêtre = refaire test (J+45, J+90). Positif = Confirmation obligatoire (T2) avant annonce.
\item Ne jamais confondre "négatif" (pas d'infection détectée) et "séronégatif certifié" (hors fenêtre).
\end{enumerate}
\end{methode}

\begin{exoresolu}[Interprétation algorithme dépistage]
Un jeune se présente au CTA 3 semaines après un rapport non protégé. Test rapide (4ème géné Ag/Ab) = Négatif.
\begin{enumerate}
\item Ce résultat exclut-il l'infection ? Justifier.
\item Quelle conduite tenir ?
\end{enumerate}
\tcblower
\textbf{1. Non.} 3 semaines $<$ délai médian séroconversion (3-4 sem) et $<$ fenêtre max (8 sem pour Ag, 12 sem pour Ab seuls). Le test peut être \textbf{faux négatif} (infection en cours, pas encore de marqueurs détectables). Le risque n'est pas écarté.

\textbf{2. Conduite :}
\begin{itemize}
\item Expliquer la fenêtre sérologique (pédagogie, réduire anxiété).
\item \textbf{Refaire test à 6 semaines (J+45)} : détection Ag p24 + Ac précoces.
\item \textbf{Refaire test à 3 mois (J+90)} : certitude quasi absolue (fin fenêtre pour tous tests).
\item Conseiller prévention (préservatif, PrEP si risque récurrent) pendant cette période.
\item Si symptômes primo-infection (fièvre, adénopathies, éruption) $\rightarrow$ dosage \textbf{charge virale (PCR ARN)} ou \textbf{Ag p24 quantitatif} (diagnostic précoce direct, avant anticorps).
\end{itemize}
\end{exoresolu}

\courssub{Prise en charge : Traitements Antirétroviraux (ARV) et suivi}

\begin{definition}[Traitement Antirétroviral (ARV) / Trithérapie]
Association d'au moins \textbf{3 molécules} actives (généralement 2 INRTI + 1 INNTI \textbf{ou} 1 II \textbf{ou} 1 IP boosté) pour \textbf{bloquer la réplication virale} à plusieurs étapes, prévenir la résistance et obtenir une \textbf{suppression virale durable}. Au Cameroun (PNLS 2023/Guide OMS) : \textbf{Ligne 1 standard} = \textbf{TDF + 3TC (ou FTC) + DTG} (Dolutégravir, Inhibiteur Intégrase, haute barrière résistance, bien toléré). Alternatives : TDF+3TC+EFV (grossesse 1er trimestre si DTG contre-indiqué), AZT+3TC+DTG (insuffisance rénale). \textbf{Objectifs :} Charge virale \textbf{indétectable} ($< 50$ copies/mL), CD4 $> 500$/mm$^3$, qualité de vie, prévention transmission (I=I).
\end{definition}

\begin{propriete}[Indétectable = Intransmissible (I = I)]
Une personne séropositive sous ARV efficace, avec \textbf{charge virale indétectable ($< 50$ copies/mL) confirmée depuis au moins 6 mois} et \textbf{observance parfaite}, \textbf{ne transmet pas le VIH par voie sexuelle}. Preuves scientifiques solides (études PARTNER, Opposites Attract, HPTN 052). \textbf{Ne s'applique pas} à la transmission mère-enfant (allaitement) ni au partage de seringues. \textbf{Condition sine qua non :} Observance $\ge 95\%$ (prise quotidienne, heure fixe, pas d'oubli).
\end{propriete}

\begin{methode}[Évaluer l'efficacité d'un traitement ARV à partir du suivi biologique]
\begin{enumerate}
\item \textbf{Indicateurs} : Charge virale (cible : $< 50$ copies/mL = indétectable) ; CD4 (cible : $> 500$ cellules/mm$^3$, augmentation $\ge 50-100$/an).
\item \textbf{Calcul variation relative (\%)} :
\[
\text{Variation (\%)} = \frac{\text{Valeur finale} - \text{Valeur initiale}}{\text{Valeur initiale}} \times 100
\]
\item \textbf{Interprétation} :
\begin{itemize}
\item VL indétectable + CD4 en hausse $\rightarrow$ \textbf{Succès virologique et immunologique} (observance bonne).
\item VL détectable ($> 50$ copies/mL) à 6 mois $\rightarrow$ \textbf{Suspicion échec virologique} $\rightarrow$ Vérifier observance (entretien non jugeant, comptage comprimés, dosage plasmatique ARV si dispo) $\rightarrow$ Résistance (génotypage) $\rightarrow$ Adaptation ligne 2.
\item CD4 stagnants/baisse malgré VL indétectable $\rightarrow$ \textbf{Échec immunologique} (rare, facteurs : âge, co-infections, fibrose ganglionnaire).
\end{itemize}
\item \textbf{Règle d'or} : Ne \textbf{jamais} arrêter ou modifier seul son traitement. Risque résistance majeur.
\end{enumerate}
\end{methode}

\begin{exoresolu}[Suivi d'une patiente sous ARV à Yaoundé (HDJ Essos)]
Mme N., 32 ans, séropositive. Bilan initial (J0) : VL = \num{250000} copies/mL ; CD4 = 280/mm$^3$. Début trithérapie TDF+3TC+DTG (Ligne 1). 6 mois (M6) : VL = \num{40} copies/mL ; CD4 = 520/mm$^3$.
\begin{enumerate}
\item Calculer la variation relative (\%) de la charge virale.
\item Le traitement est-il efficace ? Justifier (2 arguments chiffrés).
\item Mme N. demande d'arrêter car "virus indétectable". Réponse argumentée.
\end{enumerate}
\tcblower
\textbf{1. Variation relative VL.}
\[
\frac{40 - 250000}{250000} \times 100 = \frac{-249960}{250000} \times 100 = \mathbf{-99,98\,\%}
\]
Chute quasi totale de la réplication virale.

\textbf{2. Efficacité : OUI, traitement \textbf{efficace} (Succès virologique + immunologique).}
\begin{itemize}
\item \textbf{Virologique} : VL passée de \num{250000} à \num{40} copies/mL $\rightarrow$ \textbf{Indétectable} ($< 50$ copies/mL). Réplication bloquée.
\item \textbf{Immunologique} : CD4 passés de 280 à 520 cellules/mm$^3$ ($+240$, $+86\%$). Franchissement seuil 500 $\rightarrow$ restauration immunitaire, protection contre IO.
\end{itemize}

\textbf{3. Réponse à Mme N. : \textbf{NON, arrêt strictement interdit.}}
\begin{itemize}
\item "Indétectable" $\neq$ "Guéri". Les ARV bloquent la réplication \textbf{extracellulaire} mais n'éliminent pas le \textbf{provirus intégré} (ADN viral) dans les \textbf{réservoirs} (lymphocytes T CD4 mémoire quiescents, macrophages, SNC, tissu lymphoïde).
\item À l'arrêt : réactivation transcription provirus $\rightarrow$ rebond viral en \textbf{2 à 4 semaines} (VL remonte, CD4 rechutent, résistance possible si mono/bithérapie résiduelle, contagiosité rétablie).
\item \textbf{Traitement à vie}, observance quotidienne rigoureuse. Seule une \textbf{rémisson fonctionnelle} (cas exceptionnels "patients de Berlin/Londres/Düsseldorf" post greffe moelle CCR5$\Delta$32) ou future thérapie curative (CRISPR, vaccins thérapeutiques, "shock and kill") pourrait changer la donne.
\end{itemize}
\textbf{Conclusion} : Succès thérapeutique majeur (VL -99,98\%, CD4 x2). Maintien ARV à vie = condition de la survie et de la non-transmission (I=I).
\end{exoresolu}

\courssub{Aspects psychosociaux, droits et riposte nationale}

\begin{savaistu}[Riposte nationale au VIH/Sida au Cameroun]
\textbf{Contexte :} Prévalence 2,7\% (EDS 2018, 15-49 ans), $\sim$ 500 000 PVVIH. Femmes plus touchées (3,9\% vs 1,8\%). Jeunes filles 15-24 ans vulnérables.
\textbf{Stratégie nationale (PNLS, Plan Stratégique National 2020-2024) :} Objectif \textbf{95-95-95} (ONUSIDA 2025) : 95\% PVVIH connaissent statut, 95\% d'entre eux sous ARV, 95\% d'entre eux charge virale supprimée.
\textbf{Piliers :}
\begin{itemize}
\item \textbf{Prévention combinée} : Dépistage massif (CTA, dépistage communautaire, auto-test VIH), Préservatifs (gratuits structures publiques), PrEP (populations clés : HSH, TS, UD, partenaires séronégatifs couples sérodifférents), PEP, PTME (Option B+), Circoncision médicale volontaire.
\item \textbf{Prise en charge universelle gratuite} : ARV gratuits (Fonds Mondial, Budget Etat, Clinton Health Access Initiative), Bilan biologique (CD4, VL) gratuit, Prise en charge IO (Cotrimoxazole, Fluconazole, TB), Soins palliatifs.
\item \textbf{Lutte contre la stigmatisation/discrimination} : Loi n°2005/007 (protection droits PVVIH), Confidentialité médicale, Accès emploi/soins/éducation. Journée Mondiale 1er Décembre, campagnes "Zéro discrimination".
\item \textbf{Approche différenciée} : Modèles de distribution ARV (DMD - Distribution Médicaments à Domicile, Groupes d'adhésion, Points de distribution communautaires) pour décongestionner hôpitaux et améliorer observance.
\end{itemize}
\textbf{Défis persistants :} Perte de vue patients, résistance ARV montante (surveillance DRV), financement durable, prise en charge enfants/adolescents (faible couverture PTME, dépistage pédiatrique), co-infection VIH/TB.
\end{savaistu}

\begin{methode}[Construire un argumentaire pour le dépistage volontaire (Épreuve écrite type)]
\begin{enumerate}
\item \textbf{Accroche} : Confronter l'idée reçue ("Je me sens bien") à la réalité biologique (phase asymptomatique).
\item \textbf{Argument 1 (Biologie)} : Phase asymptomatique = 8-10 ans sans symptômes MAIS virus actif, CD4 baissent, contagiosité réelle. Seul le test (Ac anti-VIH) révèle le statut.
\item \textbf{Argument 2 (Bénéfice individuel)} : Dépistage précoce $\rightarrow$ ARV précoces (ligne 1 TDF/3TC/DTG) $\rightarrow$ VL indétectable, CD4 conservés, espérance de vie normale, éviter Sida.
\item \textbf{Argument 3 (Bénéfice collectif / I=I)} : Sous ARV efficace (VL $< 50$ copies/mL, 6 mois) $\rightarrow$ \textbf{Zéro transmission sexuelle}. Dépistage = protection partenaires + communauté.
\item \textbf{Argument 4 (Accès droits/soins)} : Gratuité tests/ARV/biologie au Cameroun. Confidentialité garantie. Prise en charge globale (médicale, psychologique, sociale, juridique).
\item \textbf{Réponse aux freins} : Peur résultat $\rightarrow$ conseil pré/post-test ; Stigmatisation $\rightarrow$ loi protège, associations soutien ; "Pas de risque" $\rightarrow$ tout actif sexuel a un risque (préservatif pas 100\% utilisé, rupture, viol).
\item \textbf{Conclusion} : Appel à l'action : "Connais ton statut, protège ton avenir, protège les autres". Dépistage = acte responsable, citoyen, gratuit, confidentiel.
\end{enumerate}
\end{methode}

\begin{exoresolu}[Convaincre lors de la Journée Mondiale Sida (Argumentation)]
Énoncé : Un élève dit : « Je me sens en pleine forme, je n'ai aucune raison de faire le test de dépistage du VIH. » Rédiger un texte argumentatif (10-15 lignes) scientifique.
\tcblower
\textbf{Proposition de réponse type (niveau Probatoire) :}

Se sentir en bonne santé ne garantit pas d'être séronégatif : c'est précisément le piège biologique de l'infection à VIH. En effet, après la primo-infection souvent inaperçue, le virus entre dans une \textbf{phase asymptomatique} qui dure en moyenne huit à dix ans. Pendant toute cette période, la personne infectée ne présente \textbf{aucun signe clinique}, mais le VIH se réplique activement, détruit progressivement ses \textbf{lymphocytes T CD4} et, fait majeur, \textbf{elle peut transmettre le virus} à ses partenaires sexuels sans le savoir.

D'abord, le dépistage volontaire est le \textbf{seul moyen} de connaître son statut sérologique. Il repose sur une simple prise de sang (recherche d'anticorps anti-VIH, test rapide 4ème génération Ag/Ac), réalisée \textbf{gratuitement} dans les centres de santé publics et agréés du Cameroun (CTA, HDJ), dans le strict respect du \textbf{consentement} et de la \textbf{confidentialité}.

Ensuite, connaître tôt sa séropositivité permet d'accéder \textbf{immédiatement} au traitement antirétroviral (trithérapie standard TDF+3TC+DTG). Les ARV bloquent la multiplication virale, maintiennent les CD4 à un niveau protecteur ($> 500$/mm$^3$) et évitent l'évolution vers le stade Sida. Une personne traitée précocement a une \textbf{espérance de vie proche de la normale}.

Enfin, le dépistage protège la collectivité : une personne sous traitement efficace dont la \textbf{charge virale est durablement indétectable ($< 50$ copies/mL) ne transmet plus le VIH par voie sexuelle} (principe \textbf{I = I}, Indétectable = Intransmissible, validé scientifiquement). Se faire dépister est donc un acte de responsabilité partagée.

Par conséquent, l'argument « je me sens bien » est trompeur : c'est \textbf{justement parce qu'on se sent bien} qu'il faut se faire dépister, \textbf{avant} que la maladie ne se déclare. Connaître son statut, c'est se donner les moyens de vivre longtemps en bonne santé et de ne contaminer personne.

\textbf{Barème indicatif (10 pts) :} Phase asymptomatique/contagiosité (3) ; Principe test/gratuité/3C (2) ; Bénéfice individuel ARV précoces (2) ; I=I / bénéfice collectif (2) ; Structure/connecteurs/vocabulaire précis (1).
\end{exoresolu}

\begin{attention}[Les 5 erreurs fatales au Probatoire (SVT 1ère C/TI)]
\begin{enumerate}
\item \textbf{Confondre VIH, Séropositivité, Sida.} VIH = virus ; Séropositivité = présence Ac anti-VIH (infection confirmée) ; Sida = \textbf{stade ultime} (CD4 $< 200$/mm$^3$ \textbf{ET/OU} Infections Opportunistes définissant le Sida). \textbf{Jamais} écrire "attraper le Sida" $\rightarrow$ on est contaminé par le \textbf{VIH}.
\item \textbf{Modes de transmission fantaisistes.} VIH \textbf{ne se transmet JAMAIS} par : moustiques (ni piqûre, ni écrasement), salive/baiser social, larmes, sueur, urine, selles, poignées de main, repas partagés, piscine, toilettes, air, objets quotidiens. \textbf{Seuls} : Sang, Sperme, Sécrétions vaginales, Lait maternel (+ liquides souillés par le sang). Citer le moustique = \textbf{0 point} à la question.
\item \textbf{Inverser/Confondre courbes Charge Virale / CD4.} Relation \textbf{inverse} : VL $\uparrow$ $\Rightarrow$ CD4 $\downarrow$. Oublier unités (copies/mL ; cellules/mm$^3$) ou seuil \textbf{200 CD4/mm$^3$} = perte points. Savoir lire un graphique (log vs linéaire).
\item \textbf{Oublier la fenêtre sérologique.} Test négatif $< 6$ semaines post-risque = \textbf{non concluant}. Séroconversion (apparition Ac) = 2 à 8 semaines (médiane 3-4 sem). \textbf{Toujours} mentionner contrôle à \textbf{3 mois (12 semaines)} pour certitude. Test 4ème géné (Ag p24 + Ac) raccourcit fenêtre mais ne l'annule pas.
\item \textbf{Dire que les ARV "guérissent" ou "Indétectable = Guéri".} ARV = \textbf{suppression virale}, pas éradication. Provirus intégré dans réservoirs = persistance. Arrêt = rebond viral certain (semaines). Traitement = \textbf{À VIE}, observance $\ge 95\%$. "Indétectable" = VL $< 50$ copies/mL = non transmissible sexuellement (I=I), \textbf{PAS} virus éliminé.
\end{enumerate}
\end{attention}

\newpage

\coursec{Exercices}

\courssub{Je vérifie mes connaissances}

\begin{exercice}[QCM : le VIH et son cycle]
Pour chaque question, une seule réponse est exacte. Recopie la lettre correspondante.
\begin{enumerate}
\item Le VIH est un virus :
\begin{enumerate}
\item à ADN double brin ;
\item à ARN simple brin possédant une transcriptase inverse ;
\item à ARN double brin sans enveloppe ;
\item à ADN circulaire.
\end{enumerate}
\item La cible principale du VIH est :
\begin{enumerate}
\item le globule rouge ;
\item le lymphocyte T4 (CD4) ;
\item le neurone ;
\item la cellule musculaire.
\end{enumerate}
\item L'enzyme qui copie l'ARN viral en ADN dans la cellule hôte est :
\begin{enumerate}
\item l'ADN polymérase ;
\item la protéase ;
\item la transcriptase inverse ;
\item l'intégrase.
\end{enumerate}
\item La séroconversion correspond à :
\begin{enumerate}
\item l'apparition du Sida déclaré ;
\item l'apparition d'anticorps anti-VIH détectables dans le sang ;
\item la disparition du virus de l'organisme ;
\item la guérison de l'infection.
\end{enumerate}
\end{enumerate}
\end{exercice}

\begin{exercice}[Vrai ou faux justifié]
Pour chaque affirmation, indique si elle est vraie ou fausse, puis justifie en une ou deux phrases.
\begin{enumerate}
\item On peut contracter le VIH en partageant un repas avec une personne séropositive.
\item Une personne séropositive sous traitement antirétroviral efficace, avec une charge virale indétectable durable, ne transmet pas le virus par voie sexuelle.
\item La phase asymptomatique signifie que le virus est devenu inactif et ne se multiplie plus.
\item Le préservatif, correctement utilisé, est un moyen efficace de prévention de la transmission sexuelle du VIH.
\end{enumerate}
\end{exercice}

\begin{exercice}[Définitions]
Définis en deux lignes maximum chacun des termes suivants : \cle{charge virale}, \cle{lymphocytes CD4}, \cle{séropositivité}, \cle{Sida}, \cle{trithérapie}.
\end{exercice}

\begin{exercice}[Calcul direct : taux de CD4]
Lors d'un suivi au Centre de Traitement Agréé (CTA) de l'hôpital de district de Bafoussam, un patient présente $280$ lymphocytes CD4 par mm$^3$ de sang, alors que la valeur normale est comprise entre $500$ et \num{1500} CD4/mm$^3$.
\begin{enumerate}
\item Calcule le pourcentage de CD4 restants par rapport à la borne inférieure de la normale ($500$ CD4/mm$^3$).
\item Sachant que le seuil définissant le Sida est de $200$ CD4/mm$^3$, ce patient a-t-il atteint ce seuil ? Justifie.
\end{enumerate}
\end{exercice}

\courssub{Je m'entraîne}

\begin{exercice}[Modes de transmission et prévention]
Voici des situations de la vie courante au Cameroun :
\begin{enumerate}
\item[a.] Utilisation de la même lame de rasoir non stérilisée lors d'une scarification traditionnelle ;
\item[b.] Poignée de main avec une personne séropositive ;
\item[c.] Rapport sexuel non protégé avec un partenaire séropositif ;
\item[d.] Allaitement d'un nourrisson par une mère séropositive non traitée ;
\item[e.] Piqûre de moustique après avoir piqué une personne séropositive ;
\item[f.] Transfusion de sang non testé.
\end{enumerate}
\begin{enumerate}
\item Classe ces situations en deux groupes : situations à risque de transmission du VIH et situations sans risque.
\item Pour chaque situation à risque, propose une mesure de prévention adaptée.
\end{enumerate}
\end{exercice}

\begin{exercice}[Lecture de marqueurs biologiques]
Le tableau ci-dessous donne les résultats du suivi biologique d'un patient nouvellement diagnostiqué, avant toute mise sous traitement.

\begin{center}
\begin{tabularx}{\linewidth}{|l|X|X|X|X|}\hline
\rowcolor{popblueL} Date du bilan & Janvier 2020 & Janvier 2021 & Janvier 2022 & Janvier 2023 \\ \hline
CD4 (cellules/mm$^3$) & 620 & 540 & 430 & 310 \\ \hline
Charge virale (copies/mL) & \num{25000} & \num{40000} & \num{85000} & \num{150000} \\ \hline
\end{tabularx}
\end{center}

\begin{enumerate}
\item Décris l'évolution des deux marqueurs entre 2020 et 2023.
\item Calcule le taux de diminution annuel moyen des CD4 sur cette période (en cellules/mm$^3$/an).
\item Que traduit l'augmentation concomitante de la charge virale ?
\item À partir de quelle année le patient franchit-il le seuil de $500$ CD4/mm$^3$ ? Quelle décision médicale cela aurait-il dû entraîner ?
\end{enumerate}
\end{exercice}

\begin{exercice}[Interpréter un test de dépistage]
Trois personnes se présentent à une campagne de dépistage volontaire organisée par une ONG à Yaoundé. Les résultats du test rapide de $1^{\text{re}}$ intention sont :
\begin{itemize}
\item Personne A : test négatif, dernier risque pris il y a 5 mois ;
\item Personne B : test négatif, dernier risque pris il y a 10 jours ;
\item Personne C : test positif.
\end{itemize}
\begin{enumerate}
\item Explique ce qu'est la \cle{fenêtre sérologique} et sa durée approximative.
\item Peut-on conclure définitivement à la séronégativité de la personne B ? Justifie et propose une conduite à tenir.
\item Pourquoi le résultat de la personne C doit-il être confirmé par un second test (algorithme national de dépistage) ?
\end{enumerate}
\end{exercice}

\begin{exercice}[Classes d'antirétroviraux]
Voici quatre classes d'antirétroviraux : inhibiteurs nucléosidiques de la transcriptase inverse (INTI), inhibiteurs non nucléosidiques de la transcriptase inverse (INNTI), inhibiteurs de la protéase (IP), inhibiteurs de l'intégrase (INI).
\begin{enumerate}
\item Pour chaque classe, indique l'étape du cycle infectieux du VIH qu'elle bloque.
\item Explique pourquoi on associe toujours au moins trois molécules de classes différentes (trithérapie) plutôt qu'une seule molécule.
\end{enumerate}
\end{exercice}

\courssub{Je me prépare au Probatoire}

\begin{exercice}[Type Probatoire — Analyse de graphique : évolution naturelle de l'infection]
Le graphique ci-dessous représente, chez un patient non traité, l'évolution de la charge virale (CV, en copies/mL, échelle logarithmique) et du taux de CD4 (cellules/mm$^3$) au cours des 10 premières années suivant la contamination.

\begin{center}
\begin{tikzpicture}
\begin{axis}[
width=0.85\linewidth, height=7cm,
xlabel={Temps (années)}, ylabel={CD4 (cellules/mm$^3$)},
xmin=0, xmax=10, ymin=0, ymax=1400,
xtick={0,1,2,3,4,5,6,7,8,9,10},
legend style={at={(0.02,0.02)},anchor=south west, font=\small},
grid=both, grid style={popline!40}
]
\addplot[popcurve, color=popblue, very thick, domain=0:10, samples=100] {1200 - 85*x};
\addlegendentry{CD4}
\addplot[color=poporange, very thick, domain=0:10, samples=100] {900*exp(-3*x) + 60 + 28*x};
\addlegendentry{Charge virale (unités arbitraires)}
\end{axis}
\end{tikzpicture}
\end{center}

On rappelle les valeurs repères : pic de la primo-infection vers la $3^{\text{e}}$ semaine ; \emph{set point} viral atteint vers 6 mois ; seuil du Sida : $200$ CD4/mm$^3$.
\begin{enumerate}
\item À partir du graphique et de tes connaissances, identifie et délimite les trois phases de l'évolution naturelle de l'infection à VIH. Pour chaque phase, donne un argument tiré des courbes.
\item Estime graphiquement la valeur du \emph{set point} viral et explique sa signification biologique.
\item Détermine la durée approximative de la phase asymptomatique chez ce patient.
\item Explique, en t'appuyant sur le cycle infectieux du VIH, la relation entre l'augmentation de la charge virale et la chute des CD4.
\item Le médecin propose de débuter les ARV dès le diagnostic. Justifie cette attitude par deux arguments.
\end{enumerate}
\end{exercice}

\begin{exercice}[Type Probatoire — Analyse de données : riposte nationale et objectifs 95-95-95]
Une enquête menée dans la ville de Douala auprès d'une population de \num{2000} personnes vivant avec le VIH (PVVIH) estimées a donné les résultats suivants :

\begin{center}
\begin{tabularx}{\linewidth}{|X|c|}\hline
\rowcolor{popblueL} Indicateur & Nombre de personnes \\ \hline
PVVIH connaissant leur statut sérologique & \num{1640} \\ \hline
Personnes sous traitement antirétroviral & \num{1410} \\ \hline
Personnes ayant une charge virale supprimée (indétectable) & \num{1180} \\ \hline
\end{tabularx}
\end{center}

La stratégie onusidienne « 95-95-95 » fixe les cibles suivantes : 95 \% des PVVIH connaissent leur statut ; 95 \% des personnes diagnostiquées sont sous traitement ; 95 \% des personnes sous traitement ont une charge virale supprimée.
\begin{enumerate}
\item Calcule, pour chacun des trois « 95 », le pourcentage réellement atteint dans cette population (arrondis à l'unité).
\item Calcule l'écart (\emph{gap}) entre chaque résultat et la cible correspondante. Identifie le maillon le plus faible de la cascade.
\item Explique pourquoi la suppression de la charge virale est un double enjeu : individuel et de santé publique (notion « Indétectable = Intransmissible »).
\item Propose deux interventions ciblées et argumentées pour améliorer le maillon le plus faible, en tenant compte du contexte camerounais (dépistage communautaire, stigmatisation, accès aux CTA).
\item La stigmatisation est citée comme obstacle majeur. Rappelle le cadre juridique camerounais protégeant les PVVIH (loi n° 2007/001 du 19 avril 2007) et explique en quoi il soutient la riposte.
\end{enumerate}
\end{exercice}

\begin{exercice}[Type Probatoire — Problème biologique : résistance aux ARV et argumentation]
\textbf{Partie A.} Un patient sous traitement de première ligne contenant l'éfavirenz (EFV, un INNTI) présente un échec thérapeutique : sa charge virale redevient détectable. L'analyse du gène \emph{pol} (codant la transcriptase inverse) du virus donne les séquences partielles suivantes :

\begin{center}
\begin{tabularx}{\linewidth}{|l|X|}\hline
\rowcolor{popblueL} Virus & Séquence autour du codon 103 \\ \hline
Avant traitement (sauvage) & ... AAA ... (codant la lysine K en position 103) \\ \hline
Après échec (mutant) & ... AAT ... (codant l'asparagine N en position 103) \\ \hline
\end{tabularx}
\end{center}

\begin{enumerate}
\item Identifie la mutation apparue (notation K103N) et précise le mécanisme évolutif qui a permis sa sélection sous traitement.
\item Sachant que les INNTI se fixent sur une poche allostérique de la transcriptase inverse (site différent du site actif), explique comment la mutation K103N rend l'éfavirenz inefficace (résistance dite non compétitive).
\item Explique pourquoi cette résistance impose un changement complet de schéma thérapeutique (passage à une deuxième ligne associant d'autres classes d'ARV) et non une simple augmentation de dose.
\item Déduis-en pourquoi l'observance stricte du traitement (prise quotidienne sans oubli) est essentielle pour limiter l'apparition de résistances.
\end{enumerate}

\textbf{Partie B.} Dans un débat au sein de ton établissement, un élève affirme : « Le dépistage obligatoire du VIH pour tous les élèves avant le Probatoire serait une bonne mesure de santé publique. »

\begin{enumerate}
\item[5.] Rédige un texte argumenté de 15 lignes maximum réfutant cette proposition, en mobilisant : un argument scientifique (efficacité des stratégies de dépistage volontaire et ciblé), un argument éthique (consentement, confidentialité) et un argument juridique (loi n° 2007/001 du 19 avril 2007 portant protection des personnes en matière de VIH/Sida au Cameroun).
\end{enumerate}
\end{exercice}

\newpage

\coursec{Corrigés des exercices}

\begin{corrige}[Exercice 1 — QCM : le VIH et son cycle]
\begin{enumerate}
\item \textbf{b.} Le VIH (Virus de l'Immunodéficience Humaine) est un rétrovirus de la famille des \emph{Retroviridae}, genre \emph{Lentivirus}. Son matériel génétique est de l'\textbf{ARN simple brin} (\emph{ssRNA}) de polarité positive, \textbf{diploïde} (deux copies d'ARN identiques). Il possède dans son capside l'enzyme \textbf{transcriptase inverse} (RT) qui permet la transcription inverse de son ARN en ADN double brin (ADN proviral) une fois dans le cytoplasme de la cellule hôte.
\item \textbf{b.} La cible principale du VIH est le \textbf{lymphocyte T4} (ou lymphocyte T \cle{CD4+}), coordonnateur de la réponse immune adaptative. Le VIH se fixe via son enveloppe (gp120) sur le récepteur CD4 et un co-récepteur (CCR5 ou CXCR4) présents à la surface de ces cellules. Les globules rouges, neurones et cellules musculaires ne portent pas de récepteur CD4 fonctionnel pour l'entrée virale.
\item \textbf{c.} L'enzyme virale \textbf{transcriptase inverse} (ou rétrotranscriptase) catalyse la synthèse d'ADN sur matrice d'ARN (transcription inverse). L'ADN polymérase cellulaire réplique l'ADN cellulaire ; la protéase virale clive les polyprotéines Gag/Pol lors de la maturation ; l'intégrase insère l'ADN proviral dans le génome hôte.
\item \textbf{b.} La \textbf{séroconversion} est l'apparition dans le sang d'anticorps anti-VIH détectables par les tests de dépistage (généralement entre 2 et 12 semaines post-contamination). Elle marque le passage du statut séronégatif au statut séropositif. Ce n'est ni la guérison (infection chronique), ni la disparition du virus, ni le stade Sida (survenue tardive sans traitement).
\end{enumerate}
\end{corrige}

\begin{corrige}[Exercice 2 — Vrai ou faux justifié]
\begin{enumerate}
\item \textbf{Faux.} Le VIH ne se transmet \emph{pas} par la salive, la sueur, les larmes, ni par les contacts quotidiens (poignée de main, embrassade, partage de repas, ustensiles, toilettes, moustiques). La transmission nécessite un contact avec des liquides biologiques contaminés à forte charge virale : sang, sperme, sécrétions vaginales, lait maternel, liquide céphalo-rachidien.
\item \textbf{Vrai.} C'est le principe \textbf{« Indétectable = Intransmissible » (I = I)}, validé par de grandes études (PARTNER, Opposites Attract, HPTN 052). Une personne séropositive sous traitement antirétroviral (ARV) efficace, avec une charge virale plasmatique durablement indétectable (< \SI{50}{copies/mL}) depuis au moins 6 mois et une bonne observance, \textbf{ne transmet pas le VIH} par voie sexuelle, même sans préservatif.
\item \textbf{Faux.} La \textbf{phase asymptomatique} (ou phase de latence clinique) dure en moyenne 5 à 10 ans sans traitement. Le virus \textbf{se réplique activement} en permanence (turn-over viral élevé : $10^9$ à $10^{10}$ virions produits/jour), entraînant une destruction continue des lymphocytes CD4. L'absence de symptômes cliniques ne signifie pas inactivité virale.
\item \textbf{Vrai.} Le \textbf{préservatif} (masculin ou féminin), utilisé correctement à chaque rapport sexuel (vaginal, anal, oral), constitue une barrière physique étanche aux liquides biologiques vecteurs du VIH. C'est un moyen de prévention \textbf{efficace} (efficacité > \SI{90}{\%} en utilisation parfaite), recommandé en association avec d'autres stratégies (PrEP, TasP, dépistage).
\end{enumerate}
\end{corrige}

\begin{corrige}[Exercice 3 — Définitions]
\begin{itemize}
\item \cle{Charge virale} : Quantité d'ARN du VIH (exprimée en \SI{}{copies/mL} de plasma) circulant dans le sang. C'est le marqueur direct de l'activité de réplication virale ; elle guide la mise sous traitement et le suivi de l'efficacité thérapeutique (objectif : indétectable < \SI{50}{copies/mL}).
\item \cle{Lymphocytes CD4} : Sous-population de lymphocytes T (T4) portant le récepteur CD4, cellules cibles privilégiées du VIH. Leur numération (CD4/\SI{}{mm^3}) est le marqueur immunologique de référence : elle évalue le stade d'immunodépression et le risque de maladies opportunistes (seuil Sida : $< \SI{200}{CD4/mm^3}$).
\item \cle{Séropositivité} : Résultat positif aux tests de dépistage du VIH (recherche d'anticorps anti-VIH et/ou antigène p24), attestant l'infection par le VIH. Elle ne précise ni le stade clinique, ni la charge virale, ni l'immunité.
\item \cle{Sida} (Syndrome d'Immunodéficience Acquise) : Stade avancé de l'infection à VIH non traitée, défini par un taux de CD4 $< \SI{200}{cellules/mm^3}$ \emph{ou} la survenue d'au moins une maladie opportuniste (infection ou cancer) définissant le Sida (liste OMS/Cameroun).
\item \cle{Trithérapie} : Association thérapeutique standard de \textbf{trois molécules antirétrovirales} d'au moins deux classes pharmacologiques différentes (ex: 2 INTI + 1 INNTI ou 1 INI). Elle vise à obtenir une suppression virologique durable, prévenir la résistance et restaurer l'immunité.
\end{itemize}
\end{corrige}

\begin{corrige}[Exercice 4 — Calcul direct : taux de CD4]
\begin{enumerate}
\item Taux de CD4 du patient : \SI{280}{CD4/mm^3}. Borne inférieure de la normale : \SI{500}{CD4/mm^3}.
\[
\text{Pourcentage} = \frac{280}{500} \times 100 = 56\,\%.
\]
Le patient a conservé \textbf{56\,\%} de la borne inférieure normale de CD4.
\item Seuil définissant le Sida (immunodépression sévère) : \SI{200}{CD4/mm^3}.
Comme $280 > 200$, le patient \textbf{n'a pas atteint le seuil du Sida} au sens immunologique strict. Toutefois, avec \SI{280}{CD4/mm^3}, il présente une \textbf{immunodépression modérée à sévère} (stade 3 OMS) et une indication formelle de traitement ARV immédiat (stratégie « Treat All »).
\end{enumerate}
\end{corrige}

\begin{corrige}[Exercice 5 — Modes de transmission et prévention]
\begin{enumerate}
\item \textbf{Classification :}
\begin{center}
\begin{tabularx}{\linewidth}{|l|X|}\hline
\rowcolor{popblueL} \textbf{Situations À RISQUE} (transmission possible) & \textbf{Situations SANS RISQUE} \\ \hline
a. Lame de rasoir non stérilisée (scarification) & b. Poignée de main \\ \hline
c. Rapport sexuel non protégé & e. Piqûre de moustique \\ \hline
d. Allaitement par mère séropositive non traitée & \\ \hline
f. Transfusion de sang non testé & \\ \hline
\end{tabularx}
\end{center}
\item \textbf{Mesures de prévention adaptées :}
\begin{itemize}
\item \textbf{a.} Utilisation de matériel \textbf{stérile à usage unique} (lames, aiguilles) pour tout acte invasif (scarification, tatouage, piqûre, circoncision) ; sensibilisation des praticiens traditionnels.
\item \textbf{c.} \textbf{Préservatif} (masculin/féminin) correct et systématique ; \textbf{PrEP} (prophylaxie pré-exposition) pour personnes à risque élevé ; \textbf{TasP} (traitement comme prévention) si partenaire séropositif sous ARV efficace (charge virale indétectable).
\item \textbf{d.} \textbf{Prise en charge ARV de la mère} (idéalement avant grossesse ou dès le 1er trimestre) + \textbf{prophylaxie ARV du nourrisson} (névirapine ou zidovudine 6 semaines) + \textbf{allaitement exclusif 6 mois} (recommandation OMS/Cameroun en contexte ARV) ou lait de substitution si conditions AFASS réunies (Acceptable, Faisable, Abordable, Durable, Sûr).
\item \textbf{f.} \textbf{Dépistage systématique du VIH (et VHB, VHC, syphilis) sur tout donneur de sang} selon l'algorithme national (tests rapides + confirmation) ; promotion du don de sang bénévole, régulier, non rémunéré ; hémovigilance.
\end{itemize}
\end{enumerate}
\end{corrige}

\begin{corrige}[Exercice 6 — Lecture de marqueurs biologiques]
\begin{enumerate}
\item \textbf{Évolution 2020--2023 :}
\begin{itemize}
\item \textbf{CD4 :} Diminution \textbf{régulière et continue} : $620 \to 540 \to 430 \to 310$ cellules/\SI{}{mm^3} (perte totale de $310$ cellules en 3 ans). Le patient passe de la normale ($>500$) à une immunodépression modérée ($350-500$) puis sévère ($200-350$).
\item \textbf{Charge virale :} Augmentation \textbf{régulière et continue} : \num{25000} $\to$ \num{40000} $\to$ \num{85000} $\to$ \num{150000} \SI{}{copies/mL} (multipliée par 6 en 3 ans). Témoin d'une réplication virale non contrôlée, en progression exponentielle.
\end{itemize}
Ces deux évolutions parallèles (CD4 $\downarrow$, CV $\uparrow$) sont caractéristiques de l'\textbf{évolution naturelle de l'infection VIH non traitée}.
\item \textbf{Taux de diminution annuel moyen des CD4 :}
\[
\frac{620 - 310}{2023 - 2020} = \frac{310}{3} \approx \textbf{103,3 cellules/mm$^3$/an}.
\]
(Perte moyenne $\approx 100$ CD4/an, valeur typique de la phase asymptomatique non traitée).
\item \textbf{Signification de l'augmentation concomitante de la charge virale :} Elle traduit une \textbf{perte de contrôle immunitaire} : le système immunitaire (réponse cytotoxique CD8, anticorps) ne parvient plus à contenir la réplication virale. Le \emph{set point} viral (équilibre transitoire) est dépassé ; la réplication s'emballe, infectant et détruisant davantage de CD4, ce qui affaiblit encore la réponse immune (cercle vicieux).
\item \textbf{Franchissement du seuil de 500 CD4/mm$^3$ :} Entre janvier 2020 ($620$) et janvier 2021 ($540$). Le seuil est franchi \textbf{au cours de l'année 2020} (vers le milieu de l'année si linéaire).
\textbf{Décision médicale attendue :} Selon les recommandations OMS (2015) et Cameroun (2016) « \textbf{Treat All} » (traiter tous), \textbf{la mise sous ARV est indiquée dès le diagnostic}, quel que soit le taux de CD4. Le franchissement de 500 CD4 (ancien seuil de démarrage) aurait renforcé l'urgence, mais ne change pas la prise en charge actuelle : \textbf{initiation immédiate de la trithérapie} + conseil observance + bilan pré-thérapeutique.
\end{enumerate}
\end{corrige}

\begin{corrige}[Exercice 7 — Interpréter un test de dépistage]
\begin{enumerate}
\item \textbf{Fenêtre sérologique :} Période suivant la contamination pendant laquelle la personne est \textbf{infectée (VIH présent, réplication active, contagieuse)} mais les \textbf{anticorps anti-VIH ne sont pas encore détectables} par les tests sérologiques (tests rapides anticorps seuls ou tests 4ème génération Ag/Ac).
\textbf{Durée approximative :}
\begin{itemize}
\item Tests rapides (anticorps seuls, 1ère/3ème génération) : \textbf{3 à 12 semaines} (parfois jusqu'à 3 mois).
\item Tests 4ème génération (Ag p24 + Ac, laboratoire) : \textbf{2 à 4 semaines} (détection précoce de l'antigène p24).
\end{itemize}
Au Cameroun, l'algorithme national utilise des tests rapides (Determine, Stat-Pak, Uni-Gold) : la fenêtre est donc de \textbf{3 mois} pour une certitude totale.
\item \textbf{Personne B :} \textbf{Non}, on ne peut \emph{pas} conclure à la séronégativité. Le dernier risque date de \textbf{10 jours} : la personne est en pleine \textbf{fenêtre sérologique}. Le test rapide (anticorps) est nécessairement négatif même si l'infection a eu lieu (pas encore de séroconversion).
\textbf{Conduite à tenir :}
\begin{enumerate}
\item Informer la personne du résultat « négatif à ce jour » mais \emph{non concluant} du fait de la fenêtre.
\item Conseiller la \textbf{prévention} (préservatifs, PrEP si éligible) en attendant.
\item Programmer un \textbf{test de contrôle} : à \textbf{6 semaines} (test 4ème génération Ag/Ac au laboratoire) ou à \textbf{3 mois} (test rapide anticorps) pour lever le doute.
\end{enumerate}
\item \textbf{Personne C (test positif) :} Tout test de dépistage (même de haute qualité) a une \textbf{spécificité $< 100\,\%$} (risque de \textbf{faux positifs} : réactions croisées, grossesse, maladies auto-immunes, erreurs techniques). Au Cameroun, l'\textbf{algorithme national de dépistage} impose une \textbf{confirmation obligatoire} avant toute annonce de séropositivité et mise sous ARV :
\begin{enumerate}
\item Test 1 (Determine/Alere) : si positif $\to$
\item Test 2 (Stat-Pak/Uni-Gold) : si positif $\to$
\item Test 3 (Tie-breaker / confirmation) ou envoi échantillon au laboratoire de référence pour Western Blot / PCR VIH-1/2.
\end{enumerate}
Seule la concordance de 2 tests (ou 3 selon algorithme) valide le diagnostic. Cela évite l'annonce erronée d'une séropositivité (impact psychologique, social, médical majeur).
\end{enumerate}
\end{corrige}

\begin{corrige}[Exercice 8 — Classes d'antirétroviraux]
\begin{enumerate}
\item \textbf{Étape du cycle infectieux bloquée par chaque classe :}
\begin{center}
\begin{tabularx}{\linewidth}{|l|X|}\hline
\rowcolor{popblueL} \textbf{Classe (exemples)} & \textbf{Étape ciblée du cycle du VIH} \\ \hline
\textbf{INTI} (Inhibiteurs Nucléosidiques de la Transcriptase Inverse) : Zidovudine (AZT), Lamivudine (3TC), Emtricitabine (FTC), Ténofovir (TDF/TAF), Abacavir (ABC) & \textbf{Transcription inverse (ARN $\to$ ADN)} : analogues de nucléosides/nucleotides $\to$ incorporation dans l'ADN viral naissant $\to$ \textbf{terminaison de chaîne} (manque groupement 3'-OH). \\ \hline
\textbf{INNTI} (Inhibiteurs Non Nucléosidiques de la Transcriptase Inverse) : Éfavirenz (EFV), Névirapine (NVP), Doravirine (DOR), Rilpivirine (RPV) & \textbf{Transcription inverse} : fixation sur une \textbf{poche allostérique} (site NNRTI) de la transcriptase inverse $\to$ changement de conformation $\to$ \textbf{inhibition non compétitive} de l'activité polymérase. \\ \hline
\textbf{IP} (Inhibiteurs de la Protéase) : Lopinavir/r (LPV/r), Atazanavir/r (ATV/r), Darunavir/r (DRV/r) & \textbf{Maturation / Bourgeonnement} : inhibition de la \textbf{protéase virale} $\to$ blocage du clivage des polyprotéines Gag et Gag-Pol $\to$ production de \textbf{virions immatures, non infectieux}. \\ \hline
\textbf{INI} (Inhibiteurs de l'Intégrase) : Dolutégravir (DTG), Raltégravir (RAL), Bictegravir (BIC), Cabotégravir (CAB) & \textbf{Intégration} : blocage de l'\textbf{intégrase} $\to$ empêche l'insertion de l'ADN proviral dans le génome de la cellule hôte $\to$ l'ADN viral reste épisomal (cercle 2-LTR) et se dégrade. \\ \hline
\end{tabularx}
\end{center}
\item \textbf{Pourquoi associer 3 molécules de classes différentes (trithérapie) ?}
\begin{enumerate}
\item \textbf{Barrière génétique élevée :} Le VIH mute très vite (erreur RT $\approx 10^{-4}$/base/cycle). La probabilité qu'un virus développe \emph{simultanément} des mutations conférant la résistance à 3 molécules ciblant des étapes \emph{différentes} est le produit des probabilités individuelles : extrèmement faible ($< 10^{-12}$). Cela empêche la sélection de mutants résistants.
\item \textbf{Synergie antivirale :} Attaque simultanée de plusieurs étapes clés (RT, Intégrase, Protéase) $\to$ suppression virologique maximale et rapide (charge virale indétectable).
\item \textbf{Éviter la monothérapie/bithérapie :} Historiquement, la monothérapie (AZT seul) ou bithérapie (2 INTI) a entraîné une \textbf{échec virologique rapide} (semaines/mois) par sélection de résistances. La trithérapie (1996) a transformé l'infection en maladie chronique contrôlée.
\end{enumerate}
\end{enumerate}
\end{corrige}

\begin{corrige}[Exercice 9 — Type Probatoire : Analyse de graphique]
\textbf{Barème indicatif : 10 points}
\begin{itemize}
\item Q1 : 3 pts (1 pt par phase bien identifiée + argument graphique).
\item Q2 : 2 pts (valeur + signification).
\item Q3 : 1 pt (durée).
\item Q4 : 2 pts (mécanisme cycle infectieux).
\item Q5 : 2 pts (2 arguments distincts).
\end{itemize}

\textbf{Points de vigilance :}
\begin{itemize}
\item Lire les courbes sur le graphique fourni (axes, unités, échelles).
\item Distinguer \emph{primo-infection} (aiguë), \emph{phase asymptomatique} (chronique), \emph{phase Sida} (terminale).
\item Relier mécaniquement réplication virale (CV) et destruction CD4 (cycle infectieux).
\item Citer les recommandations actuelles (Treat All, I=I).
\end{itemize}

\begin{enumerate}
\item \textbf{Identification des trois phases (arguments graphiques) :}
\begin{enumerate}
\item \textbf{Phase de primo-infection (0 -- $\sim$ 6 mois / 0,5 an) :}
\begin{itemize}
\item \textbf{Courbe CV (orange) :} Pic très élevé au départ ($t=0$, $\sim 960$ u.a.), chute brutale exponentielle pour atteindre un plateau bas (\emph{set point}) vers $t=0,5$ an ($\sim 120$ u.a.).
\item \textbf{Courbe CD4 (bleu) :} Chute aiguë et rapide depuis la normale ($\sim 1200$) vers $\sim 800$ cellules/\SI{}{mm^3}.
\end{itemize}
\item \textbf{Phase asymptomatique / latence clinique ($\sim$ 6 mois -- $\sim$ 8--9 ans) :}
\begin{itemize}
\item \textbf{Courbe CV :} Plateau relativement stable (\emph{set point} viral) entre $\sim 100$ et $\sim 200$ u.a. (légère remontée linéaire due au terme $+28x$).
\item \textbf{Courbe CD4 :} Déclin \textbf{linéaire, lent et continu} (pente $\approx -85$ cellules/\SI{}{mm^3}/an), passant de $\sim 800$ à $\sim 400$ cellules/\SI{}{mm^3}.
\end{itemize}
\item \textbf{Phase d'accélération / Sida déclaré ($\sim$ 9--10 ans et au-delà) :}
\begin{itemize}
\item \textbf{Courbe CV :} \textbf{Remontée exponentielle nette} (terme $+28x$ devient dominant, ou reprise réplication non contrôlée), dépassant $\sim 300$ u.a. à 10 ans.
\item \textbf{Courbe CD4 :} Poursuite de la chute linéaire, franchissement du \textbf{seuil de 200 CD4/mm$^3$} (seuil Sida) prévu vers $t \approx 11,7$ ans (extrapolation $1200 - 85t = 200 \Rightarrow t = 11,76$). Sur la fenêtre graphique (10 ans), le patient est en \textbf{immunodépression sévère} ($< 350$ CD4) pré-Sida.
\end{itemize}
\end{enumerate}
\item \textbf{Set point viral :} Estimé graphiquement vers \textbf{6 mois ($t=0,5$ an)}. Valeur lue sur la courbe orange : $\approx \textbf{120 unités arbitraires}$ (minimum de la courbe avant remontée).
\textbf{Signification biologique :} Niveau d'équilibre transitoire entre \textbf{réplication virale} et \textbf{réponse immune cytotoxique (CD8+)}. Il prédit la vitesse de progression : un \emph{set point} bas = progression lente ; un \emph{set point} haut = progression rapide. Ici, $\sim 120$ u.a. correspond à un \emph{set point} modéré.
\item \textbf{Durée de la phase asymptomatique :} De la fin de la primo-infection ($\sim 0,5$ an) au début de la remontée de la charge virale / chute CD4 accélérée ($\sim 8,5$--$9$ ans). Soit environ \textbf{8 à 8,5 ans} pour ce patient (durée typique 5--10 ans).
\item \textbf{Relation CV $\uparrow$ / CD4 $\downarrow$ (mécanisme cycle infectieux) :}
\begin{enumerate}
\item Le VIH infecte les lymphocytes CD4+ via fixation gp120-CD4/co-récepteur.
\item Après transcription inverse (RT) et intégration (Intégrase), la cellule devient \textbf{usine à virions} (transcription/translation génome viral).
\item L'assemblage et le bourgeonnement (maturation par Protéase) libèrent des milliers de nouveaux virions $\to$ \textbf{lyse de la cellule productrice} (mort directe).
\item De plus, protéines virales (gp120, Nef, Tat) et activation immune chronique induisent \textbf{apoptose} des CD4 non infectés (bystander effect) et dysfonction thymique.
\item \textbf{Bilan :} Plus la charge virale (CV) est élevée $\to$ plus de cellules infectées par unité de temps $\to$ plus de destruction CD4 $\to$ chute du taux CD4.
\end{enumerate}
\item \textbf{Justification ARV dès le diagnostic (2 arguments) :}
\begin{enumerate}
\item \textbf{Préservation/restauration immunitaire :} L'initiation précoce (CD4 $> 500$) évite la perte irréversible de diversité du répertoire T (thymus), réduit le réservoir viral, diminue l'inflammation chronique et le risque de maladies non-Sida (cardiovasculaires, rénales, cancers). Données : étude START, TEMPRANO (ANRS 12136, Cameroun/Côte d'Ivoire).
\item \textbf{Prévention de la transmission (TasP / I=I) :} La suppression virologique durable (CV indétectable) élimine le risque de transmission sexuelle (études PARTNER, HPTN 052). C'est un pilier de la stratégie \textbf{95-95-95} et de l'élimination de l'épidémie.
\end{enumerate}
(Autres arguments acceptables : réduction morbidité/mortalité, simplification schémas (INI : DTG), coût-efficacité à long terme).
\end{enumerate}
\end{corrige}

\begin{corrige}[Exercice 10 — Type Probatoire : Riposte nationale et objectifs 95-95-95]
\textbf{Barème indicatif : 12 points}
\begin{itemize}
\item Q1 : 3 pts (3 calculs justes + arrondi).
\item Q2 : 3 pts (3 écarts + identification maillon faible justifiée).
\item Q3 : 2 pts (double enjeu individuel + santé publique + I=I).
\item Q4 : 2 pts (2 interventions ciblées, argumentées, contextuelles Cameroun).
\item Q5 : 2 pts (loi 2007/001 citée + lien avec riposte).
\end{itemize}

\textbf{Points de vigilance :}
\begin{itemize}
\item Respecter la définition en cascade des 95-95-95 (dénominateurs différents).
\item Arrondir à l'unité (pourcentage entier).
\item « Maillon le plus faible » = plus bas pourcentage atteint ou plus grand gap.
\item Interventions : concrètes, réalistes au Cameroun (communautaire, stigmatisation, CTA).
\item Loi 2007/001 : citer articles clés (consentement, confidentialité, non-discrimination).
\end{itemize}

\begin{enumerate}
\item \textbf{Calcul des pourcentages atteints (cascade 95-95-95) :}
\begin{itemize}
\item \textbf{1er 95 (Connaissance du statut) :} $\frac{1640}{2000} \times 100 = \textbf{82\,\%}$.
\item \textbf{2ème 95 (Sous ARV parmi diagnostiqués) :} $\frac{1410}{1640} \times 100 = 85,97\,\% \approx \textbf{86\,\%}$.
\item \textbf{3ème 95 (CV supprimée parmi sous ARV) :} $\frac{1180}{1410} \times 100 = 83,69\,\% \approx \textbf{84\,\%}$.
\end{itemize}
\item \textbf{Ecarts (gaps) par rapport aux cibles (95\,\%) :}
\begin{itemize}
\item 1er 95 : $95 - 82 = \textbf{13 points}$.
\item 2ème 95 : $95 - 86 = \textbf{9 points}$.
\item 3ème 95 : $95 - 84 = \textbf{11 points}$.
\end{itemize}
\textbf{Maillon le plus faible :} Le \textbf{1er 95 (dépistage / connaissance du statut)} avec seulement \textbf{82\,\%} atteint (gap de 13 points). C'est l'entrée de la cascade : sans diagnostic, pas de traitement, pas de suppression virale. Améliorer ce maillon tire toute la cascade vers le haut.
\item \textbf{Double enjeu de la suppression de la charge virale (CV indétectable) :}
\begin{itemize}
\item \textbf{Enjeu individuel (bénéfice clinique) :} Arrêt de la réplication virale $\to$ reconstruction immunitaire (hausse CD4) $\to$ prévention des maladies opportunistes et comorbidités $\to$ \textbf{espérance de vie quasi normale}, qualité de vie préservée.
\item \textbf{Enjeu de santé publique (bénéfice collectif) :} \textbf{Indétectable = Intransmissible (I = I).} Une personne sous ARV efficace (CV $< \SI{50}{copies/mL}$ durablement) \textbf{ne transmet pas le VIH} par voie sexuelle. La suppression virale à l'échelle populationnelle réduit l'incidence (nouvelles infections) : c'est la \textbf{prévention par le traitement (TasP)}, pilier de l'élimination de l'épidémie à l'horizon 2030 (ONUSIDA).
\end{itemize}
\item \textbf{Deux interventions ciblées pour améliorer le 1er 95 (dépistage) au Cameroun :}
\begin{enumerate}
\item \textbf{Dépistage communautaire et auto-test VIH (VIHAT) :} Déployer des \textbf{équipes mobiles} et \textbf{pairs-éducateurs} dans les lieux de vie des populations clés (HSH, travailleurs du sexe, usagers de drogues, jeunes, hommes) et zones rurales mal desservies. Proposer l'\textbf{auto-test VIH} (OraQuick, test salivaire) : confidentiel, non invasif, résultat en 20 min, réalisé par la personne elle-même. Cela lève les freins : stigmatisation en structure de santé, distance géographique, horaires incompatibles, peur du jugement. \emph{Preuve :} Projets ATLAS, STAR en Afrique subsaharienne ont doublé le taux de dépistage chez les hommes et populations clés.
\item \textbf{Dépistage index (partner notification) optimisé et linkage immédiat :} Pour \textbf{chaque nouveau cas diagnostiqué} au CTA, proposer systématiquement le dépistage aux \textbf{partenaires sexuels et enfants} (avec consentement, confidentialité, conseil). Assurer un \textbf{liaison (linkage) active} vers le CTA le jour même ou sous 7 jours : navigateur patient, rendez-vous programmé, accompagnement physique, suppression des frais utilisateurs (gratuité ARV + bilans). \emph{Contexte Cameroun :} Taux de linkage encore perfectible ($\sim 70-80\,\%$) ; perte de vue entre diagnostic et 1ère consultation ARV.
\end{enumerate}
\item \textbf{Cadre juridique camerounais (Loi n° 2007/001 du 19 avril 2007) :}
Cette loi \textbf{protège les droits des personnes vivant avec le VIH (PVVIH)} et soutient la riposte par :
\begin{itemize}
\item \textbf{Consentement libre et éclairé (Art. 3) :} « Nul ne peut être soumis à un test de dépistage du VIH sans son consentement libre et éclairé. » $\to$ Garantit le dépistage \textbf{volontaire} (VCT), base éthique de l'approche « 95-95-95 ».
\item \textbf{Confidentialité (Art. 8) :} « Les résultats des tests... sont confidentiels. » $\to$ Rassure les PVVIH, encourage le recours aux services sans peur de divulgation (famille, employeur, communauté).
\item \textbf{Non-discrimination (Art. 12, 13, 14) :} Interdiction de toute discrimination dans l'emploi, l'éducation, l'accès aux soins, au logement, aux assurances. Sanctions pénales prévues. $\to$ Lutte contre la \textbf{stigmatisation}, obstacle majeur au dépistage et à l'observance.
\item \textbf{Accès aux soins et ARV (Art. 16) :} Droit à la prise en charge médicale, psychologique, sociale ; gratuité des ARV (décret d'application). $\to$ Soutient le 2ème et 3ème 95.
\end{itemize}
En protégeant les droits humains, cette loi crée un \textbf{environnement légal favorable} (enabling environment) indispensable pour que les PVVIH acceptent de se faire dépister, se lient aux soins et observent leur traitement.
\end{enumerate}
\end{corrige}

\begin{corrige}[Exercice 11 — Type Probatoire : Résistance aux ARV et argumentation]
\textbf{Barème indicatif : 12 points}
\begin{itemize}
\item Partie A (Q1--Q4) : 8 pts (2 pts par question : mutation + mécanisme, explication résistance, justification changement schéma, lien observance/résistance).
\item Partie B (Q5) : 4 pts (argument scientifique 1,5 pt + éthique 1,5 pt + juridique 1 pt, rédaction cohérente $< 15$ lignes).
\end{itemize}

\textbf{Points de vigilance :}
\begin{itemize}
\item Notation mutation : K103N (acide aminé position).
\item Mécanisme : sélection naturelle (pression thérapeutique), pas « mutation induite par le médicament ».
\item Résistance INNTI : non-compétitive (poche allostérique), pas site actif.
\item Résistance croisée 1ère génération INNTI (EFV/NVP) $\to$ changement classe (INI).
\item Observance : concentration plasmatique, fenêtre de sélection, réplication résiduelle.
\item Argumentation Q5 : 3 piliers distincts (scientifique, éthique, juridique), concis, sans répétition.
\end{itemize}

\textbf{Partie A.}
\begin{enumerate}
\item \textbf{Mutation :} \textbf{K103N} (substitution de la \textbf{Lysine K} par l'\textbf{Asparagine N} en position 103 de la transcriptase inverse).
\textbf{Mécanisme évolutif :} \textbf{Sélection naturelle (darwinienne) sous pression thérapeutique.}
\begin{itemize}
\item La transcriptase inverse du VIH n'a pas d'activité de relecture (proofreading) $\to$ taux d'erreur élevé ($\sim 3 \times 10^{-5}$/base/cycle) $\to$ diversité génétique énorme (quasi-espèces).
\item \textbf{Avant traitement :} Le mutant K103N existe à très basse fréquence (mutation spontanée) mais est \textbf{minoritaire} (coût de fitness possible).
\item \textbf{Sous traitement EFV :} Le virus sauvage (K103) est \textbf{inhibé} (réplication bloquée). Le mutant K103N, \textbf{résistant}, continue de se répliquer $\to$ \textbf{avantage sélectif massif} $\to$ devient \textbf{majoritaire} (population virale dominante) en quelques semaines. C'est une \textbf{sélection de mutants préexistants}, pas une mutation induite par le médicament (théorie de Luria-Delbrück).
\end{itemize}
\item \textbf{Mécanisme de résistance K103N (non-compétitive) :}
Les INNTI (EFV, NVP) se fixent sur une \textbf{poche allostérique} (site NNRTI) de la sous-unité p66 de la transcriptase inverse, \emph{loin du site actif} (liaison dNTP). Cette fixation verrouille l'enzyme dans une conformation inactive.
La mutation \textbf{K103N} (Lys $\to$ Asn) modifie la structure de cette poche :
\begin{itemize}
\item Perte d'une liaison hydrogène clé entre la chaîne latérale de K103 et l'EFV.
\item Encombrement stérique / modification hydrophobie $\to$ \textbf{diminution drastique de l'affinité} de l'EFV pour sa poche ($K_d$ augmenté 100--1000x).
\item L'enzyme conserve son activité catalytique (site actif intact) mais l'inhibiteur ne peut plus s'y fixer efficacement.
\end{itemize}
C'est une résistance \textbf{non-compétitive} : l'inhibiteur ne compete pas avec le substrat (dNTP) ; il ne peut plus bloquer l'enzyme.
\item \textbf{Changement de schéma complet (2ème ligne) indispensable :}
\begin{itemize}
\item \textbf{Résistance croisée totale :} K103N confère une résistance de haut niveau à \textbf{tous les INNTI de 1ère génération} (EFV, NVP). La rilpivirine (RPV, 2ème gen) conserve une activité partielle mais n'est pas recommandée en cas d'échec sur EFV sans test de résistance.
\item \textbf{Augmentation de dose inefficace :} La résistance est due à une \textbf{perte d'affinité} (liaison), pas à une concentration insuffisante. Augmenter la dose ne restaure pas la liaison (plafonnement de l'effet) et augmente la \textbf{toxicité} (neuropsychiatrique pour EFV, hépatique/cutanée pour NVP).
\item \textbf{Stratégie 2ème ligne (OMS/Cameroun) :} Passer à un schéma à \textbf{barrière génétique élevée} sans résistance croisée : \textbf{INI (Dolutégravir DTG) + 2 INTI (TDF + 3TC ou TAF + FTC)}. L'INI (DTG) a une barrière génétique très haute (mutations multiples nécessaires) et aucune résistance croisée avec les INNTI/INTI.
\end{itemize}
\item \textbf{Lien Observance -- Résistance :}
L'observance stricte (prise \textbf{quotidienne, sans oubli, à heure fixe}) maintient la \textbf{concentration plasmatique des ARV au-dessus de la CMI (Concentration Minimale Inhibitrice) pour le virus sauvage} en permanence.
\begin{itemize}
\item \textbf{Oubli / sous-dosage :} Concentration chute dans la « \textbf{fenêtre de sélection mutante} » (entre CMI virus sauvage et CMI mutant résistant).
\item Dans cette fenêtre : le virus sauvage est inhibé, mais le mutant résistant (s'il existe) \textbf{se réplique} $\to$ sélection \emph{in vivo}.
\item \textbf{Observance parfaite :} Charge virale indétectable ($< \SI{50}{copies/mL}$) $\to$ \textbf{arrêt quasi total de la réplication} $\to$ \textbf{pas de nouvelle mutation} $\to$ \textbf{pas de sélection de résistance}.
\end{itemize}
\emph{Conclusion :} L'observance est le \textbf{premier pilier de la durabilité} du traitement de 1ère ligne.
\end{enumerate}

\textbf{Partie B.}
\begin{enumerate}
\setcounter{enumi}{4}
\item \textbf{Réfutation argumentée (proposition de rédaction) :}
\par\noindent\textbf{« Imposer un dépistage VIH obligatoire à tous les élèves avant le Probatoire est une mesure \textbf{contre-productive et illégale}.}
\par\textbf{Scientifiquement}, le dépistage volontaire, confidentiel et accompagné (conseil pré/post-test) atteint les populations à risque avec un meilleur \emph{linkage to care} (entrée en soins). Le dépistage obligatoire génère de la fuite, des faux négatifs (fenêtre sérologique non respectée) et une rupture de confiance, réduisant l'efficacité globale de la riposte (cibles 95-95-95).
\par\textbf{Éthiquement}, il viole le \textbf{consentement libre et éclairé} (droit à l'autonomie corporelle) et la \textbf{confidentialité} médicale, piliers de la relation soignant-soigné. La contrainte stigmatise les élèves, alimente la peur et la discrimination, poussant les jeunes \emph{loin} des services de santé.
\par\textbf{Juridiquement}, l'\textbf{Article 3 de la Loi n° 2007/001 du 19 avril 2007} dispose : ``Nul ne peut être soumis à un test de dépistage du VIH sans son consentement libre et éclairé''. L'Article 8 garantit la confidentialité des résultats. Une telle mesure exposerait l'État et les établissements à des poursuites judiciaires et contredit l'engagement du Cameroun dans la riposte nationale respectueuse des droits humains.''
\end{enumerate}
\end{corrige}

\newpage

\coursec{Fiche bilan}

\begin{popfigure}
\centering
\begin{tikzpicture}[
  mindmap,
  concept color=popblue,
  level 1 concept/.append style={level distance=4.5cm, sibling angle=60, font=\footnotesize\bfseries},
  level 2 concept/.append style={level distance=3.2cm, font=\scriptsize},
  every node/.style={concept, execute at begin node=\hskip0pt},
  branch1/.style={concept color=popgreen},
  branch2/.style={concept color=poporange},
  branch3/.style={concept color=poppurple},
  branch4/.style={concept color=poppink},
  branch5/.style={concept color=popgold},
  branch6/.style={concept color=popdark}
]
\node[concept, text width=3.2cm, align=center] {Lutte contre\\ le VIH/Sida}
  child[branch1] { node[concept, text width=2.8cm, align=center] {1. Structure \& Cycle infectieux}
    child { node[concept, text width=2.2cm, align=center] {Enveloppe, Capside, ARN, RT, Intégrase, Protéase} }
    child { node[concept, text width=2.2cm, align=center] {Fixation CD4 $\to$ Fusion $\to$ RT $\to$ Intégration $\to$ Transcription/Traduction $\to$ Assemblage $\to$ Bourgeonnement} }
  }
  child[branch2] { node[concept, text width=2.8cm, align=center] {2. Transmission \& Prévention}
    child { node[concept, text width=2.2cm, align=center] {Sang, Sexuel, Mère-Enfant (PTME)} }
    child { node[concept, text width=2.2cm, align=center] {ABC : Abstinence, Fidélité, Préservatif ; PrEP/PEP ; Stérilisation ; Dépistage} }
  }
  child[branch3] { node[concept, text width=2.8cm, align=center] {3. Évolution \& Marqueurs}
    child { node[concept, text width=2.2cm, align=center] {Séroconversion $\to$ Phase asymptomatique $\to$ Sida} }
    child { node[concept, text width=2.2cm, align=center] {Charge virale $\uparrow$, CD4 $\downarrow$ ; Infections opportunistes} }
  }
  child[branch4] { node[concept, text width=2.8cm, align=center] {4. Dépistage}
    child { node[concept, text width=2.2cm, align=center] {Tests rapides (Déterminer, SD Bioline), ELISA, Western Blot} }
    child { node[concept, text width=2.2cm, align=center] {Algorithme national : T1 $\to$ T2 $\to$ T3 (discordance)} }
  }
  child[branch5] { node[concept, text width=2.8cm, align=center] {5. Prise en charge (ARV)}
    child { node[concept, text width=2.2cm, align=center] {TAR : 2 INI + 1 INNTI (ex: TDF+3TC+DTG)} }
    child { node[concept, text width=2.2cm, align=center] {Indétectable = Intransmissible (I=I) ; Observance $\ge$ 95\%} }
  }
  child[branch6] { node[concept, text width=2.8cm, align=center] {6. Psychosocial \& Riposte}
    child { node[concept, text width=2.2cm, align=center] {Stigmatisation, Droits, Conseil, CNLS, Gratuité ARV} }
  };
\end{tikzpicture}
\legende{Carte mentale récapitulative de la leçon NCT04 : structure du VIH, transmission/prévention, évolution clinique/biologique, dépistage, traitement ARV et riposte nationale (Cameroun).}
\end{popfigure}

\begin{bilan}
\courssub{Le VIH : Structure et Cycle Infectieux}
\begin{definition}[VIH (Virus de l'Immunodéficience Humaine)]
Rétrovirus à ARN simple brin ($+$), enveloppé, de la famille des \emph{Lentiviridae}. Deux types : VIH-1 (pandémique, dont groupe M au Cameroun) et VIH-2 (Afrique de l'Ouest, moins virulent).
\end{definition}
\begin{itemize}
  \item \textbf{Structure :} Enveloppe lipidique (glycoprotéines \cle{gp120} et \cle{gp41}), Capside protéique (p24), 2 brins d'ARN génomique, 3 enzymes virales : \cle{Transcriptase Inverse (TI/RT)}, \cle{Intégrase}, \cle{Protéase}.
  \item \textbf{Cible principale :} Lymphocytes T \cle{CD4$^+$} (récepteur CD4 + co-récepteurs CCR5/CXCR4). Aussi macrophages, cellules dendritiques.
\end{itemize}
\begin{propriete}[Cycle Infectieux (7 étapes clés)]
\begin{enumerate}
  \item \textbf{Fixation :} gp120 $\leftrightarrow$ CD4 + co-récepteur.
  \item \textbf{Fusion :} gp41 médie la fusion enveloppe/membrane $\to$ entrée capside.
  \item \textbf{Transcription Inverse :} TI synthétise ADN double brin (ADN proviral) à partir de l'ARN viral. \textit{Point d'action des INNTI (TDF, 3TC) et INI (DTG, RAL)}.
  \item \textbf{Intégration :} Intégrase insère l'ADN proviral dans l'ADN hôte. \textit{Point d'action des II (DTG, RAL, BIC)}.
  \item \textbf{Transcription/Traduction :} Machinerie hôte $\to$ ARNm viraux $\to$ Protéines polyprotéines (Gag, Pol, Env).
  \item \textbf{Assemblage/Maturation :} Protéase clive les polyprotéines $\to$ virions matures infectieux. \textit{Point d'action des IP (LPV/r, ATV/r, DRV/r)}.
  \item \textbf{Bourgeonnement :} Libération du nouveau virion.
\end{enumerate}
\end{propriete}

\courssub{Modes de Transmission et Prévention}
\begin{center}
\begin{tabularx}{\linewidth}{|l|X|X|}\hline
\rowcolor{popblueL}
\textbf{Voie} & \textbf{Mécanisme / Situation à risque} & \textbf{Prévention Principale (Contexte Cameroun)} \\ \hline
\textbf{Sanguine} & Transfusion (sécurisée au Cameroun), partage seringues/aiguilles (UDI), scarifications, excisions, accidents d'exposition (AES). & Stérilisation matériel, usage unique, \cle{PrEP} (Prophylaxie Pré-Exposition), \cle{PEP} (Post-Exposition $<$ 72h), sécurité transfusionnelle (CNTS). \\ \hline
\textbf{Sexuelle} & Rapports non protégés (vaginaux, anaux, oraux), IST non traitées (ulcérations facilitent entrée). & \cle{Préservatif} (masculin/féminin), Fidélité/Abstinence, Dépistage couple, Traitement IST, \cle{PrEP} (populations clés). \\ \hline
\textbf{Mère-Enfant (PTME)} & Grossesse (transplacentaire), Accouchement (contact sang/sécrétions), Allaitement. & \cle{PTME Option B+ : TAR à vie pour toute femme enceinte/allaitante VIH+}, AZT nourrisson 6 sem, Allaitement exclusif 6 mois + TAR mère $\to$ risque $<$ 5\%. \\ \hline
\end{tabularx}
\end{center}
\begin{attention}[Pièges Fréquents]
Ne pas confondre \textbf{PrEP} (prise \emph{avant} risque continu, ex: couple sérodifférent) et \textbf{PEP} (urgence \emph{après} exposition accidentelle, 28 jours de TAR). Le VIH ne se transmet \textbf{PAS} par : salive, sueur, larmes, piqûres moustiques, contacts quotidiens (mains, repas, toilettes, piscine).
\end{attention}

\courssub{Évolution Naturelle et Marqueurs Biologiques}
\begin{center}
\begin{tabularx}{\linewidth}{|l|X|X|X|}\hline
\rowcolor{popblueL}
\textbf{Phase} & \textbf{Clinique} & \textbf{Charge Virale (CV) / ARN VIH} & \textbf{CD4 (cellules/mm$^3$)} \\ \hline
\textbf{Infection Primaire / Séroconversion} (2--8 sem) & Syndrome pseudo-grippal / mononucléosique (fièvre, adénopathies, éruption), souvent méconnu. & \textbf{Pic extrême} ($>10^6$ copies/mL), très contagieux. & Chute brutale transitoire. \\ \hline
\textbf{Phase Asymptomatique / Latence Clinique} (années) & Asymptomatique ou adénopathies généralisées persistantes (AGP). & \textbf{Point de consigne} (équilibre), variable ($10^3$--$10^5$). & Déclin lent et continu ($\sim$50--80/an). \\ \hline
\textbf{Sida Déclaré (Stade 3/4 OMS)} & Infections opportunistes (TB, Pneumocystose, Toxoplasmose, Candidose œsophagienne...), cancers (Sarcome Kaposi, Lymphome), amaigrissement $>$10\%. & \textbf{Très élevée} ($>10^5$), progression rapide. & \textbf{$< 200$} (immunodépression sévère). \\ \hline
\end{tabularx}
\end{center}
\begin{aretenir}[Marqueurs de Suivi]
\begin{itemize}
  \item \cle{CD4} : Indicateur de l'\textbf{immunité} (risque IO). Seuil traitement : \textbf{TAR pour tous (Treat All)} depuis 2016 (OMS/Cameroun), quel que soit le taux.
  \item \cle{Charge Virale (CV)} : Indicateur de la \textbf{réplication virale} et de l'\textbf{efficacité du TAR}. Objectif : \textbf{Indétectable ($< 50$ ou $< 200$ copies/mL)} $\to$ \cle{I = I} (Indétectable = Intransmissible).
\end{itemize}
\end{aretenir}

\courssub{Dépistage : Principes, Stratégies, Algorithme National}
\begin{methode}[Algorithme de Dépistage Cameroun (3 Tests Rapides)]
\begin{enumerate}
  \item \textbf{Test 1 (T1) :} Haute sensibilité (ex: \textit{Determine HIV-1/2}). Si Négatif $\to$ \textbf{Négatif} (sauf fenêtre sérologique).
  \item \textbf{Test 2 (T2) :} Haute spécificité (ex: \textit{SD Bioline HIV-1/2 3.0}). Si T1 Positif \textbf{ET} T2 Positif $\to$ \textbf{Positif}.
  \item \textbf{Test 3 (T3) :} Arbitre (ex: \textit{Uni-Gold HIV} ou \textit{First Response}). Si T1 Positif \textbf{ET} T2 Négatif (Discordance) $\to$ Faire T3.
  \begin{itemize}
    \item T3 Positif $\to$ \textbf{Positif}.
    \item T3 Négatif $\to$ \textbf{Indéterminé} $\to$ Contrôle à 4--6 semaines (fenêtre sérologique) ou PCR ADN/ARN (nourrisson PTME).
  \end{itemize}
\end{enumerate}
\end{methode}
\begin{definition}[Fenêtre Sérologique]
Période post-infection où les anticorps ne sont pas encore détectables (tests sérologiques négatifs) mais où le virus est présent (CV détectable par PCR). Durée : \textbf{3 à 6 semaines} (tests 3$^{\text{ème}}$/4$^{\text{ème}}$ gén). \textbf{Conduite :} Re-test à 6 semaines si exposition récente.
\end{definition}
\begin{itemize}
  \item \textbf{Stratégies :} Dépistage Volontaire (DV), Dépistage Proposé par le Prestataire (DPP/PICT), Autotest (VIH Auto-test), Dépistage Communautaire, PTME (systématique CPN).
  \item \textbf{Conseil Pré/Post-Test :} Obligatoire, confidentiel, non directif. Clé pour l'adhésion au soin.
\end{itemize}

\courssub{Prise en Charge : TAR et Suivi}
\begin{propriete}[Lignes Directrices Cameroun (Alignées OMS 2016+) : \textbf{Treat All}]
\textbf{Initiation TAR :} \textbf{Impérative et Rapide} (idéalement J0 ou J7) pour \textbf{TOUS} les PVVIH, quel que soit le stade clinique ou le taux de CD4.
\end{propriete}
\begin{center}
\begin{tabularx}{\linewidth}{|l|X|X|}\hline
\rowcolor{popblueL}
\textbf{Ligne} & \textbf{Schéma Recommandé (Adultes/Ados $\ge$ 30kg/10ans)} & \textbf{Commentaires} \\ \hline
\textbf{1$^{\text{ère}}$ Ligne (Préférée)} & \textbf{TDF + 3TC (ou FTC) + DTG} (TLD) & \cle{DTG (Dolutégravir)} : II de 2$^{\text{ème}}$ gén, barrière génétique haute, peu d'effets secondaires, compatible grossesse (donnée Tsepamo rassurantes). Coût faible. \\ \hline
\textbf{1$^{\text{ère}}$ Ligne (Alternative)} & TDF + 3TC + EFV 400mg (TLE400) ou AZT + 3TC + DTG & Si CI DTG (rare) ou Intolérance TDF (rénal/ostéoporose) $\to$ AZT. EFV 600mg évité (effets neuropsychiaques). \\ \hline
\textbf{2$^{\text{nde}}$ Ligne (Échec 1$^{\text{ère}}$)} & \textbf{AZT (ou TDF) + 3TC + ATV/r (ou LPV/r ou DRV/r)} & Passage aux \cle{IP (Inhibiteurs Protéase)} boostés au Ritonavir. Nécessite \textbf{Résistance (Génotypage)} idéalement. \\ \hline
\textbf{3$^{\text{ème}}$ Ligne} & Nouveaux ARV (DTG si non utilisé, DRV/r, ETV, MVC, FTR...) & Comité d'experts, génotypage obligatoire. \\ \hline
\end{tabularx}
\end{center}
\begin{aretenir}[Suivi Biologique Standard (Cameroun)]
\begin{itemize}
  \item \textbf{Mois 3 (M3) :} CV (cible $< 50$ ou $< 200$), Créatinine (si TDF), Hb/Glycémie (si AZT).
  \item \textbf{M6, M12 puis tous les 6--12 mois :} CV (indicateur principal succès), CD4 (si CV détectable ou stade avancé initial).
  \item \textbf{Échec Virologique :} CV $\ge 1000$ copies/mL à deux reprises (3 mois d'intervalle) \textbf{APRÈS} renforcement observance $\to$ Passage 2$^{\text{nde}}$ ligne.
\end{itemize}
\end{aretenir}
\begin{attention}[Observance et Résistance]
L'observance $\ge 95\%$ est la \textbf{seule} garantie de durabilité du TAR. Oubli $\to$ Réplication $\to$ Mutations $\to$ Résistance. \textbf{Conseil :} Intégrer prise dans routine quotidienne, boîtes piluliers, groupes de soutien (GAS), disclosure sélectif.
\end{attention}

\courssub{Aspects Psychosociaux, Droits et Riposte Nationale}
\begin{itemize}
  \item \textbf{Stigmatisation/Discrimination :} Frein majeur au dépistage et à l'observance. Lutte : éducation, médias, leaders d'opinion, loi n°2007/009 (protection PVVIH).
  \item \textbf{Droits :} Accès aux soins (gratuité TAR, examens de suivi au Cameroun depuis 2007), confidentialité, non-discrimination travail/école, droit à la procréation (PTME, PMA).
  \item \textbf{Riposte Nationale :} \cle{CNLS} (Comité National de Lutte contre le Sida) coordonne la réponse multisectorielle. \cle{Fonds Mondial}, \cle{PEPFAR}, \cle{ONUSIDA} partenaires clés. Objectifs \cle{95-95-95} (2025) : 95\% diagnostiqués, 95\% sous TAR, 95\% charge virale supprimée.
  \item \textbf{Populations Clés (PK) :} HSH, TS, UD, Travailleurs du sexe, Détenus. Prévention combinée (PrEP, préservatifs, lubrifiants, réduction des risques, services adaptés sans stigma).
\end{itemize}
\end{bilan}

\begin{aretenir}[Checklist avant le Probatoire]
$\square$ Je sais décrire la structure du VIH (enveloppe, capside, ARN, 3 enzymes) et nommer sa cible cellulaire (LT CD4$^+$).
$\square$ Je restitue les 7 étapes du cycle infectieux et j'associe chaque classe d'ARV à sa cible enzymatique (TI, Intégrase, Protéase).
$\square$ Je cite les 3 modes de transmission (Sang, Sexuel, Mère-Enfant) et je donne une mesure de prévention spécifique \textbf{par mode} (ex: PTME Option B+, PrEP, PEP, Préservatif, Stérilisation).
$\square$ Je distingue clairement \textbf{PrEP} (avant risque) de \textbf{PEP} (après exposition $<$ 72h, 28 jours).
$\square$ Je décris l'évolution clinique en 3 phases (Primaire/Séroconversion, Asymptomatique, Sida) et je relie les variations de \textbf{Charge Virale} (pic, plateau, hausse) et \textbf{CD4} (chute, déclin, $<200$).
$\square$ Je sais interpréter un tableau de suivi : CV indétectable = Succès TAR = I=I ; CV $\ge 1000$ copies/mL (confirmé) = Échec virologique.
$\square$ Je maîtrise l'algorithme national de dépistage à 3 tests rapides (T1, T2, T3) et j'interprète les 4 issues possibles (Négatif, Positif, Indéterminé, Discordance).
$\square$ Je connais la définition de la \textbf{fenêtre sérologique} (3-6 sem) et la conduite à tenir (re-test à 6 sem / PCR nourrisson).
$\square$ Je cite le schéma de 1$^{\text{ère}}$ ligne préférée au Cameroun (\textbf{TDF + 3TC + DTG}) et j'explique pourquoi DTG est privilégié (barrière génétique, tolérance, grossesse).
$\square$ Je définis les critères de passage en 2$^{\text{nde}}$ ligne (Échec virologique confirmé + observance renforcée) et la classe thérapeutique utilisée (IP boostés).
$\square$ J'argumente en faveur du dépistage volontaire (connaissance statut, entrée précoce soins, prévention transmission, PTME, I=I).
$\square$ Je cite les objectifs \textbf{95-95-95} et le rôle du \textbf{CNLS} dans la riposte nationale (gratuité ARV, coordination).
\end{aretenir}

\findecours

\end{document}
